% 复变函数与积分变换课程笔记 —— 主文档（正文排版集中在 0.样式.tex，章节文件只放内容）
% 编译：xelatex 1.main.tex 两遍；成品输出见 build.ps1
\documentclass[11pt,a4paper,oneside,openany]{book}
\usepackage[UTF8,fontset=windows,scheme=plain]{ctex}
\usepackage{fontspec}
\setmainfont{Times New Roman}
\setsansfont{Microsoft YaHei}
\newcommand{\YHBbooktrue}{}
\newcommand{\YHBcourse}{复变函数与积分变换}
\newcommand{\YHBversion}{v0.9}
\newcommand{\YHBquote}{以复数为桥，连接分析与几何；以积分变换为器，求解工程与物理问题。}
\newcommand{\YHBcoverSize}{33pt}
\newcommand{\YHBcoverLead}{40pt}
\usepackage{xcolor}
\definecolor{yhbAccent}{HTML}{1F3A5F}
\definecolor{yhbAccent2}{HTML}{3F6E8C}
% ============================================================
% 0.样式.tex —— 课程笔记统一正文排版
% 用法：主文件先定义课程信息与主色，再 % ============================================================
% 0.样式.tex —— 课程笔记统一正文排版
% 用法：主文件先定义课程信息与主色，再 % ============================================================
% 0.样式.tex —— 课程笔记统一正文排版
% 用法：主文件先定义课程信息与主色，再 \input{0.样式.tex}
%   \newcommand{\YHBcourse}{微积分A}      % 课程名（页眉与书签用）
%   \newcommand{\YHBversion}{v1.2}        % 版本号（封面用）
%   \newcommand{\YHBquote}{……}            % 封面题句
%   \definecolor{yhbAccent}{HTML}{24465C} % 主色（深）
%   \definecolor{yhbAccent2}{HTML}{347A80}% 辅色（浅）
%   \YHBbooktrue / \YHBbookfalse          % 是否 book 文档类（有 \chapter）
% 章节文件只写内容，不要重复 \usepackage、不要改字号。
% ============================================================

% ---------- 页面与字体 ----------
\usepackage{geometry}
\geometry{a4paper,left=24mm,right=24mm,top=23mm,bottom=24mm,headheight=15pt,headsep=7mm,footskip=13mm}
\usepackage{microtype}
\usepackage{amsmath,amssymb,amsfonts,mathtools}
\usepackage{esint}
\usepackage{graphicx}
\usepackage{adjustbox}
\usepackage{array,booktabs,multirow,makecell,tabularx,longtable}
\usepackage{xcolor,colortbl}
\usepackage{enumitem}
\usepackage{needspace,etoolbox}
\usepackage{titlesec,titletoc}
\usepackage{fancyhdr}
\usepackage{caption}
\usepackage[most]{tcolorbox}

% ---------- 颜色别名（章节里用的旧名继续有效） ----------
\colorlet{mainblue}{yhbAccent}
\colorlet{secondcyan}{yhbAccent2}
\colorlet{thirdteal}{yhbAccent2}
\colorlet{highlight}{yhbAccent!6!white}
\colorlet{tablehead}{yhbAccent!12!white}
\definecolor{textgray}{HTML}{4A5057}

% ---------- 课程标签 ----------
\providecommand{\YHBdoclabel}{课程笔记}
% ---------- 正文节奏 ----------
\linespread{1.3}
\setlength{\parindent}{0pt}
\setlength{\parskip}{6pt plus 1.5pt minus 1.5pt}
\setlength{\emergencystretch}{2.5em}
\raggedbottom
\widowpenalty=10000
\clubpenalty=10000
\setlength{\tabcolsep}{7pt}
\renewcommand{\arraystretch}{1.3}
\allowdisplaybreaks[2]
\AtBeginDocument{%
  \setlength{\abovedisplayskip}{10pt plus 2pt minus 2pt}%
  \setlength{\belowdisplayskip}{10pt plus 2pt minus 2pt}%
  \setlength{\abovedisplayshortskip}{6pt}%
  \setlength{\belowdisplayshortskip}{6pt}%
}
\captionsetup{font=small,labelfont=bf,skip=6pt}

% ---------- 列表 ----------
\setlist{leftmargin=2em,labelsep=0.6em,itemsep=4pt,parsep=2pt,topsep=6pt,partopsep=0pt}
\setlist[itemize,2]{label=\textbullet,leftmargin=1.6em}

% ---------- 标题层级 ----------
\ifdefined\YHBbooktrue
  \titleformat{\chapter}[display]
    {\sffamily\bfseries\color{yhbAccent}}
    {\normalsize\sffamily 第\,\thechapter\,章}
    {6pt}
    {\LARGE}
    [{\vspace{5pt}{\color{yhbAccent!55}\titlerule[1pt]}}]
  \titlespacing*{\chapter}{0pt}{14pt}{24pt}
  \titleformat{\section}{\Large\sffamily\bfseries\color{yhbAccent}}{\thesection}{0.7em}{}
  \titlespacing*{\section}{0pt}{18pt plus 3pt minus 2pt}{8pt}
\else
  \titleformat{\section}{\Large\sffamily\bfseries\color{yhbAccent}}{\thesection}{0.7em}{}[{\vspace{3pt}{\color{yhbAccent!45}\titlerule[0.7pt]}}]
  \titlespacing*{\section}{0pt}{18pt plus 3pt minus 2pt}{9pt}
\fi
\titleformat{\subsection}{\large\sffamily\bfseries\color{yhbAccent!85!black}}{\thesubsection}{0.7em}{}
\titleformat{\subsubsection}{\normalsize\bfseries\color{black!85}}{\thesubsubsection}{0.7em}{}
\titlespacing*{\subsection}{0pt}{12pt plus 2pt minus 1pt}{5pt}
\titlespacing*{\subsubsection}{0pt}{9pt plus 2pt minus 1pt}{4pt}
\pretocmd{\section}{\Needspace{6\baselineskip}}{}{}
\pretocmd{\subsection}{\Needspace{5\baselineskip}}{}{}
\pretocmd{\subsubsection}{\Needspace{4\baselineskip}}{}{}

\ifdefined\YHBbooktrue
  \setcounter{secnumdepth}{2}
  \setcounter{tocdepth}{1}
\else
  \setcounter{secnumdepth}{3}
  \setcounter{tocdepth}{2}
\fi
\renewcommand{\contentsname}{目录}
% 目录编号颜色由 titletoc 的 titlecontents 控制，不再重定义 \numberline

% ---------- 目录样式 ----------
\ifdefined\YHBbooktrue
  \titlecontents{chapter}[0pt]
    {\addvspace{6pt}\sffamily\bfseries\color{yhbAccent}}
    {\contentslabel{2.2em}}{}{\titlerule*[0.6em]{.}\contentspage}
  \titlecontents{section}[2.2em]
    {\addvspace{2pt}}
    {\contentslabel{2.2em}}{}{\titlerule*[0.6em]{.}\contentspage}
\else
  \titlecontents{section}[0pt]
    {\addvspace{5pt}\sffamily\bfseries\color{yhbAccent}}
    {\contentslabel{2.2em}}{}{\titlerule*[0.6em]{.}\contentspage}
  \titlecontents{subsection}[2.2em]
    {\addvspace{2pt}}
    {\contentslabel{2.4em}}{}{\titlerule*[0.6em]{.}\contentspage}
\fi

% ---------- 页眉页脚 ----------
\pagestyle{fancy}
\fancyhf{}
\fancyhead[L]{\footnotesize\sffamily\color{textgray}\YHBcourse\ \YHBdoclabel}
\ifdefined\YHBbooktrue
  \fancyhead[R]{\footnotesize\sffamily\color{textgray}\nouppercase{\leftmark}}
  \renewcommand{\chaptermark}[1]{\markboth{#1}{}}
\else
  \fancyhead[R]{\footnotesize\sffamily\color{textgray}\nouppercase{\rightmark}}
  \renewcommand{\sectionmark}[1]{\markright{#1}}
\fi
\fancyfoot[C]{\footnotesize\color{textgray}\thepage}
\renewcommand{\headrulewidth}{0.35pt}
\renewcommand{\headrule}{{\color{yhbAccent!45}\hrule width\headwidth height\headrulewidth}}
\fancypagestyle{plain}{%
  \fancyhf{}%
  \fancyfoot[C]{\footnotesize\color{textgray}\thepage}%
  \renewcommand{\headrulewidth}{0pt}%
}

% ---------- 公式与考点框 ----------
\newcommand{\keyformula}[1]{%
  \begin{tcolorbox}[enhanced,breakable,colback=highlight,colframe=yhbAccent!30!white,boxrule=0.5pt,arc=2.5pt,left=11pt,right=11pt,top=9pt,bottom=9pt,boxsep=0pt,before skip=11pt,after skip=11pt,halign=center,width=\linewidth]
  \adjustbox{max width=\linewidth-2em}{$\displaystyle #1$}
  \end{tcolorbox}}
\newcommand{\corepoint}[1]{%
  \begin{tcolorbox}[enhanced,breakable,colback=highlight,colframe=yhbAccent2,boxrule=0pt,borderline west={2.4pt}{0pt}{yhbAccent2},arc=0pt,left=11pt,right=11pt,top=8pt,bottom=8pt,before skip=11pt,after skip=11pt]
  \textbf{\color{yhbAccent}核心考点}\quad #1
  \end{tcolorbox}}
\newcommand{\symb}[2]{\textbf{#1}：#2\par}

% ---------- 超链接（务必最后加载） ----------
\usepackage[unicode=true,colorlinks=true,linkcolor=yhbAccent,citecolor=yhbAccent,urlcolor=yhbAccent!85!black,pdfborder={0 0 0},bookmarksnumbered=true,pdfstartview=FitH]{hyperref}
\usepackage{hypcap}
\hypersetup{pdftitle={\YHBcourse 课程笔记},pdfauthor={浮萍},pdfsubject={课程学习与复习笔记},pdfcreator={XeLaTeX}}

%   \newcommand{\YHBcourse}{微积分A}      % 课程名（页眉与书签用）
%   \newcommand{\YHBversion}{v1.2}        % 版本号（封面用）
%   \newcommand{\YHBquote}{……}            % 封面题句
%   \definecolor{yhbAccent}{HTML}{24465C} % 主色（深）
%   \definecolor{yhbAccent2}{HTML}{347A80}% 辅色（浅）
%   \YHBbooktrue / \YHBbookfalse          % 是否 book 文档类（有 \chapter）
% 章节文件只写内容，不要重复 \usepackage、不要改字号。
% ============================================================

% ---------- 页面与字体 ----------
\usepackage{geometry}
\geometry{a4paper,left=24mm,right=24mm,top=23mm,bottom=24mm,headheight=15pt,headsep=7mm,footskip=13mm}
\usepackage{microtype}
\usepackage{amsmath,amssymb,amsfonts,mathtools}
\usepackage{esint}
\usepackage{graphicx}
\usepackage{adjustbox}
\usepackage{array,booktabs,multirow,makecell,tabularx,longtable}
\usepackage{xcolor,colortbl}
\usepackage{enumitem}
\usepackage{needspace,etoolbox}
\usepackage{titlesec,titletoc}
\usepackage{fancyhdr}
\usepackage{caption}
\usepackage[most]{tcolorbox}

% ---------- 颜色别名（章节里用的旧名继续有效） ----------
\colorlet{mainblue}{yhbAccent}
\colorlet{secondcyan}{yhbAccent2}
\colorlet{thirdteal}{yhbAccent2}
\colorlet{highlight}{yhbAccent!6!white}
\colorlet{tablehead}{yhbAccent!12!white}
\definecolor{textgray}{HTML}{4A5057}

% ---------- 课程标签 ----------
\providecommand{\YHBdoclabel}{课程笔记}
% ---------- 正文节奏 ----------
\linespread{1.3}
\setlength{\parindent}{0pt}
\setlength{\parskip}{6pt plus 1.5pt minus 1.5pt}
\setlength{\emergencystretch}{2.5em}
\raggedbottom
\widowpenalty=10000
\clubpenalty=10000
\setlength{\tabcolsep}{7pt}
\renewcommand{\arraystretch}{1.3}
\allowdisplaybreaks[2]
\AtBeginDocument{%
  \setlength{\abovedisplayskip}{10pt plus 2pt minus 2pt}%
  \setlength{\belowdisplayskip}{10pt plus 2pt minus 2pt}%
  \setlength{\abovedisplayshortskip}{6pt}%
  \setlength{\belowdisplayshortskip}{6pt}%
}
\captionsetup{font=small,labelfont=bf,skip=6pt}

% ---------- 列表 ----------
\setlist{leftmargin=2em,labelsep=0.6em,itemsep=4pt,parsep=2pt,topsep=6pt,partopsep=0pt}
\setlist[itemize,2]{label=\textbullet,leftmargin=1.6em}

% ---------- 标题层级 ----------
\ifdefined\YHBbooktrue
  \titleformat{\chapter}[display]
    {\sffamily\bfseries\color{yhbAccent}}
    {\normalsize\sffamily 第\,\thechapter\,章}
    {6pt}
    {\LARGE}
    [{\vspace{5pt}{\color{yhbAccent!55}\titlerule[1pt]}}]
  \titlespacing*{\chapter}{0pt}{14pt}{24pt}
  \titleformat{\section}{\Large\sffamily\bfseries\color{yhbAccent}}{\thesection}{0.7em}{}
  \titlespacing*{\section}{0pt}{18pt plus 3pt minus 2pt}{8pt}
\else
  \titleformat{\section}{\Large\sffamily\bfseries\color{yhbAccent}}{\thesection}{0.7em}{}[{\vspace{3pt}{\color{yhbAccent!45}\titlerule[0.7pt]}}]
  \titlespacing*{\section}{0pt}{18pt plus 3pt minus 2pt}{9pt}
\fi
\titleformat{\subsection}{\large\sffamily\bfseries\color{yhbAccent!85!black}}{\thesubsection}{0.7em}{}
\titleformat{\subsubsection}{\normalsize\bfseries\color{black!85}}{\thesubsubsection}{0.7em}{}
\titlespacing*{\subsection}{0pt}{12pt plus 2pt minus 1pt}{5pt}
\titlespacing*{\subsubsection}{0pt}{9pt plus 2pt minus 1pt}{4pt}
\pretocmd{\section}{\Needspace{6\baselineskip}}{}{}
\pretocmd{\subsection}{\Needspace{5\baselineskip}}{}{}
\pretocmd{\subsubsection}{\Needspace{4\baselineskip}}{}{}

\ifdefined\YHBbooktrue
  \setcounter{secnumdepth}{2}
  \setcounter{tocdepth}{1}
\else
  \setcounter{secnumdepth}{3}
  \setcounter{tocdepth}{2}
\fi
\renewcommand{\contentsname}{目录}
% 目录编号颜色由 titletoc 的 titlecontents 控制，不再重定义 \numberline

% ---------- 目录样式 ----------
\ifdefined\YHBbooktrue
  \titlecontents{chapter}[0pt]
    {\addvspace{6pt}\sffamily\bfseries\color{yhbAccent}}
    {\contentslabel{2.2em}}{}{\titlerule*[0.6em]{.}\contentspage}
  \titlecontents{section}[2.2em]
    {\addvspace{2pt}}
    {\contentslabel{2.2em}}{}{\titlerule*[0.6em]{.}\contentspage}
\else
  \titlecontents{section}[0pt]
    {\addvspace{5pt}\sffamily\bfseries\color{yhbAccent}}
    {\contentslabel{2.2em}}{}{\titlerule*[0.6em]{.}\contentspage}
  \titlecontents{subsection}[2.2em]
    {\addvspace{2pt}}
    {\contentslabel{2.4em}}{}{\titlerule*[0.6em]{.}\contentspage}
\fi

% ---------- 页眉页脚 ----------
\pagestyle{fancy}
\fancyhf{}
\fancyhead[L]{\footnotesize\sffamily\color{textgray}\YHBcourse\ \YHBdoclabel}
\ifdefined\YHBbooktrue
  \fancyhead[R]{\footnotesize\sffamily\color{textgray}\nouppercase{\leftmark}}
  \renewcommand{\chaptermark}[1]{\markboth{#1}{}}
\else
  \fancyhead[R]{\footnotesize\sffamily\color{textgray}\nouppercase{\rightmark}}
  \renewcommand{\sectionmark}[1]{\markright{#1}}
\fi
\fancyfoot[C]{\footnotesize\color{textgray}\thepage}
\renewcommand{\headrulewidth}{0.35pt}
\renewcommand{\headrule}{{\color{yhbAccent!45}\hrule width\headwidth height\headrulewidth}}
\fancypagestyle{plain}{%
  \fancyhf{}%
  \fancyfoot[C]{\footnotesize\color{textgray}\thepage}%
  \renewcommand{\headrulewidth}{0pt}%
}

% ---------- 公式与考点框 ----------
\newcommand{\keyformula}[1]{%
  \begin{tcolorbox}[enhanced,breakable,colback=highlight,colframe=yhbAccent!30!white,boxrule=0.5pt,arc=2.5pt,left=11pt,right=11pt,top=9pt,bottom=9pt,boxsep=0pt,before skip=11pt,after skip=11pt,halign=center,width=\linewidth]
  \adjustbox{max width=\linewidth-2em}{$\displaystyle #1$}
  \end{tcolorbox}}
\newcommand{\corepoint}[1]{%
  \begin{tcolorbox}[enhanced,breakable,colback=highlight,colframe=yhbAccent2,boxrule=0pt,borderline west={2.4pt}{0pt}{yhbAccent2},arc=0pt,left=11pt,right=11pt,top=8pt,bottom=8pt,before skip=11pt,after skip=11pt]
  \textbf{\color{yhbAccent}核心考点}\quad #1
  \end{tcolorbox}}
\newcommand{\symb}[2]{\textbf{#1}：#2\par}

% ---------- 超链接（务必最后加载） ----------
\usepackage[unicode=true,colorlinks=true,linkcolor=yhbAccent,citecolor=yhbAccent,urlcolor=yhbAccent!85!black,pdfborder={0 0 0},bookmarksnumbered=true,pdfstartview=FitH]{hyperref}
\usepackage{hypcap}
\hypersetup{pdftitle={\YHBcourse 课程笔记},pdfauthor={浮萍},pdfsubject={课程学习与复习笔记},pdfcreator={XeLaTeX}}

%   \newcommand{\YHBcourse}{微积分A}      % 课程名（页眉与书签用）
%   \newcommand{\YHBversion}{v1.2}        % 版本号（封面用）
%   \newcommand{\YHBquote}{……}            % 封面题句
%   \definecolor{yhbAccent}{HTML}{24465C} % 主色（深）
%   \definecolor{yhbAccent2}{HTML}{347A80}% 辅色（浅）
%   \YHBbooktrue / \YHBbookfalse          % 是否 book 文档类（有 \chapter）
% 章节文件只写内容，不要重复 \usepackage、不要改字号。
% ============================================================

% ---------- 页面与字体 ----------
\usepackage{geometry}
\geometry{a4paper,left=24mm,right=24mm,top=23mm,bottom=24mm,headheight=15pt,headsep=7mm,footskip=13mm}
\usepackage{microtype}
\usepackage{amsmath,amssymb,amsfonts,mathtools}
\usepackage{esint}
\usepackage{graphicx}
\usepackage{adjustbox}
\usepackage{array,booktabs,multirow,makecell,tabularx,longtable}
\usepackage{xcolor,colortbl}
\usepackage{enumitem}
\usepackage{needspace,etoolbox}
\usepackage{titlesec,titletoc}
\usepackage{fancyhdr}
\usepackage{caption}
\usepackage[most]{tcolorbox}

% ---------- 颜色别名（章节里用的旧名继续有效） ----------
\colorlet{mainblue}{yhbAccent}
\colorlet{secondcyan}{yhbAccent2}
\colorlet{thirdteal}{yhbAccent2}
\colorlet{highlight}{yhbAccent!6!white}
\colorlet{tablehead}{yhbAccent!12!white}
\definecolor{textgray}{HTML}{4A5057}

% ---------- 课程标签 ----------
\providecommand{\YHBdoclabel}{课程笔记}
% ---------- 正文节奏 ----------
\linespread{1.3}
\setlength{\parindent}{0pt}
\setlength{\parskip}{6pt plus 1.5pt minus 1.5pt}
\setlength{\emergencystretch}{2.5em}
\raggedbottom
\widowpenalty=10000
\clubpenalty=10000
\setlength{\tabcolsep}{7pt}
\renewcommand{\arraystretch}{1.3}
\allowdisplaybreaks[2]
\AtBeginDocument{%
  \setlength{\abovedisplayskip}{10pt plus 2pt minus 2pt}%
  \setlength{\belowdisplayskip}{10pt plus 2pt minus 2pt}%
  \setlength{\abovedisplayshortskip}{6pt}%
  \setlength{\belowdisplayshortskip}{6pt}%
}
\captionsetup{font=small,labelfont=bf,skip=6pt}

% ---------- 列表 ----------
\setlist{leftmargin=2em,labelsep=0.6em,itemsep=4pt,parsep=2pt,topsep=6pt,partopsep=0pt}
\setlist[itemize,2]{label=\textbullet,leftmargin=1.6em}

% ---------- 标题层级 ----------
\ifdefined\YHBbooktrue
  \titleformat{\chapter}[display]
    {\sffamily\bfseries\color{yhbAccent}}
    {\normalsize\sffamily 第\,\thechapter\,章}
    {6pt}
    {\LARGE}
    [{\vspace{5pt}{\color{yhbAccent!55}\titlerule[1pt]}}]
  \titlespacing*{\chapter}{0pt}{14pt}{24pt}
  \titleformat{\section}{\Large\sffamily\bfseries\color{yhbAccent}}{\thesection}{0.7em}{}
  \titlespacing*{\section}{0pt}{18pt plus 3pt minus 2pt}{8pt}
\else
  \titleformat{\section}{\Large\sffamily\bfseries\color{yhbAccent}}{\thesection}{0.7em}{}[{\vspace{3pt}{\color{yhbAccent!45}\titlerule[0.7pt]}}]
  \titlespacing*{\section}{0pt}{18pt plus 3pt minus 2pt}{9pt}
\fi
\titleformat{\subsection}{\large\sffamily\bfseries\color{yhbAccent!85!black}}{\thesubsection}{0.7em}{}
\titleformat{\subsubsection}{\normalsize\bfseries\color{black!85}}{\thesubsubsection}{0.7em}{}
\titlespacing*{\subsection}{0pt}{12pt plus 2pt minus 1pt}{5pt}
\titlespacing*{\subsubsection}{0pt}{9pt plus 2pt minus 1pt}{4pt}
\pretocmd{\section}{\Needspace{6\baselineskip}}{}{}
\pretocmd{\subsection}{\Needspace{5\baselineskip}}{}{}
\pretocmd{\subsubsection}{\Needspace{4\baselineskip}}{}{}

\ifdefined\YHBbooktrue
  \setcounter{secnumdepth}{2}
  \setcounter{tocdepth}{1}
\else
  \setcounter{secnumdepth}{3}
  \setcounter{tocdepth}{2}
\fi
\renewcommand{\contentsname}{目录}
% 目录编号颜色由 titletoc 的 titlecontents 控制，不再重定义 \numberline

% ---------- 目录样式 ----------
\ifdefined\YHBbooktrue
  \titlecontents{chapter}[0pt]
    {\addvspace{6pt}\sffamily\bfseries\color{yhbAccent}}
    {\contentslabel{2.2em}}{}{\titlerule*[0.6em]{.}\contentspage}
  \titlecontents{section}[2.2em]
    {\addvspace{2pt}}
    {\contentslabel{2.2em}}{}{\titlerule*[0.6em]{.}\contentspage}
\else
  \titlecontents{section}[0pt]
    {\addvspace{5pt}\sffamily\bfseries\color{yhbAccent}}
    {\contentslabel{2.2em}}{}{\titlerule*[0.6em]{.}\contentspage}
  \titlecontents{subsection}[2.2em]
    {\addvspace{2pt}}
    {\contentslabel{2.4em}}{}{\titlerule*[0.6em]{.}\contentspage}
\fi

% ---------- 页眉页脚 ----------
\pagestyle{fancy}
\fancyhf{}
\fancyhead[L]{\footnotesize\sffamily\color{textgray}\YHBcourse\ \YHBdoclabel}
\ifdefined\YHBbooktrue
  \fancyhead[R]{\footnotesize\sffamily\color{textgray}\nouppercase{\leftmark}}
  \renewcommand{\chaptermark}[1]{\markboth{#1}{}}
\else
  \fancyhead[R]{\footnotesize\sffamily\color{textgray}\nouppercase{\rightmark}}
  \renewcommand{\sectionmark}[1]{\markright{#1}}
\fi
\fancyfoot[C]{\footnotesize\color{textgray}\thepage}
\renewcommand{\headrulewidth}{0.35pt}
\renewcommand{\headrule}{{\color{yhbAccent!45}\hrule width\headwidth height\headrulewidth}}
\fancypagestyle{plain}{%
  \fancyhf{}%
  \fancyfoot[C]{\footnotesize\color{textgray}\thepage}%
  \renewcommand{\headrulewidth}{0pt}%
}

% ---------- 公式与考点框 ----------
\newcommand{\keyformula}[1]{%
  \begin{tcolorbox}[enhanced,breakable,colback=highlight,colframe=yhbAccent!30!white,boxrule=0.5pt,arc=2.5pt,left=11pt,right=11pt,top=9pt,bottom=9pt,boxsep=0pt,before skip=11pt,after skip=11pt,halign=center,width=\linewidth]
  \adjustbox{max width=\linewidth-2em}{$\displaystyle #1$}
  \end{tcolorbox}}
\newcommand{\corepoint}[1]{%
  \begin{tcolorbox}[enhanced,breakable,colback=highlight,colframe=yhbAccent2,boxrule=0pt,borderline west={2.4pt}{0pt}{yhbAccent2},arc=0pt,left=11pt,right=11pt,top=8pt,bottom=8pt,before skip=11pt,after skip=11pt]
  \textbf{\color{yhbAccent}核心考点}\quad #1
  \end{tcolorbox}}
\newcommand{\symb}[2]{\textbf{#1}：#2\par}

% ---------- 超链接（务必最后加载） ----------
\usepackage[unicode=true,colorlinks=true,linkcolor=yhbAccent,citecolor=yhbAccent,urlcolor=yhbAccent!85!black,pdfborder={0 0 0},bookmarksnumbered=true,pdfstartview=FitH]{hyperref}
\usepackage{hypcap}
\hypersetup{pdftitle={\YHBcourse 课程笔记},pdfauthor={浮萍},pdfsubject={课程学习与复习笔记},pdfcreator={XeLaTeX}}

\usepackage{mathrsfs}
\graphicspath{{../图片/}}
\usepackage{tikz}
\begin{document}
% 引入封面
% ============================================================
% 2.cover.tex —— 统一封面（信息全部来自主文件定义）
% 需要的宏：\YHBcourse、\YHBversion、\YHBquote、yhbAccent、yhbAccent2
% 可选：\YHBcoverSize / \YHBcoverLead（课程名字号与行距）、\YHBcoversub（副标题）
% ============================================================
\providecommand{\YHBcoverSize}{38pt}
\providecommand{\YHBcoverLead}{46pt}
\providecommand{\YHBcoversub}{课程笔记}
\begin{titlepage}
\thispagestyle{empty}
\pdfbookmark[0]{封面}{cover}

\begin{tikzpicture}[remember picture,overlay]
  % 顶部深色块
  \fill[yhbAccent] (current page.north west) rectangle ([yshift=-132mm]current page.north east);
  % 分隔细条（辅色）
  \fill[yhbAccent2] ([yshift=-132mm]current page.north west) rectangle ([yshift=-135mm]current page.north east);
  % 极细横线，弱装饰，不压字
  \foreach \y in {20,34,48,62,76,90,104,118}{%
    \draw[white,opacity=0.08,line width=0.5pt]
      ([xshift=16mm,yshift=-\y mm]current page.north west) --
      ([xshift=-16mm,yshift=-\y mm]current page.north east);}
\end{tikzpicture}

\vspace*{32mm}
\noindent
{\xeCJKsetup{CJKecglue={}}\color{white}\sffamily\bfseries\fontsize{\YHBcoverSize}{\YHBcoverLead}\selectfont \YHBcourse\par}
\vspace{6mm}
{\color{white}\sffamily\fontsize{25}{32}\selectfont \YHBcoversub\par}
\vspace{9mm}
{\color{white}\fontsize{12.5}{18}\selectfont 作者\quad 浮萍\qquad 版本\quad \YHBversion\par}

\vspace{20mm}
{\sffamily\bfseries\fontsize{13}{20}\selectfont\color{yhbAccent}\YHBquote\par}

\vfill
{\color{yhbAccent}\rule{\textwidth}{1.4pt}}\par
\vspace{7mm}
{\fontsize{11.5}{17}\selectfont\sffamily
\textbf{QQ 邮箱}\quad\href{mailto:3071974740@qq.com}{3071974740@qq.com}\par
\vspace{2mm}
\textbf{GitHub}\quad\href{https://github.com/Duckweed-yhb}{github.com/Duckweed-yhb}\par}
\vspace{7mm}
{\fontsize{10.5}{16}\selectfont\color{black!78}
本笔记为个人学习整理，内容仍在持续完善。如发现内容、公式、文字或排版问题，欢迎通过上述 QQ 邮箱反馈；请尽量注明版本、页码与问题描述，便于核对和修订。\par}
\vspace*{12mm}
\end{titlepage}


% 生成目录（极简写法，避免解析错误）
\newpage
\thispagestyle{empty}
\renewcommand{\contentsname}{目录}
\tableofcontents
\clearpage
\setcounter{page}{1}

% 引入章节文件
% 章序与教材《复变函数与积分变换》(哈工大·第三版) 8 章一一对应：
%   3  → 教材第1章 复数与复变函数      4  → 教材第2章 解析函数
%   5  → 教材第3章 复变函数的积分      6  → 教材第4章 级数（§4.5 孤立奇点取自教材 §5.1）
%   7  → 教材第5章 留数                7a → 教材第6章 保形映射
%   8  → 教材第7章 傅里叶变换          9  → 教材第8章 拉普拉斯变换
% 子文件：3.复数与复变函数.tex
% 对应教材：《复变函数与积分变换》(哈工大数学系编, 第三版, 科学出版社) 第 1 章
% 对应讲义：包益欣《复变函数》第 1 章
% 节次：§1.1 复数运算及几何表示（1.1.1 概念与四则运算 / 1.1.2 几何表示 / 1.1.3 共轭复数
%       / 1.1.4 乘除、乘方与开方 / 1.1.5 复球面与无穷远点）
%       §1.2 复平面上的点集（1.2.1 基本概念 / 1.2.2 区域和曲线）
%       §1.3 复变函数（1.3.1 定义与几何意义 / 1.3.2 极限与连续性）
% 由 1.main.tex 用 \input 引入：只写内容，不写 \documentclass、不重复 \usepackage

\chapter{复数与复变函数}

\section{复数运算及几何表示}

\subsection{复数概念及四则运算}
实系数方程 \(x^{2}+1=0\) 在实数域内无解，为此引进虚单位 \(\mathrm{i}\)，规定
\[
\mathrm{i}^{2}=-1 .
\]
称形如
\[
z=x+\mathrm{i}y\qquad (x,y\in\mathbb{R})
\]
的数为\textbf{复数}，\(x,y\) 分别称为 \(z\) 的\textbf{实部}与\textbf{虚部}，记作
\[
x=\operatorname{Re}z,\qquad y=\operatorname{Im}z .
\]
虚部为零的复数即实数，故 \(\mathbb{R}\subset\mathbb{C}\)；实部为零而虚部不为零的复数称为\textbf{纯虚数}。两个复数相等当且仅当实部与虚部分别相等：
\[
z_{1}=z_{2}\iff \operatorname{Re}z_{1}=\operatorname{Re}z_{2}\ \text{且}\ \operatorname{Im}z_{1}=\operatorname{Im}z_{2}.
\]

设 \(z_{1}=x_{1}+\mathrm{i}y_{1}\)，\(z_{2}=x_{2}+\mathrm{i}y_{2}\)，四则运算定义为
\[
z_{1}\pm z_{2}=(x_{1}\pm x_{2})+\mathrm{i}(y_{1}\pm y_{2}),
\]
\keyformula{z_{1}z_{2}=(x_{1}x_{2}-y_{1}y_{2})+\mathrm{i}\,(x_{1}y_{2}+x_{2}y_{1})}
\keyformula{\frac{z_{1}}{z_{2}}=\frac{x_{1}x_{2}+y_{1}y_{2}}{x_{2}^{2}+y_{2}^{2}}
+\mathrm{i}\,\frac{x_{2}y_{1}-x_{1}y_{2}}{x_{2}^{2}+y_{2}^{2}}\qquad (z_{2}\neq 0)}
其中乘法由分配律与 \(\mathrm{i}^{2}=-1\) 直接得到，除法则是“分子分母同乘分母的共轭”的结果。四则运算的结果仍是复数，且加法与乘法满足交换律、结合律与分配律，故 \(\mathbb{C}\) 构成一个域。

由此立得
\[
\operatorname{Re}(z_{1}\pm z_{2})=\operatorname{Re}z_{1}\pm\operatorname{Re}z_{2},\qquad
\operatorname{Im}(z_{1}\pm z_{2})=\operatorname{Im}z_{1}\pm\operatorname{Im}z_{2}.
\]

\corepoint{复数域\textbf{不是有序域}：两个复数（除非都是实数）“谁大谁小”没有意义，不能写 \(z_{1}<z_{2}\)；能比较大小的只有模（非负实数）。另外，\(z_{1}z_{2}=0\iff z_{1}=0\) 或 \(z_{2}=0\) 这一“无零因子”性质在复数域中依然成立，做题时可直接用。}

\subsection{复数的几何表示}
取平面直角坐标系，让复数 \(z=x+\mathrm{i}y\) 与点 \((x,y)\) 对应，此平面称为\textbf{复平面}（\(z\) 平面），\(x\) 轴、\(y\) 轴分别称为\textbf{实轴}、\textbf{虚轴}。于是“点 \(z\)”与“复数 \(z\)”不加区别；点 \(z\) 与原点 \(O\) 的连线所成向量 \(\overrightarrow{Oz}\) 也与 \(z\) 一一对应（零向量对应 \(z=0\)），故复数又可看作平面向量。

向量 \(\overrightarrow{Oz}\) 的长度称为 \(z\) 的\textbf{模}（绝对值）：
\[
|z|=r=\sqrt{x^{2}+y^{2}}\ \ (\geq 0),
\]
实轴正向转到与 \(\overrightarrow{Oz}\) 方向一致时所成的角 \(\theta\) 称为 \(z\) 的\textbf{辐角}，记 \(\operatorname{Arg}z\)。\(z\neq 0\) 时辐角有无穷多个值，彼此相差 \(2\pi\) 的整数倍；\(z=0\) 的模为零，辐角无意义。满足 \(-\pi<\arg z\leqslant \pi\) 的那一个值称为\textbf{辐角主值}，于是
\[
\operatorname{Arg}z=\arg z+2k\pi\qquad (k=0,\pm 1,\pm 2,\cdots).
\]
主值不能只写 \(\arctan\frac{y}{x}\)，须按象限分情况：

\begin{center}
\begin{tabular}{ll}
\toprule
条件 & \(\arg z\) \\
\midrule
\(x>0\) & \(\arctan\dfrac{y}{x}\) \\
\(x<0,\ y\geqslant 0\) & \(\arctan\dfrac{y}{x}+\pi\) \\
\(x<0,\ y<0\) & \(\arctan\dfrac{y}{x}-\pi\) \\
\(x=0,\ y>0\) & \(\dfrac{\pi}{2}\) \\
\(x=0,\ y<0\) & \(-\dfrac{\pi}{2}\) \\
\bottomrule
\end{tabular}
\end{center}

由直角坐标与极坐标的关系 \(x=r\cos\theta,\ y=r\sin\theta\) 得\textbf{三角表示式}
\[
z=r(\cos\theta+\mathrm{i}\sin\theta),
\]
再由欧拉公式 \(\mathrm{e}^{\mathrm{i}\theta}=\cos\theta+\mathrm{i}\sin\theta\) 得\textbf{指数表示式}
\[
z=r\,\mathrm{e}^{\mathrm{i}\theta}.
\]
两复数相等 \(\iff\) 模相等且辐角相差 \(2\pi\) 的整数倍（主值相等）。

几何上，\(z_{1}+z_{2}\) 与 \(z_{1}-z_{2}\) 分别是平行四边形（以 \(\overrightarrow{Oz_{1}},\overrightarrow{Oz_{2}}\) 为邻边）的两条对角线向量；\(|z_{1}-z_{2}|\) 就是复平面上两点 \(z_{1},z_{2}\) 之间的距离。由此得到两个基本不等式：
\keyformula{\bigl||z_{1}|-|z_{2}|\bigr|\leqslant |z_{1}\pm z_{2}|\leqslant |z_{1}|+|z_{2}|}
同时有 \(|x|\leqslant |z|\)，\(|y|\leqslant |z|\)，\(|z|\leqslant |x|+|y|\)，即
\[
|\operatorname{Re}z|\leqslant |z|,\qquad |\operatorname{Im}z|\leqslant |z|.
\]

\subsection{共轭复数}
实部相同、虚部互为相反数的两个复数称为\textbf{共轭复数}，\(z=x+\mathrm{i}y\) 的共轭记作
\[
\bar z=x-\mathrm{i}y .
\]
共轭运算与四则运算可交换次序：
\[
\overline{z_{1}\pm z_{2}}=\bar z_{1}\pm\bar z_{2},\qquad
\overline{z_{1}z_{2}}=\bar z_{1}\bar z_{2},\qquad
\overline{\left(\frac{z_{1}}{z_{2}}\right)}=\frac{\bar z_{1}}{\bar z_{2}}\quad(z_{2}\neq 0).
\]
并且
\keyformula{z+\bar z=2\operatorname{Re}z,\qquad z-\bar z=2\mathrm{i}\operatorname{Im}z,\qquad z\bar z=|z|^{2}}
由最后一式，除法可写成 \(\dfrac{z_{1}}{z_{2}}=\dfrac{z_{1}\bar z_{2}}{|z_{2}|^{2}}\)，这是实际计算中最常用的形式。又 \(z\) 为实数 \(\iff z=\bar z\)。利用共轭还可证明 \(|z_{1}z_{2}|=|z_{1}||z_{2}|\)、\(|z_{1}/z_{2}|=|z_{1}|/|z_{2}|\)（模是可乘的）。

\subsection{乘除、乘方与开方}
设 \(z_{1}=r_{1}\mathrm{e}^{\mathrm{i}\theta_{1}}\)，\(z_{2}=r_{2}\mathrm{e}^{\mathrm{i}\theta_{2}}\)，则
\keyformula{z_{1}z_{2}=r_{1}r_{2}\,\mathrm{e}^{\mathrm{i}(\theta_{1}+\theta_{2})},\qquad
\frac{z_{1}}{z_{2}}=\frac{r_{1}}{r_{2}}\,\mathrm{e}^{\mathrm{i}(\theta_{1}-\theta_{2})}\quad(z_{2}\neq 0)}
即：\textbf{两复数相乘，模相乘、辐角相加；相除则模相除、辐角相减}。几何意义是“伸缩 \(+\) 旋转”：乘 \(z_{2}\) 把 \(z_{1}\) 的模放大 \(r_{2}\) 倍并逆时针旋转 \(\theta_{2}\)。

反复应用乘法法则即得\textbf{棣莫弗（De Moivre）公式}
\[
z^{n}=r^{n}\mathrm{e}^{\mathrm{i}n\theta}=r^{n}(\cos n\theta+\mathrm{i}\sin n\theta),\qquad n=1,2,\cdots
\]
以及 \((\cos\theta+\mathrm{i}\sin\theta)^{n}=\cos n\theta+\mathrm{i}\sin n\theta\)。

\textbf{开方}是乘方的逆运算：求 \(w\) 使 \(w^{n}=z\)（\(z\neq 0\)）。设 \(w=\rho\mathrm{e}^{\mathrm{i}\varphi}\)，则 \(\rho^{n}=r\)、\(n\varphi=\theta+2k\pi\)，故
\keyformula{w_{k}=\sqrt[n]{r}\;\mathrm{e}^{\mathrm{i}\frac{\theta+2k\pi}{n}},\qquad k=0,1,2,\cdots,n-1}
非零复数开 \(n\) 次方恰有 \(n\) 个互不相同的值，它们的模都等于 \(\sqrt[n]{r}\)，辐角依次相差 \(\frac{2\pi}{n}\)，即均匀分布在以原点为心、半径 \(\sqrt[n]{r}\) 的圆周上（内接正 \(n\) 边形的顶点）。特别地，\(\sqrt[n]{1}\) 的 \(n\) 个值称为 \(n\) 次\textbf{单位根}，即
\[
\omega_{k}=\mathrm{e}^{\mathrm{i}\frac{2k\pi}{n}}\quad (k=0,1,\cdots,n-1),\qquad
1+z+z^{2}+\cdots+z^{n-1}=\frac{z^{n}-1}{z-1}\ (z\neq 1).
\]

\corepoint{\(\operatorname{Arg}(z_{1}z_{2})=\operatorname{Arg}z_{1}+\operatorname{Arg}z_{2}\) 只在“集合意义”下成立（两边都是无穷集合）；对\textbf{主值} \(\arg\) 一般\textbf{不成立}：取 \(z_{1}=z_{2}=-1\)，则 \(\arg(z_{1}z_{2})=\arg 1=0\)，而 \(\arg z_{1}+\arg z_{2}=2\pi\)，两者相差 \(2\pi\)。凡涉及主值的等式，都要先验证是否跨越了主值区间 \((-\pi,\pi]\) 的“割线”。

开方千万别漏根：解 \(z^{n}=a\) 时答案必须是 \(n\) 个值，写成 \(\sqrt[n]{r}\,\mathrm{e}^{\mathrm{i}\theta/n}\) 只给出一个根，这是最常见的失分点。}

\subsection{复球面与无穷远点}
在复平面上方放一个与之相切于原点的球面 \(S\)（如取球面 \(x^{2}+y^{2}+(u-1)^{2}=1\)，北极点 \(N=(0,0,2)\)），作\textbf{球极投影}：复平面上任一点 \(z\) 与 \(N\) 的连线交球面于唯一异于 \(N\) 的点 \(P(z)\)。这样得到 \(\mathbb{C}\) 到 \(S\setminus\{N\}\) 的一一对应，且该对应是连续同胚（逆映射也连续）。

再在 \(\mathbb{C}\) 中添加一个\textbf{无穷远点} \(\infty\)，规定 \(z=\infty\) 对应北极 \(N\)，便得
\[
\mathbb{C}_{\infty}=\mathbb{C}\cup\{\infty\},
\]
它与整个球面 \(S\) 一一对应，故称 \(S\) 为\textbf{复球面}，\(\mathbb{C}_{\infty}\) 为\textbf{扩充复平面}。

\begin{itemize}[leftmargin=2em]
  \item \(z\to\infty\) 指 \(|z|\to+\infty\)，与趋近方向无关；因此扩充复平面中\textbf{只有一个} \(\infty\)，没有 \(+\infty\)、\(-\infty\)、\(\mathrm{i}\infty\) 之分。
  \item \(\infty\) 的 \(\delta\) 邻域规定为 \(\{z: |z|>M\}\cup\{\infty\}\)（\(M=1/\delta>0\)），其去心邻域即 \(\{z: |z|>M\}\)。
  \item 复平面上任一直线（作为点集）都以 \(\infty\) 为聚点；无界区域在球面上以 \(\infty\) 为边界点或内点。
  \item 运算约定：\(\infty\pm z=\infty\)，\(\infty\cdot z=\infty\ (z\neq 0)\)，\(\frac{z}{\infty}=0\)，\(\frac{\infty}{0}=\infty\)；而 \(\infty\pm\infty\)、\(0\cdot\infty\)、\(\frac{0}{0}\)、\(\frac{\infty}{\infty}\) 无意义。
\end{itemize}

\section{复平面上的点集}

复变函数的定义域是复平面（或复球面）上的点集，故先约定点集的术语。

\subsection{基本概念}
设 \(z_{0}\in\mathbb{C}\)，\(\delta>0\)，称
\[
B(z_{0},\delta)=\{z: |z-z_{0}|<\delta\},\qquad
B^{\circ}(z_{0},\delta)=\{z: 0<|z-z_{0}|<\delta\}
\]
为 \(z_{0}\) 的 \(\delta\) \textbf{邻域}与\textbf{去心 \(\delta\) 邻域}。

设 \(E\subset\mathbb{C}\)：
\begin{itemize}[leftmargin=2em]
  \item \(z\) 是 \(E\) 的\textbf{内点}：存在 \(\delta>0\) 使 \(B(z,\delta)\subset E\)；\(E\) 的内点全体称为 \(E\) 的\textbf{内部}。
  \item \(E\) 为\textbf{开集}：\(E\) 的每一点都是内点。
  \item \(z\) 是 \(E\) 的\textbf{边界点}：对任意 \(\delta>0\)，\(B(z,\delta)\) 内既有 \(E\) 的点又有不属于 \(E\) 的点；边界点全体记 \(\partial E\)。
  \item \(E\) 的\textbf{闭包} \(\overline{E}=E\cup\partial E\)；若 \(\overline{E}=E\)（即 \(E\) 含其全部边界点）则 \(E\) 为\textbf{闭集}。
  \item \(z\) 是 \(E\) 的\textbf{聚点}（极限点）：对任意 \(\delta>0\)，\(B^{\circ}(z,\delta)\cap E\neq\varnothing\)。
  \item \(z\) 是 \(E\) 的\textbf{孤立点}：\(z\in E\) 且存在 \(\delta>0\) 使 \(B^{\circ}(z,\delta)\cap E=\varnothing\)。
  \item \(E\) \textbf{有界}：存在 \(M>0\) 使 \(|z|\leqslant M\) 对一切 \(z\in E\) 成立。
\end{itemize}

\begin{center}
\begin{tabular}{ll}
\toprule
命题 & 内容 \\
\midrule
内点与聚点 & \(E\) 的内点必是 \(E\) 的聚点 \\
孤立点与边界 & \(E\) 的孤立点必是 \(E\) 的边界点，且不是聚点 \\
闭集判别 & \(E\) 闭 \(\iff E\) 含其一切聚点 \\
\bottomrule
\end{tabular}
\end{center}

点 \(z_{0}\) 到集合 \(E\) 的\textbf{距离}定义为 \(\operatorname{dist}(z_{0},E)=\inf\{|z-z_{0}|: z\in E\}\)。

\subsection{区域和曲线}
\textbf{区域}（开区域）指\textbf{连通的开集}，其中“连通”指集合内任意两点可用集合内的折线连接。区域 \(D\) 连同其边界 \(\partial D\) 称为\textbf{闭区域} \(\overline{D}=D\cup\partial D\)。含于某个圆内的集合为有界集，相应的区域称为有界区域，否则为无界区域。

\textbf{曲线}是闭区间 \([a,b]\) 到复平面的连续映射
\[
C:\ z=z(t)=x(t)+\mathrm{i}y(t),\qquad a\leqslant t\leqslant b,
\]
\(z(a)\)、\(z(b)\) 称为起点与终点。若 \(z(t)\) 在 \([a,b]\) 上除端点外处处不同，称为\textbf{简单曲线}（无重点）；若还有 \(z(a)=z(b)\)，称为\textbf{简单闭曲线}。若 \(z'(t)\) 在 \([a,b]\) 上连续且处处不为零，称 \(C\) 为\textbf{光滑曲线}；有限条光滑曲线首尾相接称为\textbf{逐段光滑曲线}。

\textbf{若尔当（Jordan）定理}\quad 一条简单闭曲线把复平面分成内部（有界区域）与外部（无界区域）两部分，曲线本身是二者的公共边界。

\textbf{单连通区域}：区域 \(D\) 内任一条简单闭曲线的内部都含于 \(D\)。圆盘、半平面、带形区域、角形区域都是单连通区域；去心圆盘、圆环是多连通区域（其边界由两条以上互不相交的闭曲线组成）。

对区域的边界定向：沿边界前进时\textbf{区域始终在左侧}的方向为\textbf{正向}。于是单连通区域边界逆时针为正；圆环区域的外边界逆时针为正、内边界顺时针为正。

\corepoint{“区域”必须同时满足\textbf{开}与\textbf{连通}，缺一不可：\(\{z: |z|\leqslant 1\}\) 不开，两个不相交圆盘的并集不连通，都不是区域。求解析区域、判断积分能否用柯西定理时，第一步永远是看清定义域是开区域还是闭区域、单连通还是多连通——多连通区域上“绕洞”的积分一般不为零。}

\section{复变函数}

\subsection{定义与几何意义}
设 \(E\subset\mathbb{C}\)。若对每个 \(z\in E\)，按某一法则 \(f\) 有确定的复数 \(w\) 与之对应，则称 \(w=f(z)\) 是定义在 \(E\) 上的（单值）\textbf{复变函数}，\(E\) 为定义域，\(\{f(z): z\in E\}\) 为值域。若对每个 \(z\in E\) 对应的是一个复数集合 \(f(z)\)，则称 \(f\) 为\textbf{多值函数}。多值函数 \(f\) 在 \(E\) 上的\textbf{单值分支}指这样的单值函数 \(f_{0}\)：对每个 \(z\in E\) 都有 \(f_{0}(z)\in f(z)\)。

记 \(z=x+\mathrm{i}y\)，\(w=u+\mathrm{i}v\)，则
\[
w=f(z)=u(x,y)+\mathrm{i}v(x,y),
\]
即一个复变函数等价于两个实值二元函数 \(u,v\)。

\textbf{几何意义}：\(w=f(z)\) 把 \(z\) 平面上的点集 \(E\) 映为 \(w\) 平面上的点集 \(f(E)\)，称为\textbf{映射}（变换）。若 \(f\) 是 \(E\) 到 \(f(E)\) 的一一对应，则存在\textbf{反函数} \(z=f^{-1}(w)\)；两个映射可复合。研究映射 \(w=f(z)\) 时，通常同时画出 \(z\) 平面与 \(w\) 平面，看“把什么图形映成什么图形”——这正是第 6 章保形映射的出发点。

\subsection{极限与连续性}
设 \(f\) 在 \(z_{0}\) 的去心邻域内有定义。若存在复数 \(A\)，使对任意 \(\varepsilon>0\)，存在 \(\delta>0\)，当 \(0<|z-z_{0}|<\delta\) 时恒有
\[
|f(z)-A|<\varepsilon,
\]
则称 \(z\to z_{0}\) 时 \(f(z)\) 以 \(A\) 为\textbf{极限}，记 \(\lim\limits_{z\to z_{0}}f(z)=A\)。这里 \(z_{0}\) 与 \(A\) 都可以是 \(\infty\)（此时把相应的绝对值不等式换成 \(|z|>M\) 或 \(|f(z)|>M\)）。

极限的四则运算、唯一性与实函数相同；关键差别在于 \(z\to z_{0}\) 是\textbf{在平面上沿任意路径}趋近，故有\textbf{归结原则}：\(\lim\limits_{z\to z_{0}}f(z)=A\) 的充要条件是对任一趋于 \(z_{0}\) 的点列 \(\{z_{n}\}\) 都有 \(f(z_{n})\to A\)。

由于 \(f=u+\mathrm{i}v\)，极限问题可化为两个实二元函数的极限：
\keyformula{\lim_{z\to z_{0}}f(z)=A+\mathrm{i}B\iff
\lim_{(x,y)\to(x_{0},y_{0})}u(x,y)=A\ \text{且}\ \lim_{(x,y)\to(x_{0},y_{0})}v(x,y)=B}
若 \(\lim\limits_{z\to z_{0}}f(z)=f(z_{0})\)（要求 \(f\) 在 \(z_{0}\) 的某邻域内有定义），则称 \(f\) 在 \(z_{0}\) 处\textbf{连续}；在区域 \(D\) 内每点连续则称 \(f\) 在 \(D\) 上连续。

由复极限与实极限的等价性立得：
\begin{itemize}[leftmargin=2em]
  \item \(f\) 在 \(z_{0}\) 连续 \(\iff u,v\) 在 \((x_{0},y_{0})\) 连续；
  \item 连续函数的和、差、积、商（分母不为零处）与复合仍连续；
  \item 在有界闭集（紧集）上连续的复函数有界，且 \(|f(z)|\) 能取到最大值；并且一致连续。
\end{itemize}

\corepoint{判断“极限不存在”的标准手段是\textbf{沿特殊路径}：只要找到两条不同的路径使 \(f(z)\) 趋于不同的值（或沿某条路径无极限），极限即不存在。典型反例：
\[
f(z)=\frac{\bar z}{z}\ (z\neq 0),
\]
沿实轴 \(z=x\to 0\) 得 \(f\to 1\)，沿虚轴 \(z=\mathrm{i}y\to 0\) 得 \(f\to -1\)，故 \(z=0\) 处极限不存在。

另一个常考陷阱是\textbf{辐角主值不连续}：\(\arg z\) 在负实轴上的点（含 \(z<0\) 的实轴）不连续（上半平面趋近得 \(\pi\)，下半平面趋近得 \(-\pi\)），但在割开的区域 \(\mathbb{C}\setminus(-\infty,0]\) 上连续。它直接影响 \(w=\ln z\)、\(w=z^{a}\) 的解析区域，第 2 章会反复用到。}

% 子文件：4.解析函数.tex
% 对应教材：《复变函数与积分变换》(哈工大, 第三版) 第 2 章
% 对应讲义：包益欣《复变函数》第 2 章（可导、解析的定义 / 导数性质 / C-R 条件 / 初等函数导数 / 调和函数）
% 节次：§2.1 解析函数的概念（2.1.1 导数 / 2.1.2 解析的概念）
%       §2.2 函数解析的充要条件｜§2.3 解析函数与调和函数
%       §2.4 初等函数（2.4.1 指数 / 2.4.2 三角与双曲 / 2.4.3 对数 / 2.4.4 幂 / 2.4.5 反三角与反双曲）
%       §2.5 解析函数的物理意义
% 由 1.main.tex 用 \input 引入

\chapter{解析函数}

\section{解析函数的概念}

\subsection{复变函数的导数}
设 \(w=f(z)\) 在点 \(z_{0}\) 的某邻域内有定义。若极限
\[
\lim_{\Delta z\to 0}\frac{f(z_{0}+\Delta z)-f(z_{0})}{\Delta z}
\]
存在（\(\Delta z\to 0\) 指在复平面上沿任意方式趋近于零），则称 \(f(z)\) 在 \(z_{0}\) 处\textbf{可导}，该极限值称为 \(f\) 在 \(z_{0}\) 处的\textbf{导数}，记作
\[
f'(z_{0})=\left.\frac{\mathrm{d}w}{\mathrm{d}z}\right|_{z=z_{0}}.
\]
形式上它与一元实函数的导数定义完全相同，但 \(\Delta z\) 是\textbf{复}增量，可以从任何方向趋于零，这一要求比实函数强得多。

由此立得：\(f\) 在 \(z_{0}\) 可导 \(\Rightarrow f\) 在 \(z_{0}\) 连续；反之不成立。

求导法则与实函数一致（都可按定义直接验证）：
\begin{itemize}[leftmargin=2em]
  \item \((f\pm g)'=f'\pm g'\)，\((fg)'=f'g+fg'\)，\(\left(\dfrac{f}{g}\right)'=\dfrac{f'g-fg'}{g^{2}}\ (g\neq 0)\)；
  \item 复合函数：若 \(f\) 在 \(z_{0}\) 可导、\(g\) 在 \(w_{0}=f(z_{0})\) 可导，则 \((g\circ f)'(z_{0})=g'(w_{0})f'(z_{0})\)；
  \item 反函数：若 \(f\) 在 \(z_{0}\) 可导且 \(f'(z_{0})\neq 0\)，\(w_{0}=f(z_{0})\)，则反函数 \(g=f^{-1}\) 在 \(w_{0}\) 可导，且
  \[
  g'(w_{0})=\frac{1}{f'(z_{0})}.
  \]
\end{itemize}
特别地 \((z^{n})'=nz^{n-1}\)（\(n\) 为正整数），多项式与有理函数在分母不为零处可导。

\corepoint{“可导”是极强的条件：\(\Delta z\) 可沿任意方向趋于零，而实函数只要求左右两个方向。以 \(f(z)=\bar z\) 为例，沿实轴取 \(\Delta z=\Delta x\) 差商为 \(1\)，沿虚轴取 \(\Delta z=\mathrm{i}\Delta y\) 差商为 \(-1\)，故导数处处不存在，尽管其实部 \(x\)、虚部 \(-y\) 都光滑。凡含 \(\bar z\)、\(|z|\)、\(\operatorname{Re}z\)、\(\operatorname{Im}z\)、\(\arg z\) 而未与 \(z\) 组合成解析式的函数，先怀疑它不解析。}

\subsection{复变函数解析的概念}
若 \(f(z)\) 在 \(z_{0}\) 的\textbf{某一个邻域内处处可导}，则称 \(f\) 在 \(z_{0}\) 处\textbf{解析}（正则、全纯）。若 \(f\) 在区域 \(D\) 内处处可导（等价于在 \(D\) 内每点解析），称 \(f\) 在 \(D\) 内解析，或称 \(f\) 为 \(D\) 内的\textbf{解析函数}。在复平面上处处解析的函数称为\textbf{整函数}（如多项式、\(\mathrm{e}^{z}\)、\(\sin z\)）。

若 \(f\) 在 \(z_{0}\) 不解析，则称 \(z_{0}\) 为 \(f\) 的\textbf{奇点}。

三者的关系是（箭头反向均不成立）：
\[
\text{在区域内解析}\ \Longrightarrow\ \text{在某点解析}\ \Longrightarrow\ \text{在该点可导}\ \Longrightarrow\ \text{在该点连续}.
\]

\corepoint{“在某点可导”与“在该点解析”\textbf{不等价}：解析是\textbf{邻域}性质。例如 \(f(z)=|z|^{2}\) 只在 \(z=0\) 可导，而在 \(0\) 的任何邻域内都有不可导的点，故 \(f\) 在 \(z=0\) 处\textbf{不解析}，\(z=0\) 是它的奇点。判卷时“可导但无处解析”是高频陷阱。}

\section{函数解析的充要条件}

设 \(f(z)=u(x,y)+\mathrm{i}v(x,y)\) 定义在区域 \(D\) 内。

\textbf{定理（可导的必要条件）}\quad 若 \(f\) 在 \(z_{0}=x_{0}+\mathrm{i}y_{0}\) 处可导，则 \(u,v\) 在 \((x_{0},y_{0})\) 处的一阶偏导数都存在，且满足\textbf{柯西—黎曼（Cauchy--Riemann）条件}
\keyformula{\frac{\partial u}{\partial x}=\frac{\partial v}{\partial y},\qquad
\frac{\partial u}{\partial y}=-\frac{\partial v}{\partial x}}
此时
\keyformula{f'(z_{0})=\frac{\partial u}{\partial x}+\mathrm{i}\frac{\partial v}{\partial x}
=\frac{\partial v}{\partial y}-\mathrm{i}\frac{\partial u}{\partial y}}
推导思路：分别沿实轴方向与虚轴方向计算差商的极限，两者必须相等，比较实部、虚部即得 C-R 条件。

\textbf{定理（可导的充要条件）}\quad \(f\) 在 \(z_{0}\) 处可导 \(\iff u,v\) 在 \((x_{0},y_{0})\) 处\textbf{可微}，且在该点满足 C-R 条件。

\textbf{推论（最常用的判别法）}\quad 若 \(u,v\) 在 \((x_{0},y_{0})\) 的某邻域内有\textbf{连续}的一阶偏导数，且在该点满足 C-R 条件，则 \(f\) 在 \(z_{0}\) 可导；若这些条件在区域 \(D\) 内处处成立，则 \(f\) 在 \(D\) 内解析。

注意：C-R 条件本身只是\textbf{必要}条件，不是充分的——偏导数存在 + C-R 并不能保证可微，须补“偏导连续”或“可微”。

\subsection*{判定步骤}
\begin{enumerate}[leftmargin=2em]
  \item 把 \(f\) 化成 \(x,y\) 的显式：\(u=\operatorname{Re}f(z)\)，\(v=\operatorname{Im}f(z)\)（切忌只写 \(f(z)=u+\mathrm{i}v\) 而不分出 \(u,v\)）。
  \item 求出四个偏导 \(u_{x},u_{y},v_{x},v_{y}\)。
  \item 令 \(u_{x}=v_{y}\)、\(u_{y}=-v_{x}\)，解出\textbf{满足 C-R 条件的点集} \(E\)。
  \item 若 \(E\) 无内点（常见情形：只是若干孤立点或一条曲线），则 \(f\) 只在 \(E\) 上可导，\textbf{处处不解析}。
  \item 若 \(E\) 含某个区域且偏导连续，则 \(f\) 在该区域内解析，并用 \(f'=u_{x}+\mathrm{i}v_{x}\) 求导。
\end{enumerate}

\textbf{极坐标形式}\quad 记 \(z=r\mathrm{e}^{\mathrm{i}\theta}\)，\(u=u(r,\theta)\)，\(v=v(r,\theta)\)，则 C-R 条件为
\keyformula{\frac{\partial u}{\partial r}=\frac{1}{r}\frac{\partial v}{\partial \theta},\qquad
\frac{\partial v}{\partial r}=-\frac{1}{r}\frac{\partial u}{\partial \theta},\qquad
f'(z)=\mathrm{e}^{-\mathrm{i}\theta}\left(\frac{\partial u}{\partial r}+\mathrm{i}\frac{\partial v}{\partial r}\right)}
它在处理含 \(z\bar z\)、\(|z|\)、\(\arg z\) 的函数（如 \(f(z)=z^{2}\bar z\)、对数与幂函数）时比直角坐标方便。

\textbf{两个常用推论}
\begin{itemize}[leftmargin=2em]
  \item 若 \(f\) 在区域 \(D\) 内解析且 \(f'(z)\equiv 0\)，则 \(f\) 在 \(D\) 内为常数（由 \(u_{x}=u_{y}=v_{x}=v_{y}=0\) 及区域连通性）。
  \item 若 \(f=u+\mathrm{i}v\) 在 \(D\) 内解析，则 \(u,v\) 都满足拉普拉斯方程（见 §2.3）。
\end{itemize}

\corepoint{三类典型反例要记牢，考“是否解析”几乎都由它们变形而来：\(f(z)=\bar z\)（处处不可导）；\(f(z)=|z|^{2}\)（\(u_{x}=2x,\ v_{y}=0\Rightarrow x=0\)；\(u_{y}=2y,\ v_{x}=0\Rightarrow y=0\)，只在原点可导）；\(f(z)=z\operatorname{Re}z=x^{2}+\mathrm{i}xy\)（同样只在原点可导）。共同特征：C-R 条件只在\textbf{孤立点}成立，“可导点”不构成区域，故处处不解析。}

\section{解析函数与调和函数}

\textbf{定义}\quad 若二元实函数 \(u(x,y)\) 在区域 \(D\) 内有二阶连续偏导数，且满足拉普拉斯方程
\keyformula{\Delta u=\frac{\partial^{2}u}{\partial x^{2}}+\frac{\partial^{2}u}{\partial y^{2}}=0}
则称 \(u\) 为 \(D\) 内的\textbf{调和函数}。

\textbf{定理}\quad 若 \(f(z)=u+\mathrm{i}v\) 在区域 \(D\) 内解析，则 \(u,v\) 都是 \(D\) 内的调和函数。
（解析函数有任意阶导数，故 \(u,v\) 的二阶偏导存在且连续；由 C-R 条件
\(u_{xx}=v_{yx}=v_{xy}=-u_{yy}\)，即 \(u_{xx}+u_{yy}=0\)；对 \(v\) 同理。）

\textbf{定义}\quad 设 \(u,v\) 都是区域 \(D\) 内的调和函数，且满足 C-R 条件 \(u_{x}=v_{y}\)、\(u_{y}=-v_{x}\)，则称 \(v\) 是 \(u\) 的\textbf{共轭调和函数}。此时 \(u+\mathrm{i}v\) 在 \(D\) 内解析。

\textbf{定理}\quad 若 \(D\) 是\textbf{单连通}区域，则 \(D\) 内任一调和函数 \(u\) 都有共轭调和函数 \(v\)，且在相差一个实常数的意义下唯一；从而 \(u\) 必是某个解析函数的实部。

\subsection*{由 \(u\) 求 \(v\)（构造解析函数）的步骤}
\begin{enumerate}[leftmargin=2em]
  \item 验证 \(\Delta u=0\)（不满足则 \(u\) 不是调和函数，题目无解）。
  \item 由 \(v_{y}=u_{x}\) 对 \(y\) 积分：\(v=\displaystyle\int u_{x}(x,y)\,\mathrm{d}y+\varphi(x)\)，其中 \(\varphi(x)\) 待定。
  \item 代入 \(v_{x}=-u_{y}\)，比较两端的 \(x\) 函数，解出 \(\varphi'(x)\) 并积分得 \(\varphi(x)\)。
  \item 写出 \(f=u+\mathrm{i}v\)；若要求表成 \(z\) 的函数，可用 \(x=z,\ y=0\) 化简易式（因解析函数的表达式与坐标无关）。
\end{enumerate}
等价的曲线积分写法：
\keyformula{v(x,y)=\int_{(x_{0},y_{0})}^{(x,y)}-u_{y}\,\mathrm{d}x+u_{x}\,\mathrm{d}y+C}
（被积表达式全微分存在，正是 \(\Delta u=0\)，故积分与路径无关。）

\corepoint{共轭调和是\textbf{有方向}的：\(v\) 是 \(u\) 的共轭调和函数 \(\iff u+\mathrm{i}v\) 解析；此时 \(-u\)（而非 \(u\)）才是 \(v\) 的共轭调和函数，因为 \(v-\mathrm{i}u=-\mathrm{i}(u+\mathrm{i}v)\) 解析。另外，“存在共轭调和函数”依赖区域\textbf{单连通}：圆环 \(0<|z|<1\) 上 \(u=\ln|z|\) 调和却无共轭调和函数，否则将得到 \(\ln z\) 的单值解析分支，而绕原点一周其值改变 \(2\pi\mathrm{i}\)。}

\section{初等函数}

\subsection{指数函数}
对 \(z=x+\mathrm{i}y\)，定义
\keyformula{\mathrm{e}^{z}=\mathrm{e}^{x}(\cos y+\mathrm{i}\sin y)}
它由实指数函数的泰勒展开推广而来，并保持实指数的基本性质：
\begin{itemize}[leftmargin=2em]
  \item 在全平面解析，且 \((\mathrm{e}^{z})'=\mathrm{e}^{z}\)；
  \item \(\mathrm{e}^{z_{1}+z_{2}}=\mathrm{e}^{z_{1}}\mathrm{e}^{z_{2}}\)，\(\mathrm{e}^{0}=1\)；
  \item 以 \(2\pi\mathrm{i}\) 为周期：\(\mathrm{e}^{z+2\pi\mathrm{i}}=\mathrm{e}^{z}\)；
  \item 处处不为零：\(\mathrm{e}^{z}\neq 0\)，且 \(\mathrm{e}^{z}=1\iff z=2k\pi\mathrm{i}\ (k\in\mathbb{Z})\)；
  \item \(|\mathrm{e}^{z}|=\mathrm{e}^{x}\)，\(\operatorname{Arg}\mathrm{e}^{z}=y+2k\pi\)。
\end{itemize}
与实指数函数的重要差别：\(\mathrm{e}^{z}\) 是\textbf{周期}函数，故不是一一映射，没有单值反函数（这正是对数函数多值性的根源）；又当 \(y\to\pm\infty\) 时 \(|\mathrm{e}^{z}|=\mathrm{e}^{x}\) 与 \(y\) 无关而 \(\mathrm{e}^{z}\) 沿虚方向震荡，故 \(\lim\limits_{z\to\infty}\mathrm{e}^{z}\) 不存在。

\subsection{三角函数与双曲函数}
由欧拉公式解出 \(\cos x,\sin x\) 的指数表示，再把实变量换成复变量，定义
\keyformula{\sin z=\frac{\mathrm{e}^{\mathrm{i}z}-\mathrm{e}^{-\mathrm{i}z}}{2\mathrm{i}},\qquad
\cos z=\frac{\mathrm{e}^{\mathrm{i}z}+\mathrm{e}^{-\mathrm{i}z}}{2},\qquad
\tan z=\frac{\sin z}{\cos z},\quad \cot z=\frac{\cos z}{\sin z}}
（正割、余割类似）。性质：
\begin{itemize}[leftmargin=2em]
  \item 在全平面解析，\((\sin z)'=\cos z\)，\((\cos z)'=-\sin z\)；\(\tan z,\cot z\) 在分母不为零处解析；
  \item 以 \(2\pi\) 为周期；\(\sin z\) 为奇函数、\(\cos z\) 为偶函数；和角公式、\(\sin^{2}z+\cos^{2}z=1\) 等恒等式全部成立；
  \item 零点与实情形相同：\(\sin z=0\iff z=k\pi\)，\(\cos z=0\iff z=\frac{\pi}{2}+k\pi\)（\(k\in\mathbb{Z}\)）；
  \item \textbf{无界性}：\(|\cos(\mathrm{i}y)|=\dfrac{\mathrm{e}^{y}+\mathrm{e}^{-y}}{2}\to+\infty\ (y\to\infty)\)，故 \(|\sin z|\leqslant 1\)、\(|\cos z|\leqslant 1\) 在复平面上\textbf{不再成立}。
\end{itemize}

\textbf{双曲函数}直接由指数定义：
\keyformula{\operatorname{sh}z=\frac{\mathrm{e}^{z}-\mathrm{e}^{-z}}{2},\qquad
\operatorname{ch}z=\frac{\mathrm{e}^{z}+\mathrm{e}^{-z}}{2},\qquad
\operatorname{th}z=\frac{\operatorname{sh}z}{\operatorname{ch}z}}
它们在全平面解析，\((\operatorname{sh}z)'=\operatorname{ch}z\)，\((\operatorname{ch}z)'=\operatorname{sh}z\)，以 \(2\pi\mathrm{i}\) 为周期；且
\[
\operatorname{ch}z=\cos(\mathrm{i}z),\qquad \operatorname{sh}z=-\mathrm{i}\sin(\mathrm{i}z),
\]
零点为 \(\operatorname{sh}z=0\iff z=k\pi\mathrm{i}\)，\(\operatorname{ch}z=0\iff z=\left(\frac{\pi}{2}+k\pi\right)\mathrm{i}\)。

\corepoint{复三角函数只在\textbf{实轴}上有零点，却在虚方向上无界——这与实三角函数“有界、零点周期性分布”的直觉相反。求 \(\sin z=0\) 时不要只写 \(z=k\pi\) 却忘了可以设 \(z=x+\mathrm{i}y\) 验证：由 \(\mathrm{e}^{2\mathrm{i}z}=1\) 可得 \(\mathrm{e}^{-2y}\mathrm{e}^{2\mathrm{i}x}=1\)，故 \(y=0\)，零点的确都在实轴上。}

\subsection{对数函数}
对数函数定义为指数函数的反函数：满足 \(\mathrm{e}^{w}=z\ (z\neq 0)\) 的 \(w\) 称为 \(z\) 的\textbf{对数}，记 \(w=\operatorname{Ln}z\)。设 \(w=u+\mathrm{i}v\)，则 \(\mathrm{e}^{w}=\mathrm{e}^{u}(\cos v+\mathrm{i}\sin v)=z\)，比较两端的模与辐角得
\keyformula{\operatorname{Ln}z=\ln|z|+\mathrm{i}\operatorname{Arg}z=\ln|z|+\mathrm{i}(\arg z+2k\pi),\qquad k=0,\pm 1,\pm 2,\cdots}
它是\textbf{无穷多值}函数，各值相差 \(2\pi\mathrm{i}\) 的整数倍。取 \(\operatorname{Arg}z\) 的主值即得\textbf{对数主值}
\[
\ln z=\ln|z|+\mathrm{i}\arg z\qquad (-\pi<\arg z\leqslant \pi).
\]
性质：
\begin{itemize}[leftmargin=2em]
  \item 集合意义下的运算律：\(\operatorname{Ln}(z_{1}z_{2})=\operatorname{Ln}z_{1}+\operatorname{Ln}z_{2}\)，\(\operatorname{Ln}\dfrac{z_{1}}{z_{2}}=\operatorname{Ln}z_{1}-\operatorname{Ln}z_{2}\)；
  \item 主值 \(\ln z\) 在割开平面 \(D=\mathbb{C}\setminus(-\infty,0]\)（即 \(|z|>0,\ -\pi<\arg z<\pi\)）内单值解析，且
  \[
  \frac{\mathrm{d}}{\mathrm{d}z}\ln z=\frac{1}{z};
  \]
  \item \(\operatorname{Ln}z\) 的\textbf{任一}连续单值分支 \(f(z)\) 都满足 \(f'(z)=\dfrac{1}{z}\)（不同分支只差常数 \(2k\pi\mathrm{i}\)）；
  \item 当 \(z=x>0\) 时 \(\ln z=\ln x\)，即复对数是实对数的推广；负数也有对数，如 \(\operatorname{Ln}(-1)=(2k+1)\pi\mathrm{i}\)。
\end{itemize}

\subsection{幂函数}
设 \(a\in\mathbb{C}\)，对 \(z\neq 0\) 定义
\keyformula{z^{a}=\mathrm{e}^{a\operatorname{Ln}z}=\mathrm{e}^{a(\ln|z|+\mathrm{i}\arg z+2k\pi\mathrm{i})}}
其多值性完全由 \(a\) 决定：
\begin{center}
\begin{tabular}{lll}
\toprule
\(a\) 的类型 & \(z^{a}\) 的取值个数 & 说明 \\
\midrule
正整数 \(n\) & 1（单值） & 即通常的 \(z^{n}\)，全平面解析 \\
负整数 \(-n\) & 1（单值） & \(z^{-n}=1/z^{n}\)，除 \(z=0\) 外解析 \\
既约分数 \(p/q\ (q>0)\) & \(q\)（有限多值） & \(q\) 个分支，如 \(\sqrt[n]{z}\) 有 \(n\) 个值 \\
其他（无理数或虚部非零） & 无穷多值 & 各值相差 \(2\pi a\mathrm{i}\) 的整数倍 \\
\bottomrule
\end{tabular}
\end{center}
取定 \(\operatorname{Ln}z\) 的一个单值解析分支即得 \(z^{a}\) 的一个\textbf{单值解析分支}，在该分支上
\[
(z^{a})'=a\,z^{a-1}.
\]
\(a\) 为整数时上式对通常的幂函数也成立。

\subsection{反三角函数与反双曲函数}
由三角函数的指数表示解出 \(w\)，即得反三角函数的对数表示（均为多值函数）：
\keyformula{\operatorname{Arcsin}z=-\mathrm{i}\operatorname{Ln}\left(\mathrm{i}z+\sqrt{1-z^{2}}\right),\qquad
\operatorname{Arccos}z=-\mathrm{i}\operatorname{Ln}\left(z+\sqrt{z^{2}-1}\right)}
\keyformula{\operatorname{Arctan}z=\frac{1}{2\mathrm{i}}\operatorname{Ln}\frac{1+\mathrm{i}z}{1-\mathrm{i}z}}
类似地
\[
\operatorname{Arsh}z=\operatorname{Ln}\left(z+\sqrt{z^{2}+1}\right),\quad
\operatorname{Arch}z=\operatorname{Ln}\left(z+\sqrt{z^{2}-1}\right),\quad
\operatorname{Arth}z=\frac{1}{2}\operatorname{Ln}\frac{1+z}{1-z}.
\]
它们的多值性来自两处：\textbf{对数}的多值性与\textbf{根式}的多值性（复合而成）。取定单值解析分支后有
\[
(\operatorname{Arcsin}z)'=\frac{1}{\sqrt{1-z^{2}}},\qquad
(\operatorname{Arctan}z)'=\frac{1}{1+z^{2}}.
\]

\subsection{多值函数与支点}
对多值函数 \(w=f(z)\)，若沿绕 \(z_{0}\) 一周的闭曲线连续变动时，函数值从某一分支连续变到另一分支（且绕任意小圈都如此），则称 \(z_{0}\) 为 \(f\) 的\textbf{支点}（分支点）。

\begin{center}
\begin{tabular}{ll}
\toprule
多值函数 & 支点 \\
\midrule
\(\operatorname{Ln}z\) & \(0\) 与 \(\infty\)（每绕原点一周，虚部增加 \(2\pi\)） \\
\(\sqrt[n]{z}\) & \(0\) 与 \(\infty\)（绕一周进入下一分支，绕 \(n\) 周复原） \\
\(\sqrt{(z-a)(z-b)}\) & \(a,\ b,\ \infty\)（绕 \(a\) 或 \(b\) 一周变号） \\
\bottomrule
\end{tabular}
\end{center}

在复平面上从支点引一条\textbf{支割线}（如 \(\operatorname{Ln}z\) 取负实轴），把平面割成单连通区域，多值函数在该区域上就可分出单值解析分支。保证这一点的定理是：\textbf{若 \(D\) 是不含原点的单连通区域，则 \(\operatorname{Ln}z\) 在 \(D\) 上存在无穷多个单值解析分支，彼此相差 \(2k\pi\mathrm{i}\)。}

\corepoint{易混的两个说法：\textbf{支点}是“绕行一周函数值改变”的点，与函数在该点是否有定义无关——\(z=0\) 不在 \(\operatorname{Ln}z\) 的定义域内，却是它的支点。问“\(\ln z\) 的解析区域”时，标准答案是“割去负实轴（或任一由原点出发的射线）后的单连通区域”，而不是“除原点外的平面”。}

\section{解析函数的物理意义}

解析函数的实部 \(\varphi\) 与虚部 \(\psi\) 都是调和函数且等值线正交（\(\nabla\varphi\cdot\nabla\psi=\varphi_{x}\psi_{x}+\varphi_{y}\psi_{y}=0\)），这一结构恰好对应平面稳定场“无旋且无源”的两条性质。于是三类稳态场都可用一个\textbf{复势} \(w=\varphi+\mathrm{i}\psi\)（解析函数）统一描述（教材 §2.5）：

\begin{center}
\begin{tabular}{llll}
\toprule
场 & 势函数 \(\varphi\) & 共轭调和函数 \(\psi\) & 正交的两族曲线 \\
\midrule
平面流速场 & 速度势 & 流函数 & 等势线 / 流线 \\
平面静电场 & 电势 & 力函数 & 等势线 / 电力线 \\
平面稳定温度场 & 温度 \(T\) & 热流函数 & 等温线 / 热流线 \\
\bottomrule
\end{tabular}
\end{center}

三者的势函数都满足 \(\Delta\varphi=0\)（无旋无源、无电荷、无内热源），场量由 \(w\) 的导数给出。以流速场为例，
\keyformula{v_{x}+\mathrm{i}v_{y}=\overline{w'(z)}}
即\textbf{复速度等于复势导数的共轭}；静电场中 \(\boldsymbol{E}=-\nabla\varphi\)，温度场中 \(\boldsymbol{q}=-k\nabla T\)（\(k\) 为导热系数），结构完全相同。数学上“求一个平面无旋无源场”就此化归为“求一个解析函数”。

\corepoint{三场同构，差别只是物理量的名称与符号约定。最易错的是复速度的共轭：由 \(w'=\varphi_{x}+\mathrm{i}\psi_{x}=v_{x}-\mathrm{i}v_{y}\) 可知正确式是 \(v_{x}+\mathrm{i}v_{y}=\overline{w'(z)}\)，写成 \(w'(z)\) 就把 \(v_{y}\) 的符号弄反了。}
% 子文件：5.复变函数的积分.tex
% 对应教材：《复变函数与积分变换》(哈工大, 第三版) 第 3 章
% 对应讲义：包益欣《复变函数》第 3 章（复积分的定义 / 柯西积分定理及应用 / 柯西积分公式及证明 / 其应用）
% 节次：§3.1 复变函数积分的概念（定义 / 性质 / 存在条件与计算）
%       §3.2 柯西积分定理（3.2.1 定理 / 3.2.2 不定积分 / 3.2.3 复合闭路定理）
%       §3.3 柯西积分公式（3.3.1 公式 / 3.3.2 高阶导数公式 / 3.3.3 几个重要的推论）
%       附加：§3.4 解析函数模的性质（最大模、最小模原理，据讲义 §1.4）
% 由 1.main.tex 用 \input 引入

\chapter{复变函数的积分}

\section{复变函数积分的概念}

\subsection{积分的定义}
先约定曲线。闭区间 \([a,b]\) 上的连续映射
\[
C:\ z=z(t)=x(t)+\mathrm{i}y(t),\qquad a\leqslant t\leqslant b
\]
称为复平面上由 \(z(a)\) 到 \(z(b)\) 的\textbf{有向曲线}，\(z(a)\)、\(z(b)\) 为起点与终点。称 \(z=z(a+b-t)\ (a\leqslant t\leqslant b)\) 为 \(C\) 的\textbf{反向曲线}，记 \(C^{-}\)。若 \(z'(t)\) 连续且处处不为零，称 \(C\) \textbf{光滑}；有限条光滑曲线首尾相接称为\textbf{逐段光滑}。若 \(z(a)=z(b)\) 且 \(t\in(a,b)\) 时曲线无重点，称 \(C\) 为\textbf{简单闭曲线}；不作说明时规定其\textbf{逆时针}方向为正向。

设 \(C\) 是逐段光滑有向曲线，\(f(z)\) 在 \(C\) 上有定义。在 \([a,b]\) 上插入分点
\[
a=t_{0}<t_{1}<\cdots<t_{n}=b,
\]
记 \(z_{k}=z(t_{k})\)，\(\Delta z_{k}=z_{k}-z_{k-1}\)，在每个小弧 \(z_{k-1}z_{k}\) 上任取一点 \(\zeta_{k}\)，作\textbf{黎曼和}
\[
\sum_{k=1}^{n}f(\zeta_{k})\,\Delta z_{k}.
\]
令 \(\lambda=\max\limits_{1\leqslant k\leqslant n}|\Delta z_{k}|\)，若 \(\lambda\to 0\) 时该和式的极限存在且与分法与 \(\zeta_{k}\) 的取法无关，则称此极限为 \(f(z)\) 沿 \(C\) 的\textbf{复积分}，记
\[
\int_{C}f(z)\,\mathrm{d}z=\lim_{\lambda\to 0}\sum_{k=1}^{n}f(\zeta_{k})\,\Delta z_{k},
\]
\(C\) 为积分路径，\(f\) 为被积函数。沿闭曲线的积分记作 \(\oint_{C}f(z)\,\mathrm{d}z\)。

\subsection{积分的性质}
设所涉积分都存在，则
\begin{enumerate}[leftmargin=2em]
  \item \textbf{线性性}：\(\displaystyle\int_{C}\bigl[\alpha f(z)+\beta g(z)\bigr]\mathrm{d}z=\alpha\int_{C}f\,\mathrm{d}z+\beta\int_{C}g\,\mathrm{d}z\)（\(\alpha,\beta\) 为复常数）。
  \item \textbf{反向性}：\(\displaystyle\int_{C^{-}}f(z)\,\mathrm{d}z=-\int_{C}f(z)\,\mathrm{d}z\)。
  \item \textbf{路径可加性}：若 \(C\) 由 \(C_{1}\) 与 \(C_{2}\) 首尾相接而成，则 \(\displaystyle\int_{C}f\,\mathrm{d}z=\int_{C_{1}}f\,\mathrm{d}z+\int_{C_{2}}f\,\mathrm{d}z\)。
  \item \textbf{估值不等式}：设 \(|f(z)|\leqslant M\ (z\in C)\)，\(L\) 为 \(C\) 的长度，则
  \[
  \left|\int_{C}f(z)\,\mathrm{d}z\right|\leqslant\int_{C}|f(z)|\,|\mathrm{d}z|=\int_{C}|f(z)|\,\mathrm{d}s\leqslant ML .
  \]
\end{enumerate}

\corepoint{估值不等式是复积分里\emph{唯一}的“放缩工具”，凡证极限为零、证积分有界、证刘维尔定理都靠它。使用时两个量都要算准：\(M\) 必须是 \(|f|\) 在 \(C\) 上的\textbf{上界}（常取最大值），\(L\) 是路径长度（圆周为 \(2\pi R\)）。注意被积函数在路径上含参数时 \(M\) 也随参数变，要一并估计。}

\subsection{积分的存在条件与计算}
\textbf{定理}\quad 若 \(f(z)\) 在逐段光滑曲线 \(C\) 上连续，则 \(\displaystyle\int_{C}f(z)\,\mathrm{d}z\) 存在。

\textbf{计算方法一（参数化，最基本）}\quad 设 \(C:z=z(t)\ (a\leqslant t\leqslant b)\) 光滑，则
\keyformula{\int_{C}f(z)\,\mathrm{d}z=\int_{a}^{b}f\bigl[z(t)\bigr]\,z'(t)\,\mathrm{d}t}
逐段光滑时按段积分再相加。这一步把复积分化为定积分（被积函数为复值函数，实部虚部分别积分）。

\textbf{计算方法二（化为两个第二类线积分）}\quad 设 \(f=u+\mathrm{i}v\)，\(z=x+\mathrm{i}y\)，则
\keyformula{\int_{C}f(z)\,\mathrm{d}z=\int_{C}(u\,\mathrm{d}x-v\,\mathrm{d}y)+\mathrm{i}\int_{C}(v\,\mathrm{d}x+u\,\mathrm{d}y)}
它说明复积分存在 \(\iff\) 右端两个线积分都存在；也说明复积分与路径方向有关。

\textbf{两个基本结果}（后面反复使用）
\keyformula{\int_{C}\mathrm{d}z=z(b)-z(a);\qquad
\oint_{|z-z_{0}|=r}(z-z_{0})^{n}\,\mathrm{d}z=
\begin{cases}2\pi\mathrm{i}, & n=-1,\\ 0, & n\neq -1,\ n\in\mathbb{Z}\end{cases}}
前者由参数化直接得到（\(\int_{a}^{b}z'(t)\mathrm{d}t=z(b)-z(a)\)），后者把 \(z=z_{0}+r\mathrm{e}^{\mathrm{i}\theta}\) 代入即可验证。它是“绕奇点积分不为零”的根源。

\section{柯西积分定理}

\subsection{柯西积分定理}
\textbf{定理（柯西积分定理）}\quad 设 \(f(z)\) 在\textbf{单连通区域} \(D\) 内解析，\(C\) 为 \(D\) 内任一条逐段光滑闭曲线，则
\keyformula{\oint_{C}f(z)\,\mathrm{d}z=0}
常用的等价叙述：若 \(f\) 在简单闭曲线 \(C\) 上及其内部解析，则沿 \(C\) 的积分为零。

\textbf{条件不能削弱}：只要求 \(f\) 连续不够，必须解析；区域必须是单连通的，否则绕“洞”的积分一般不为零——在去心圆盘 \(0<|z|<1\) 内 \(f(z)=1/z\) 解析，但 \(\oint_{|z|=1}\mathrm{d}z/z=2\pi\mathrm{i}\neq 0\)。

\textbf{推论（积分与路径无关）}\quad 若 \(f\) 在单连通区域 \(D\) 内解析，则 \(\int_{C}f\,\mathrm{d}z\) 只与起点、终点有关，与路径无关。

\textbf{闭路变形原理}\quad 在积分值不变的条件下，闭曲线可以在区域内部连续变形，只要变形过程中\textbf{不越过奇点}。它是复合闭路定理的几何表述，也是“挖小圆”技巧的依据。

\subsection{不定积分}
\textbf{定理（原函数存在定理）}\quad 设 \(f(z)\) 在单连通区域 \(D\) 内解析，固定 \(z_{0}\in D\)，令
\[
F(z)=\int_{z_{0}}^{z}f(\xi)\,\mathrm{d}\xi
\]
（积分沿 \(D\) 内任一路径），则 \(F(z)\) 在 \(D\) 内解析，且 \(F'(z)=f(z)\)。称 \(F\) 为 \(f\) 的\textbf{原函数}，\(F(z)+C\)（\(C\) 为复常数）为 \(f\) 的\textbf{不定积分}，记 \(\int f(z)\mathrm{d}z=F(z)+C\)。同一函数的两个原函数相差常数。

\textbf{牛顿—莱布尼茨公式}\quad 若 \(f\) 在单连通区域 \(D\) 内解析，\(F\) 是 \(f\) 的任一原函数，则对 \(z_{0},z_{1}\in D\)
\keyformula{\int_{z_{0}}^{z_{1}}f(z)\,\mathrm{d}z=F(z_{1})-F(z_{0})}
于是实积分的换元、分部积分等技巧在复积分中照旧可用（只要原函数在区域内单值解析）。

\textbf{推论}\quad 若 \(D\) 是不含零点的单连通区域，则 \(\operatorname{Ln}z\) 在 \(D\) 上存在单值解析分支，且
\[
F(z)=\ln z_{0}+\int_{z_{0}}^{z}\frac{\mathrm{d}\xi}{\xi}
\]
就是其中一个分支。更一般地，若 \(f\) 在单连通区域 \(D\) 内解析且处处不为零，则 \(\operatorname{Ln}f(z)\) 在 \(D\) 内存在单值解析分支（取 \(F'=f'/f\)）。

\subsection{复合闭路定理}
设多连通区域 \(D\) 的边界由外边界 \(C_{0}\) 与内边界 \(C_{1},C_{2},\cdots,C_{n}\) 组成，均取\textbf{正向}（沿边界前进时区域在左侧，故外边界逆时针、内边界顺时针），记复合闭路 \(C=C_{0}+C_{1}^{-}+\cdots+C_{n}^{-}\)。若 \(f\) 在 \(D\) 上解析，则
\keyformula{\oint_{C}f(z)\,\mathrm{d}z=0,\qquad\text{即}\qquad
\oint_{C_{0}}f\,\mathrm{d}z=\sum_{k=1}^{n}\oint_{C_{k}}f\,\mathrm{d}z}
（右端各内边界均取逆时针方向。）直观地说：绕外边界一周的积分等于绕各个“洞”一周的积分之和。

\corepoint{用复合闭路定理解题时\textbf{三个位置必须交代清楚}：(1) 被积函数的奇点在哪；(2) 积分曲线围住了哪些奇点；(3) 挖出的小圆周方向是逆时针还是顺时针。教材写法是“在 \(C\) 内作互不相交的小圆周 \(C_{1},\dots,C_{n}\) 分别包围各奇点”，此时 \(\oint_{C}=\sum\oint_{C_{k}}\)（都取逆时针）；若把小圆当作区域边界，则内边界是\textbf{顺时针}，符号要反过来。方向弄反是这一章最常见的错误。}

\section{柯西积分公式}

\subsection{柯西积分公式}
\textbf{定理（柯西积分公式）}\quad 设 \(f(z)\) 在简单闭曲线 \(C\) 上及其内部 \(D\) 内解析，\(z_{0}\) 为 \(D\) 内任一点，则
\keyformula{f(z_{0})=\frac{1}{2\pi\mathrm{i}}\oint_{C}\frac{f(z)}{z-z_{0}}\,\mathrm{d}z}
等价写法：\(\displaystyle\oint_{C}\frac{f(z)}{z-z_{0}}\mathrm{d}z=2\pi\mathrm{i}f(z_{0})\)。

证明要点：以 \(z_{0}\) 为心作充分小的圆 \(K:|z-z_{0}|=\rho\) 含于 \(D\)，由复合闭路定理把沿 \(C\) 的积分化为沿 \(K\) 的积分，再用 \(f\) 的连续性与估值不等式证明
\(\oint_{K}\dfrac{f(z)}{z-z_{0}}\mathrm{d}z\to 2\pi\mathrm{i}f(z_{0})\)（\(\rho\to 0\)）。

\textbf{意义}\quad 解析函数在区域\textbf{内部}的值完全由它在\textbf{边界}上的值决定。由此立得：若两个解析函数在区域边界上处处相等，则它们在整个区域上恒等（解析函数的“刚性”，实函数没有此性质）。

公式对复合闭路（多连通区域）同样成立：只要 \(z_{0}\) 在 \(C\) 所围区域内部，而 \(f\) 在 \(C\) 上及其所围区域内解析即可。

\subsection{高阶导数公式}
\textbf{定理}\quad 若 \(f(z)\) 在简单闭曲线 \(C\) 上及其内部 \(D\) 内解析，则 \(f\) 在 \(D\) 内有\textbf{任意阶导数}，且对 \(D\) 内任一点 \(z_{0}\)
\keyformula{f^{(n)}(z_{0})=\frac{n!}{2\pi\mathrm{i}}\oint_{C}\frac{f(z)}{(z-z_{0})^{n+1}}\,\mathrm{d}z,
\qquad n=1,2,\cdots}
证明思路：对 \(f(z_{0})=\frac{1}{2\pi\mathrm{i}}\oint_{C}\frac{f(z)}{z-z_{0}}\mathrm{d}z\) 关于 \(z_{0}\) 求导，用估值不等式证明“积分号下可以求导”（差商与导数的差为 \(O(|\Delta z|)\)），再对 \(n\) 归纳。

反用此式是计算积分的常用手段：
\[
\oint_{C}\frac{f(z)}{(z-z_{0})^{n+1}}\,\mathrm{d}z=\frac{2\pi\mathrm{i}}{n!}f^{(n)}(z_{0}).
\]

\corepoint{这一公式带来复分析与实分析最本质的差别：\textbf{解析 \(\Rightarrow\) 无穷次可导}（而实函数一次可导推不出二次可导）。使用时注意三点：(1) 分子 \(f\) 必须在 \(C\) 及内部\textbf{解析}——奇点要留在分母外面；(2) 分母的幂次是 \(n+1\)，对应 \(n!\)；(3) 若 \(C\) 内有多个奇点，先挖洞拆分，使每个积分只含一个奇点再套公式。}

\subsection{几个重要的推论}
\subsubsection*{（1）平均值公式}
设 \(f\) 在以 \(z_{0}\) 为心、\(R\) 为半径的圆周 \(C\) 上及其内部解析，则
\keyformula{f(z_{0})=\frac{1}{2\pi}\int_{0}^{2\pi}f\bigl(z_{0}+R\mathrm{e}^{\mathrm{i}\theta}\bigr)\,\mathrm{d}\theta
=\frac{1}{2\pi R}\oint_{C}f(z)\,\mathrm{d}s}
即解析函数在圆心的值等于它在圆周上值的平均。

\subsubsection*{（2）柯西不等式}
设 \(M=\max\limits_{z\in C}|f(z)|\)，则
\keyformula{\left|f^{(n)}(z_{0})\right|\leqslant\frac{n!\,M}{R^{n}};\qquad
\text{特别地}\ \left|f'(z_{0})\right|\leqslant\frac{M}{R}}
它把函数在一点处的各阶导数用它在圆周上的最大模控制住，是估计导数、证明刘维尔定理的工具。

\subsubsection*{（3）刘维尔（Liouville）定理}
\textbf{整函数}指在复平面上处处解析的函数（多项式、\(\mathrm{e}^{z}\)、\(\sin z\) 等）。若整函数 \(f(z)\) 在复平面上\textbf{有界}（\(|f(z)|\leqslant M\)），则 \(f\) 必为常数。

证明：对任意 \(z\) 与任意 \(R>0\)，由柯西不等式 \(|f'(z)|\leqslant M/R\)；令 \(R\to+\infty\) 得 \(f'(z)=0\)，故 \(f\) 为常数。

\subsubsection*{（4）代数基本定理}
任意 \(n\) 次（\(n\geqslant 1\)）复系数多项式
\[
f(z)=a_{0}z^{n}+a_{1}z^{n-1}+\cdots+a_{n-1}z+a_{n}\qquad (a_{0}\neq 0)
\]
在复平面上\textbf{至少有一个零点}，从而恰有 \(n\) 个零点（计重数）。

证明：若 \(f\) 无零点，则 \(F=1/f\) 是整函数，且由 \(|f(z)|\to+\infty\ (z\to\infty)\) 知 \(F\) 有界，由刘维尔定理 \(F\) 为常数，矛盾。

\subsubsection*{（5）莫累拉（Morera）定理}
设 \(f(z)\) 在区域 \(D\) 内连续，且对 \(D\) 内任一条其内部含于 \(D\) 的简单闭曲线 \(C\) 都有 \(\oint_{C}f\,\mathrm{d}z=0\)，则 \(f\) 在 \(D\) 内解析。

它是柯西积分定理的\textbf{逆定理}：柯西定理用“解析”推出“积分为零”，莫累拉定理用“积分为零”反推“解析”，常用于证明由积分或级数定义的函数是解析的。

\subsubsection*{（6）泊松（Poisson）积分公式}
设 \(f\) 在圆 \(C:|z-z_{0}|=R\) 内及 \(C\) 上解析，\(z=z_{0}+r\mathrm{e}^{\mathrm{i}\theta}\ (r<R)\) 为圆内一点，则
\[
f(z)=\frac{1}{2\pi}\int_{0}^{2\pi}
\frac{R^{2}-r^{2}}{R^{2}-2Rr\cos(\varphi-\theta)+r^{2}}\,
f\bigl(z_{0}+R\mathrm{e}^{\mathrm{i}\varphi}\bigr)\,\mathrm{d}\varphi .
\]
比较两端的实部即得\textbf{泊松积分公式}：圆内调和函数的值可由它在圆周上的值通过上述积分表示。特别取 \(r=0\) 即回到调和函数的平均值公式。

\section{解析函数模的性质}

下面的结论由平均值公式与柯西不等式得到（讲义 §1.4），它刻画了解析函数模的“边界控制”性质。

\textbf{最大模原理}\quad 设 \(D\) 是有界连通区域，\(f(z)\) 在 \(\overline{D}\) 上连续、在 \(D\) 内解析，则
\keyformula{\max_{z\in\overline{D}}|f(z)|=\max_{z\in\partial D}|f(z)|=M}
并且：若存在 \(D\) 的\textbf{内点} \(z_{0}\) 使 \(|f(z_{0})|=M\)，则 \(f\) 在 \(D\) 内为常数。

证明要点：设内点取到最大值 \(M\)，由平均值公式
\(M=|f(z_{0})|\leqslant\frac{1}{2\pi}\int_{0}^{2\pi}|f(z_{0}+r\mathrm{e}^{\mathrm{i}\theta})|\mathrm{d}\theta\leqslant M\)，
故等号处处成立、\(|f|\) 在圆内恒为 \(M\)；沿折线逐点推过去，得 \(|f|\) 在 \(D\) 内恒为 \(M\)，再由 §2.2 的推论知 \(f\) 为常数。

\textbf{最小模原理}\quad 设 \(D\) 有界连通，\(f\) 在 \(\overline{D}\) 上连续、在 \(D\) 内解析且\textbf{处处不为零}，则
\[
\min_{z\in\overline{D}}|f(z)|=\min_{z\in\partial D}|f(z)|;
\]
若在内点取到最小值，则 \(f\) 为常数。（对 \(1/f\) 用最大模原理。）

\textbf{调和函数的极值原理}\quad 若 \(u(x,y)\) 是有界闭区域 \(\overline{D}\) 上的调和函数（在 \(D\) 内调和、在 \(\overline{D}\) 上连续），则 \(u\) 的最大值与最小值都在边界 \(\partial D\) 上取到；若在内点取到，则 \(u\) 为常数。

\corepoint{最大模原理管的是 \(|f|\) 而\textbf{不是} \(f\)：\(f\) 本身（复值）谈不上大小。三个使用前提别漏：区域\textbf{有界}、\(f\) 在闭区域上\textbf{连续}、在内部\textbf{解析}。典型应用：证明“模为常数的解析函数是常数”（\(|f|\equiv c\) 时内部也取到最大值）。对照两个二元实函数能看清“调和”的关键：\(u=x^{2}+y^{2}\) 在圆盘上的最大值在圆心取到（\(\Delta u=4\neq0\)，不调和，不受定理约束）；而调和的 \(u=x^{2}-y^{2}=\operatorname{Re}(z^{2})\) 的最大值必在边界上取到。}

% 子文件：6.级数.tex
% 对应教材：《复变函数与积分变换》(哈工大, 第三版) 第 4 章（§4.5 孤立奇点分类取自教材第 5 章 §5.1）
% 对应讲义：包益欣《复变函数》第 4 章（数项级数与函数项级数 / 幂级数与泰勒展开 / 广义幂级数与洛朗展开）
% 节次：§4.1 复变函数项级数（4.1.1 复数序列 / 4.1.2 复数项级数 / 4.1.3 复变函数项级数）
%       §4.2 幂级数（4.2.1 概念 / 4.2.2 收敛圆与收敛半径 / 4.2.3 性质 / 4.2.4 运算）
%       §4.3 泰勒级数（4.3.1 展开定理 / 4.3.2 初等函数的展开式）
%       §4.4 洛朗级数（4.4.1 概念与性质 / 4.4.2 展开定理 / 4.4.3 求展开式的方法）
%       §4.5 孤立奇点的分类（洛朗展开的直接应用）
% 由 1.main.tex 用 \input 引入

\chapter{级数}

级数是研究解析函数的第二套工具（第一套是积分）。本章的逻辑主线是：复数项级数 \(\to\) 幂级数（圆盘内） \(\to\) 泰勒级数（解析函数的幂级数表示） \(\to\) 洛朗级数（圆环内的广义幂级数） \(\to\) 孤立奇点分类。

\section{复变函数项级数}

\subsection{复数序列}
设 \(\{z_{n}\}\) 为复数序列，\(z_{n}=a_{n}+\mathrm{i}b_{n}\)。若存在复数 \(z_{0}=a_{0}+\mathrm{i}b_{0}\)，使 \(\lim\limits_{n\to\infty}|z_{n}-z_{0}|=0\)，则称 \(\{z_{n}\}\) \textbf{收敛}于 \(z_{0}\)。由于
\[
|a_{n}-a_{0}|\leqslant|z_{n}-z_{0}|,\qquad |b_{n}-b_{0}|\leqslant|z_{n}-z_{0}|,
\]
故
\keyformula{\lim_{n\to\infty}z_{n}=z_{0}\iff\lim_{n\to\infty}a_{n}=a_{0}\ \text{且}\ \lim_{n\to\infty}b_{n}=b_{0}}
即复数序列的收敛问题完全化归为两个实数列的收敛问题。由此得\textbf{柯西收敛准则}：\(\{z_{n}\}\) 收敛 \(\iff\) 对任意 \(\varepsilon>0\) 存在 \(N\)，当 \(n,m>N\) 时 \(|z_{n}-z_{m}|<\varepsilon\)。收敛序列必有界。

\subsection{复数项级数}
称 \(\sum\limits_{n=1}^{\infty}z_{n}=z_{1}+z_{2}+\cdots\) 为\textbf{复数项级数}，其部分和 \(S_{n}=\sum\limits_{k=1}^{n}z_{k}\)。若 \(\{S_{n}\}\) 收敛于 \(S\)，称级数\textbf{收敛}，\(S\) 为级数的和；否则称\textbf{发散}。

\textbf{定理}\quad \(\sum z_{n}\) 收敛 \(\iff\) \(\sum a_{n}\) 与 \(\sum b_{n}\) 都收敛；且此时 \(\sum z_{n}=\sum a_{n}+\mathrm{i}\sum b_{n}\)。

\textbf{基本性质}
\begin{itemize}[leftmargin=2em]
  \item 收敛的\textbf{必要条件}：若 \(\sum z_{n}\) 收敛，则 \(\lim\limits_{n\to\infty}z_{n}=0\)（逆不成立）；
  \item 线性性：\(\sum(\alpha z_{n}+\beta w_{n})=\alpha\sum z_{n}+\beta\sum w_{n}\)；
  \item 改变或去掉有限项不改变收敛性。
\end{itemize}

若 \(\sum|z_{n}|\) 收敛，称 \(\sum z_{n}\) \textbf{绝对收敛}。因为
\[
|a_{n}|\leqslant|z_{n}|,\qquad |b_{n}|\leqslant|z_{n}|,\qquad |z_{n}|\leqslant|a_{n}|+|b_{n}|,
\]
故
\keyformula{\sum z_{n}\ \text{绝对收敛}\iff\sum|z_{n}|\ \text{收敛}
\iff\sum a_{n},\ \sum b_{n}\ \text{都绝对收敛}}
并且\textbf{绝对收敛 \(\Rightarrow\) 收敛}（反之不成立）。判断绝对收敛时可直接套用正项级数的比值法、根值法、比较法。

\subsection{复变函数项级数}
设 \(f_{n}(z)\ (n=1,2,\cdots)\) 是点集 \(E\) 上的复函数。若对每个 \(z\in E\)，数项级数 \(\sum f_{n}(z)\) 都收敛，则称 \(\sum f_{n}(z)\) 在 \(E\) 上\textbf{逐点收敛}，其和函数记 \(S(z)=\sum\limits_{n=1}^{\infty}f_{n}(z)\)，部分和 \(S_{n}(z)=\sum\limits_{k=1}^{n}f_{k}(z)\)。

若对任意 \(\varepsilon>0\) 存在\textbf{与 \(z\) 无关}的 \(N\)，当 \(n>N\) 时对一切 \(z\in E\) 都有 \(|S_{n}(z)-S(z)|<\varepsilon\)，则称 \(\sum f_{n}(z)\) 在 \(E\) 上\textbf{一致收敛}。若它在区域 \(D\) 的任意有界闭子集上都一致收敛，则称在 \(D\) 上\textbf{内闭一致收敛}。

\textbf{优级数判别法（Weierstrass）}\quad 若存在收敛的正项级数 \(\sum M_{n}\)，使对一切 \(z\in E\) 有 \(|f_{n}(z)|\leqslant M_{n}\)，则 \(\sum f_{n}(z)\) 在 \(E\) 上一致收敛（且绝对收敛）。

\textbf{内闭一致收敛的性质}\quad 设 \(f_{n}(z)\) 都在区域 \(D\) 内解析，且 \(\sum f_{n}(z)\) 在 \(D\) 上内闭一致收敛于 \(S(z)\)，则
\begin{itemize}[leftmargin=2em]
  \item \textbf{（解析性）} \(S(z)\) 在 \(D\) 内解析；
  \item \textbf{（逐项求导）} \(S'(z)=\sum\limits_{n=1}^{\infty}f_{n}'(z)\)，且导数级数在 \(D\) 上内闭一致收敛；
  \item \textbf{（逐项积分）} 对 \(D\) 内任一条逐段光滑曲线 \(C\)，\(\displaystyle\int_{C}S(z)\,\mathrm{d}z=\sum_{n=1}^{\infty}\int_{C}f_{n}(z)\,\mathrm{d}z\)。
\end{itemize}

\corepoint{“逐点收敛”远不够用：\(\sum z^{n}\) 在 \(|z|<1\) 内逐点收敛于 \(\frac{1}{1-z}\)，但在整个开圆盘上\textbf{不一致}收敛（越靠近边界收敛越慢，\(S_{n}\) 在 \(|z|<1\) 上无界）。正确说法是“在 \(|z|\leqslant r\ (r<1)\) 上一致收敛”，即\textbf{内闭一致收敛}——它恰好是在开集上做逐项求导、逐项积分所需的强度。凡写“幂级数在其收敛圆内可逐项求导”，依据都是内闭一致收敛。}

\section{幂级数}

\subsection{幂级数的概念}
形如
\[
\sum_{n=0}^{\infty}a_{n}(z-z_{0})^{n}=a_{0}+a_{1}(z-z_{0})+a_{2}(z-z_{0})^{2}+\cdots
\]
的函数项级数称为以 \(z_{0}\) 为中心的\textbf{幂级数}，\(a_{n}\) 为系数。作平移 \(w=z-z_{0}\) 即可化为中心在原点的幂级数 \(\sum a_{n}w^{n}\)，故先研究后者。

\textbf{阿贝尔（Abel）定理}\quad 若 \(\sum a_{n}w^{n}\) 在 \(w_{0}\neq 0\) 处收敛，则它在 \(|w|<|w_{0}|\) 内\textbf{绝对收敛}；若在 \(w_{1}\) 处发散，则它在 \(|w|>|w_{1}|\) 内\textbf{发散}。

证明思路：由 \(\sum a_{n}w_{0}^{n}\) 收敛得 \(a_{n}w_{0}^{n}\) 有界，设 \(|a_{n}w_{0}^{n}|\leqslant M\)；对 \(|w|<|w_{0}|\) 有 \(|a_{n}w^{n}|\leqslant M\left|\frac{w}{w_{0}}\right|^{n}\)，与收敛的几何级数比较即得。

\subsection{幂级数的收敛圆与收敛半径}
由 Abel 定理，幂级数的收敛点集必是一个以 \(z_{0}\) 为心的圆盘内部（可能添上边界上的部分点）。定义
\[
R=\sup\left\{|z-z_{0}|:\ \sum_{n=0}^{\infty}a_{n}(z-z_{0})^{n}\ \text{收敛}\right\}
\]
为幂级数的\textbf{收敛半径}，圆盘 \(|z-z_{0}|<R\) 为\textbf{收敛圆}（\(R\) 可取 \(0\) 或 \(+\infty\)）。

\textbf{定理}\quad 设收敛半径为 \(R\)，则
\begin{itemize}[leftmargin=2em]
  \item 在 \(|z-z_{0}|<R\) 内幂级数\textbf{绝对收敛}且\textbf{内闭一致收敛}；
  \item 在 \(|z-z_{0}|>R\) 内幂级数\textbf{发散}；
  \item 在边界 \(|z-z_{0}|=R\) 上收敛性\textbf{不确定}，须具体讨论。
\end{itemize}

\textbf{收敛半径的计算}
\keyformula{R=\frac{1}{\varlimsup\limits_{n\to\infty}\sqrt[n]{|a_{n}|}}}
（约定 \(\frac{1}{0}=+\infty\)，\(\frac{1}{+\infty}=0\)。）若极限存在，也可用\textbf{比值法}
\[
R=\lim_{n\to\infty}\left|\frac{a_{n}}{a_{n+1}}\right|.
\]
实际计算时还可利用：幂级数在 \(z=z_{1}\) 处收敛 \(\Rightarrow R\geqslant|z_{1}-z_{0}|\)；在 \(z_{2}\) 处发散 \(\Rightarrow R\leqslant|z_{2}-z_{0}|\)；而 \(R\) 又等于 \(f(z)\) 在 \(z_{0}\) 的泰勒展开成立的最大圆盘半径，即 \(z_{0}\) 到 \(f\) 最近奇点的距离（这一条最快，见 §4.3）。

\corepoint{三个高频陷阱：(1) \textbf{缺项级数}（如 \(\sum 2^{n}z^{2n}\)、\(\sum a_{n}z^{kn}\)）不能直接用比值法 \(R=\lim|a_{n}/a_{n+1}|\)（此时 \(\lim\) 一般不存在），应把 \(z^{k}\) 看成整体或直接用根值公式 \(R=1/\varlimsup\sqrt[n]{|a_{n}|}\)；(2) 收敛半径只由系数决定，与 \(z\) 无关，圆心是展开中心 \(z_{0}\) 而非原点；(3) 收敛半径相同，边界上的收敛性可以完全不同——\(\sum z^{n}\) 在 \(|z|=1\) 上处处发散，\(\sum z^{n}/n^{2}\) 在 \(|z|=1\) 上处处绝对收敛，\(\sum z^{n}/n\) 在 \(|z|=1,z\neq 1\) 处收敛而在 \(z=1\) 处发散。所以“\(|z|=R\) 上收敛吗”必须单独讨论。}

\subsection{幂级数的性质}
\textbf{定理}\quad 幂级数 \(\sum a_{n}(z-z_{0})^{n}\) 的和函数 \(S(z)\) 在其收敛圆 \(|z-z_{0}|<R\) 内解析，且
\begin{itemize}[leftmargin=2em]
  \item 可\textbf{逐项求导}：\(S'(z)=\sum\limits_{n=1}^{\infty}na_{n}(z-z_{0})^{n-1}\)，收敛半径仍为 \(R\)；
  \item 可\textbf{逐项积分}：\(\displaystyle\int_{z_{0}}^{z}S(\xi)\,\mathrm{d}\xi=\sum_{n=0}^{\infty}\frac{a_{n}}{n+1}(z-z_{0})^{n+1}\)，收敛半径仍为 \(R\)；
  \item 反复求导可知 \(S(z)\) 在收敛圆内\textbf{无穷次可导}，且 \(a_{n}=\dfrac{S^{(n)}(z_{0})}{n!}\)（系数被和函数唯一确定）。
\end{itemize}

\subsection{幂级数的运算}
设 \(\sum a_{n}(z-z_{0})^{n}\)、\(\sum b_{n}(z-z_{0})^{n}\) 的收敛半径分别为 \(R_{1},R_{2}\)，\(R=\min\{R_{1},R_{2}\}\)。在 \(|z-z_{0}|<R\) 内可作
\begin{itemize}[leftmargin=2em]
  \item \textbf{加减法}：\(\sum(a_{n}\pm b_{n})(z-z_{0})^{n}\)，收敛半径至少为 \(R\)；
  \item \textbf{乘法}（柯西乘积）：\(\left(\sum a_{n}w^{n}\right)\left(\sum b_{n}w^{n}\right)=\sum c_{n}w^{n}\)，其中 \(c_{n}=\sum\limits_{k=0}^{n}a_{k}b_{n-k}\)，收敛半径至少为 \(R\)；
  \item \textbf{除法}：当 \(b_{0}\neq 0\) 时商仍为幂级数，系数由 \(\sum a_{n}w^{n}=\left(\sum b_{n}w^{n}\right)\left(\sum d_{n}w^{n}\right)\) 逐项比较（\textbf{待定系数法}）确定；
  \item \textbf{复合}：若 \(g(z)=\sum b_{n}w^{n}\) 且 \(w=f(z)\) 在 \(|z-z_{0}|<\rho\) 内解析、\(|f(z)|<R_{2}\)，则 \(g[f(z)]\) 可展成幂级数。
\end{itemize}

\section{泰勒级数}

\subsection{泰勒展开定理}
\textbf{定理（泰勒 Taylor 展开）}\quad 设 \(f(z)\) 在圆盘 \(B(z_{0},R):|z-z_{0}|<R\) 内解析，则 \(f\) 在该圆盘内可唯一地展成幂级数
\keyformula{f(z)=\sum_{n=0}^{\infty}a_{n}(z-z_{0})^{n},\qquad
a_{n}=\frac{f^{(n)}(z_{0})}{n!}=\frac{1}{2\pi\mathrm{i}}\oint_{C}\frac{f(\zeta)}{(\zeta-z_{0})^{n+1}}\,\mathrm{d}\zeta}
其中 \(C\) 为 \(|z-z_{0}|=\rho\ (0<\rho<R)\)，且该级数在 \(B(z_{0},R)\) 内绝对收敛、内闭一致收敛。

证明思路：对 \(|z-z_{0}|<r<R\)，由柯西积分公式
\(f(z)=\frac{1}{2\pi\mathrm{i}}\oint_{|\zeta-z_{0}|=r}\frac{f(\zeta)}{\zeta-z_{0}}\cdot\frac{1}{1-\frac{z-z_{0}}{\zeta-z_{0}}}\mathrm{d}\zeta\)，
把 \(\frac{1}{1-w}\) 展成几何级数（\(|w|<1\)）并逐项积分，再用高阶导数公式即得系数表达式。\(z_{0}=0\) 时称为\textbf{麦克劳林（Maclaurin）展开}。

\textbf{几点说明}
\begin{itemize}[leftmargin=2em]
  \item \textbf{收敛范围}：定理只保证在 \(B(z_{0},R)\) 内成立；实际展开的收敛半径等于 \(z_{0}\) 到 \(f\) 的\textbf{最近奇点}的距离（若 \(f\) 是整函数则 \(R=+\infty\)）。
  \item \textbf{唯一性}：解析函数在一点处的幂级数展开式唯一（系数只能是 \(f^{(n)}(z_{0})/n!\)）。这条性质是“间接展开法”的合法性依据。
  \item \textbf{解析的等价刻画}：\(f\) 在 \(z_{0}\) 解析 \(\iff f\) 在 \(z_{0}\) 的某邻域内可展成幂级数。
\end{itemize}

\textbf{两个重要推论}
\begin{itemize}[leftmargin=2em]
  \item \textbf{零点孤立性定理}：设 \(f\) 在区域 \(D\) 内解析且不恒为零，则 \(f\) 的每个零点是\textbf{孤立}的——即若 \(f(z_{0})=0\)，则存在 \(z_{0}\) 的去心邻域使 \(f\) 在其中不为零。
  \item \textbf{唯一性定理}：设 \(f,g\) 在区域 \(D\) 内解析，若存在 \(D\) 内一列互不相同的点 \(\{z_{n}\}\)（有聚点 \(z^{*}\in D\)）使 \(f(z_{n})=g(z_{n})\)，则 \(f\equiv g\) 于 \(D\)。
\end{itemize}

\subsection{几个初等函数的幂级数展开式}
以下展开式必须记熟（\(z_{0}=0\)）：
\keyformula{\mathrm{e}^{z}=\sum_{n=0}^{\infty}\frac{z^{n}}{n!},
\qquad \sin z=\sum_{n=0}^{\infty}\frac{(-1)^{n}z^{2n+1}}{(2n+1)!},
\qquad \cos z=\sum_{n=0}^{\infty}\frac{(-1)^{n}z^{2n}}{(2n)!}\qquad (|z|<+\infty)}
\keyformula{\frac{1}{1-z}=\sum_{n=0}^{\infty}z^{n}\ (|z|<1),\qquad
\ln(1+z)=\sum_{n=1}^{\infty}\frac{(-1)^{n-1}}{n}z^{n}\ (|z|<1)}
\keyformula{(1+z)^{\alpha}=1+\alpha z+\frac{\alpha(\alpha-1)}{2!}z^{2}+\cdots
+\frac{\alpha(\alpha-1)\cdots(\alpha-n+1)}{n!}z^{n}+\cdots\qquad (|z|<1)}

\textbf{间接展开法（最常用）}\quad 由展开式的唯一性，可用下列手段从已知展开式得到新的展开式，而不必逐阶求导：
\begin{enumerate}[leftmargin=2em]
  \item \textbf{代换}：如 \(\mathrm{e}^{z^{2}}=\sum\frac{z^{2n}}{n!}\)、\(\dfrac{1}{1+z^{2}}=\sum(-1)^{n}z^{2n}\)；
  \item \textbf{分拆为部分分式}：如有理函数先拆成 \(\frac{A}{z-a}+\frac{B}{z-b}\) 再逐个展开；
  \item \textbf{逐项求导或逐项积分}：如由 \(\frac{1}{1-z}\) 求导得 \(\frac{1}{(1-z)^{2}}\) 的展开，由 \(\frac{1}{1+z}\) 积分得 \(\ln(1+z)\)；
  \item \textbf{乘除法与待定系数}：如求 \(\tan z\)、\(\sec z\) 的展开。
\end{enumerate}

\corepoint{展开式后面\textbf{必须写明成立范围}（\(|z|<1\) 还是 \(|z|<+\infty\)），这是泰勒展开与洛朗展开的“坐标”。另外，间接法只允许在\textbf{收敛圆内}做代数变形：由 \(\ln(1+z)=\sum(-1)^{n-1}z^{n}/n\ (|z|<1)\) 不能直接推出 \(\ln(1+z^{2})\) 在 \(|z|<+\infty\) 成立，只能得 \(|z|<1\)；正确范围由“最近奇点”判断——\(1+z^{2}=0\) 的根 \(\pm\mathrm{i}\) 在 \(|z|=1\) 上，故半径确为 \(1\)。}

\section{洛朗级数}

\subsection{洛朗级数的概念及性质}
形如
\[
\sum_{n=-\infty}^{+\infty}a_{n}(z-z_{0})^{n}
=\underbrace{\sum_{n=0}^{\infty}a_{n}(z-z_{0})^{n}}_{\text{正部}}
+\underbrace{\sum_{n=1}^{\infty}a_{-n}(z-z_{0})^{-n}}_{\text{负部}}
\]
的函数项级数称为以 \(z_{0}\) 为中心的\textbf{广义幂级数}（洛朗级数）。它收敛 \(\iff\) 正部与负部同时收敛。

设正部收敛半径为 \(R_{1}\)，负部在 \(|z-z_{0}|>R_{2}\) 收敛，则当 \(R_{2}<R_{1}\) 时，级数在\textbf{圆环}
\[
R_{2}<|z-z_{0}|<R_{1}
\]
内绝对收敛且内闭一致收敛，此圆环称为\textbf{收敛圆环}；在圆环外发散，在边界上不确定。特别当 \(R_{2}=0\) 时圆环退化为去心圆盘 \(0<|z-z_{0}|<R_{1}\)。

\subsection{洛朗展开定理}
\textbf{定理（洛朗 Laurent 展开）}\quad 设 \(f(z)\) 在圆环 \(D:r<|z-z_{0}|<R\) 内解析，则 \(f\) 在 \(D\) 内可唯一地展成洛朗级数
\keyformula{f(z)=\sum_{n=-\infty}^{+\infty}a_{n}(z-z_{0})^{n},\qquad
a_{n}=\frac{1}{2\pi\mathrm{i}}\oint_{C}\frac{f(\zeta)}{(\zeta-z_{0})^{n+1}}\,\mathrm{d}\zeta
\quad(n=0,\pm 1,\pm 2,\cdots)}
其中 \(C\) 为圆环内任一以 \(z_{0}\) 为心的圆周 \(|\zeta-z_{0}|=\rho\ (r<\rho<R)\)，且级数在 \(D\) 内绝对收敛、内闭一致收敛。

与泰勒定理的两个差别：
\begin{itemize}[leftmargin=2em]
  \item 系数一般\textbf{不能}写成 \(\frac{f^{(n)}(z_{0})}{n!}\)（\(f\) 在 \(z_{0}\) 处可能没有定义，谈不上求导）；
  \item 展开式依赖\textbf{圆环}：同一个 \(f\) 在 \(z_{0}\) 的不同圆环上有不同的洛朗展开。
\end{itemize}

\textbf{唯一性}\quad \(f\) 在给定圆环内的洛朗展开式是唯一的；故用任何合法手段（代换、分拆、乘除）得到的展开式就是它。

\subsection{求洛朗展开式的方法}
\textbf{步骤}
\begin{enumerate}[leftmargin=2em]
  \item 找出 \(f\) 的全部奇点，确定展开中心 \(z_{0}\)；
  \item 由奇点到 \(z_{0}\) 的距离把平面分成若干圆环，题目指定哪个就在哪个圆环内做；
  \item 把 \(f\) 拆成若干“简单分式”之和（有理函数用部分分式）；
  \item 对每个分式，按“\(|z-z_{0}|\) 与奇点距离的大小”选择展开方向：
  若 \(|w|<1\) 用 \(\dfrac{1}{1-w}=\sum w^{n}\)；若 \(|w|>1\) 用 \(\dfrac{1}{1-w}=-\dfrac{1}{w}\cdot\dfrac{1}{1-1/w}=-\sum\limits_{n=1}^{\infty}w^{-n}\)；
  \item 合并同类项，并\textbf{写明成立区域}。
\end{enumerate}

\textbf{判别准则}\quad 对每个因子 \(\dfrac{1}{z-a}\)：若 \(|z-z_{0}|<|a-z_{0}|\)（奇点 \(a\) 落在圆环\textbf{之外}），展成 \(z-z_{0}\) 的\textbf{正幂}；若 \(|z-z_{0}|>|a-z_{0}|\)（\(a\) 落在圆环\textbf{之内}），先提出 \((z-z_{0})\) 作公因子再展成\textbf{负幂}。同一函数在不同圆环上的展开式互不相同，但在各自圆环内都表示同一个 \(f\)。

\corepoint{洛朗展开最容易错的三处：(1) 忘了区分圆环，把 \(|z|<1\) 的展开式拿去用在 \(1<|z|<2\) 上；(2) 该出负幂时仍机械地写 \(\frac{1}{1-z}=\sum z^{n}\)（此式只在 \(|z|<1\) 成立）；(3) 展开中心不在原点时忘记把 \(z-z_{0}\) 当作整体 \(w\)。判别口诀：看 \(|z-z_{0}|\) 与“该分式的奇点到 \(z_{0}\) 的距离”谁大——\(|z-z_{0}|\) 大就用负幂，小就用正幂。}

\section{孤立奇点的分类}

本节是洛朗展开的直接应用（教材排在 §5.1，留数的计算见第 5 章）。

\textbf{定义}\quad 若 \(f(z)\) 在 \(z_{0}\) 的某个\textbf{去心邻域} \(0<|z-z_{0}|<\delta\) 内解析，而在 \(z_{0}\) 处不解析（或未定义），则称 \(z_{0}\) 为 \(f\) 的\textbf{孤立奇点}。

设 \(f\) 在 \(0<|z-z_{0}|<R\) 内的洛朗展开为 \(f(z)=\sum\limits_{n=-\infty}^{+\infty}a_{n}(z-z_{0})^{n}\)，其负幂项部分 \(\sum\limits_{n=1}^{\infty}a_{-n}(z-z_{0})^{-n}\) 称为\textbf{主要部分}。

\begin{center}
\begin{tabular}{llll}
\toprule
类型 & 主要部分 & 极限特征 & 等价刻画 \\
\midrule
可去奇点 & 无（\(a_{-n}=0\)） & \(\lim f(z)\) 存在且有限 & 去心邻域内有界 \\
\(m\) 阶极点 & 有限项，最低次 \(-m\) & \(\lim f(z)=\infty\) & \(f=\dfrac{g(z)}{(z-z_{0})^{m}}\)，\(g(z_{0})\neq0\) \\
本性奇点 & 无穷多项 & \(\lim f(z)\) 不存在 & 取值稠密（Picard） \\
\bottomrule
\end{tabular}
\end{center}

\textbf{零点与极点的关系}\quad 若 \(f(z)=(z-z_{0})^{m}\varphi(z)\)，\(\varphi\) 在 \(z_{0}\) 解析且 \(\varphi(z_{0})\neq0\)，则称 \(z_{0}\) 为 \(f\) 的 \textbf{\(m\) 阶零点}。此时
\[
z_{0}\ \text{是}\ f\ \text{的}\ m\ \text{阶零点}\iff z_{0}\ \text{是}\ \frac{1}{f}\ \text{的}\ m\ \text{阶极点}.
\]
对有理函数 \(f=\dfrac{P(z)}{Q(z)}\)（\(P,Q\) 在 \(z_{0}\) 解析，\(z_{0}\) 分别是 \(P\) 的 \(m\) 阶零点、\(Q\) 的 \(n\) 阶零点）：\(m>n\) 时 \(z_{0}\) 是 \(f\) 的 \(m-n\) 阶零点；\(m=n\) 时是可去奇点；\(m<n\) 时是 \(n-m\) 阶极点。

\corepoint{判定顺序：先看\textbf{极限}，再用\textbf{洛朗展开}定阶数。“极限不存在”不等于“本性奇点”——极限既不有限也不为 \(\infty\) 才是本性奇点（如 \(\mathrm{e}^{1/z}\) 在 \(0\) 处）；而 \(\lim f=\infty\) 只说明是极点，阶数还要靠展开式或零点阶数确定。又：孤立奇点要求“去心邻域内解析”，故 \(\operatorname{Ln}z\) 的支点 \(0\)、\(\frac{1}{\sin(1/z)}\) 在 \(z=0\) 处的奇点聚点都\textbf{不是}孤立奇点。}

\chapter{留数}

本章按讲义"定义—奇点分类—计算方法—无穷远点—实积分应用"五节展开，教材 \S5.4 的辐角原理与儒歇定理作为选学内容附在章末。

\section{留数的定义和留数基本定理}

\subsection{孤立奇点与留数}

\begin{itemize}
    \item 复平面上的点 \(z_0\) 是 \(f(z)\) 的\textbf{孤立奇点}，当且仅当存在 \(\delta>0\)，使 \(f(z)\) 在去心邻域 \(B^\circ(z_0,\delta)=\{\,0<|z-z_0|<\delta\,\}\) 内解析。
    \item "孤立"两字是关键：\(z=0\) 是 \(f(z)=\dfrac{1}{\sin(1/z)}\) 的奇点，但它是一串奇点 \(z_k=\dfrac{1}{k\pi}\) 的聚点，任何去心邻域里都有别的奇点，故 \textbf{不是}孤立奇点。
\end{itemize}

设 \(z_0\) 是 \(f(z)\) 的孤立奇点，则 \(f(z)\) 在 \(B^\circ(z_0,\delta)\) 内可展成洛朗级数
\[
f(z)=\sum_{n=-\infty}^{+\infty}a_n(z-z_0)^n .
\]
其中 \(a_{-1}\)（即 \(\dfrac{1}{z-z_0}\) 项的系数）称为 \(f(z)\) 在 \(z_0\) 处的\textbf{留数}（residue），记作 \(\operatorname{Res}[f(z),z_0]\)。留数只由 \(a_{-1}\) 决定，与展开式的其余部分无关。

\keyformula{\operatorname{Res}[f(z),z_0]=a_{-1}=\frac{1}{2\pi i}\oint_{|z-z_0|=\delta}f(z)\,dz\quad(0<\delta<R)}

这个等式来自洛朗级数的逐项积分：对任意整数 \(n\)，\(\displaystyle\oint_{|z-z_0|=\delta}(z-z_0)^n dz\) 当 \(n=-1\) 时为 \(2\pi i\)、当 \(n\neq-1\) 时为 \(0\)，故只有 \(a_{-1}\) 项在积分中存活。

\subsection{留数基本定理}

\begin{itemize}
    \item \textbf{引理}：设 \(f(z)\) 在去心开圆盘 \(B^\circ(z_0,R)\) 上解析，则当 \(0<\delta<R\) 时 \(\displaystyle\oint_{|z-z_0|=\delta}f(z)dz=2\pi i\operatorname{Res}[f(z),z_0]\)。
    \item \textbf{留数基本定理}：设 \(C\) 是一条简单光滑闭曲线（取正向），\(f(z)\) 在 \(C\) 所围闭区域上除内部有限个孤立奇点 \(a_1,a_2,\dots,a_n\) 外解析，则
\end{itemize}

\keyformula{\oint_C f(z)\,dz=2\pi i\sum_{k=1}^{n}\operatorname{Res}[f(z),a_k]}

证明思路：在 \(C\) 内作两两不交的小圆盘 \(B(a_k,\varepsilon_k)\)，在挖去这些小圆盘的闭区域上用柯西积分定理得 \(\oint_C=\sum_k\oint_{|z-a_k|=\varepsilon_k}\)，再用引理。

对\textbf{有界光滑连通闭区域} \(D\)（可以带洞）有更一般的写法：设 \(f(z)\) 在 \(D\) 上除内部点 \(a_1,\dots,a_n\) 外解析，则
\[
\int_{\partial D}f(z)\,dz=2\pi i\sum_{k=1}^{n}\operatorname{Res}[f(z),a_k],
\]
其中 \(\partial D\) 的外边界取逆时针、各内边界取顺时针（正向即"区域始终在左侧"）。

\corepoint{留数定理只统计 \textbf{闭曲线内部} 的奇点；边界上的奇点不能直接计入，也不能让积分路径穿过奇点。带洞区域的内边界要取 \textbf{顺时针}（反向），否则符号整体反号。}

\section{孤立奇点的分类和性质}

设 \(z_0\) 是 \(f(z)\) 的孤立奇点，洛朗展开中 \(n<0\) 的部分称为\textbf{负部}（主要部分），按负部形态分为三类：

\begin{center}
\begin{tabular}{@{}lll@{}}
\toprule
类型 & 洛朗展开负部 & 极限刻画 \\
\midrule
可去奇点 & 无（恒为零） & \(\lim\limits_{z\to z_0}f(z)\) 存在且有限 \\
极点 & 有限项（最高为 \((z-z_0)^{-m}\)） & \(\lim\limits_{z\to z_0}f(z)=\infty\) \\
本性奇点 & 无限项 & \(\lim\limits_{z\to z_0}f(z)\)（广义）不存在 \\
\bottomrule
\end{tabular}
\end{center}

\subsection{可去奇点}

\textbf{命题}：\(z_0\) 是 \(f\) 的可去奇点 \(\iff \lim_{z\to z_0}f(z)\) 存在有限。

由幂级数理论，负部为零时右端在 \(B(z_0,r)\) 内解析，故极限存在；反之极限有限则 \(f\) 在去心邻域有界，由积分估值不等式可证各 \(a_{-n}=0\)。补充定义 \(f(z_0)=\lim_{z\to z_0}f(z)\) 后 \(z_0\) 即为解析点，这正是"可去"的含义。

\textbf{开映射定理}：设 \(f(z)\) 在 \(z_0\) 处解析且在 \(z_0\) 的某邻域内不是常值函数，则对任一充分小的 \(\varepsilon>0\)，存在 \(\rho>0\) 使 \(B(f(z_0),\rho)\subset f\big(B(z_0,\varepsilon)\big)\)，即非常值解析函数把开集映成开集。

\subsection{零点、极点及其阶数}

\begin{itemize}
    \item \(z_0\) 是 \(f(z)\) 的 \(m\) 阶\textbf{零点}，当且仅当
    \[
    f(z_0)=f'(z_0)=\cdots=f^{(m-1)}(z_0)=0,\qquad f^{(m)}(z_0)\neq 0 .
    \]
    等价地，\(f(z)=(z-z_0)^m g(z)\)，其中 \(g(z)\) 在 \(z_0\) 处解析且 \(g(z_0)\neq 0\)；此时泰勒展开从 \(a_m(z-z_0)^m\) 项开始。
    \item \(z_0\) 是 \(f(z)\) 的 \(m\) 阶\textbf{极点}，当且仅当洛朗展开的负部恰为
    \[
    a_{-m}(z-z_0)^{-m}+\cdots+a_{-1}(z-z_0)^{-1},\qquad a_{-m}\neq 0,
    \]
    即主部首项为 \((z-z_0)^{-m}\)。
\end{itemize}

\textbf{命题}：有限阶零点都是孤立的。设 \(z_0\) 是 \(f\) 的 \(m\) 阶零点，则在 \(z_0\) 的某去心邻域内 \(f(z)\neq 0\)。

\textbf{零点与极点的对偶定理}：\(z_0\) 是 \(f(z)\) 的 \(m\) 阶零点 \(\iff\) \(z_0\) 是 \(\dfrac{1}{f(z)}\) 的 \(m\) 阶极点。

\corepoint{这个对偶是把"判断极点阶数"转化为"判断零点阶数"的桥梁：遇到 \(f(z)\) 结构复杂（如 \(1/\sin^2 z\)），先看分母 \(\sin^2 z\) 的零点阶数更省力。}

\textbf{极点判定定理}：\(z_0\) 是 \(f(z)\) 的极点 \(\iff\lim_{z\to z_0}f(z)=\infty\)。

\textbf{分式函数的阶数推论}：设 \(f(z)=\dfrac{P(z)}{Q(z)}\)，\(P,Q\) 在区域 \(D\) 内解析，\(z_0\) 分别是 \(P,Q\) 的 \(m\) 阶和 \(n\) 阶零点，则

\begin{center}
\begin{tabular}{@{}lll@{}}
\toprule
条件 & \(z_0\) 对 \(f\) 的类型 & 说明 \\
\midrule
\(m>n\) & \((m-n)\) 阶零点 & 分子"更强" \\
\(m=n\) & 可去奇点 & 极限 \(\lim f=\dfrac{P^{(m)}(z_0)}{Q^{(n)}(z_0)}\neq 0\) 有限 \\
\(m<n\) & \((n-m)\) 阶极点 & 分母"更强" \\
\bottomrule
\end{tabular}
\end{center}

\textbf{极点的局部取值}：设 \(z_0\) 是 \(f(z)\) 的极点，则在 \(z_0\) 的某去心邻域内 \(\{|f(z)|>\rho\}\) 包含一个 \(\{|z-z_0|<\delta\}\)，这就是 \(\lim f=\infty\) 的确切含义。

\subsection{本性奇点}

\textbf{定理}：\(z_0\) 是 \(f(z)\) 的本性奇点 \(\iff\) \(\lim_{z\to z_0}f(z)\) 在广义意义下不存在。

典型例子是 \(f(z)=e^{1/z}\) 在 \(z=0\)：沿正实轴 \(z\to0^+\) 时 \(e^{1/z}\to+\infty\)，沿负实轴 \(z\to0^-\) 时 \(e^{1/z}\to0\)，沿虚轴 \(z=iy\to0\) 时 \(e^{-i/y}\) 不趋于定值。方向极限互不相同，故极限不存在。

\textbf{Picard 大定理}：设 \(z_0\) 是 \(f(z)\) 的本性奇点，\(f\) 在 \(B^\circ(z_0,\delta)\) 上解析，则 \(\mathbb{C}\setminus f\big(B^\circ(z_0,\delta)\big)\) 中至多含一个复数。即本性奇点附近函数取遍几乎整个复平面（\(e^{1/z}\) 取不到 \(0\)，\(\sin(1/z)\) 取遍 \(\mathbb{C}\)）。

\textbf{复合函数判定推论}：设 \(f(z)=\dfrac{g(z)}{h(z)}\)，\(z_0\) 是 \(f\) 的本性奇点（即 \(h\) 的行为"失控"），则
\begin{enumerate}[label=(\arabic*)]
    \item 若 \(z_0\) 是 \(g\) 的本性奇点，则 \(z_0\) 也是 \(f\) 的本性奇点；
    \item 若 \(z_0\) 是 \(g\) 的解析点或可去奇点且 \(\lim\limits_{z\to z_0}g(z)\neq 0\)，则 \(z_0\) 是 \(f\) 的本性奇点；
    \item 若 \(z_0\) 是 \(g\) 的 \(m\) 阶零点或 \(m\) 阶极点，则 \(z_0\) 是 \(f\) 的本性奇点。
\end{enumerate}
第 (3) 条常用于判定 \(e^{1/z}\tan z\)、\(\dfrac{1}{\tan z}\) 这类"零点与本性奇点相乘"的情形。
\section{留数的计算}

\subsection{基本规则}

\begin{enumerate}[label=(\arabic*)]
    \item \textbf{线性性}：设 \(z_0\) 是 \(f,g\) 的孤立奇点，则
    \[
    \operatorname{Res}[f\pm g,z_0]=\operatorname{Res}[f,z_0]\pm\operatorname{Res}[g,z_0],\qquad
    \operatorname{Res}[af,z_0]=a\operatorname{Res}[f,z_0],\quad a\in\mathbb{C}.
    \]
    \item \textbf{可去奇点}：若 \(z_0\) 是 \(f\) 的可去奇点，则 \(\operatorname{Res}[f(z),z_0]=0\)。
    \item \textbf{\(m\) 阶极点}：
\end{enumerate}

\keyformula{\operatorname{Res}[f(z),z_0]=\frac{1}{(m-1)!}\lim_{z\to z_0}\frac{d^{\,m-1}}{dz^{\,m-1}}\Big[(z-z_0)^m f(z)\Big]}

证明：\((z-z_0)^mf(z)=a_{-m}+a_{-m+1}(z-z_0)+\cdots+a_{-1}(z-z_0)^{m-1}+\cdots\)，求 \(m-1\) 阶导并令 \(z\to z_0\) 即得 \((m-1)!\,a_{-1}\)。

当 \(m=1\) 时公式退化为最常用的
\[
\operatorname{Res}[f(z),z_0]=\lim_{z\to z_0}(z-z_0)f(z).
\]

\begin{enumerate}[label=(\arabic*),start=4]
    \item \textbf{一阶极点（分子分母形式）}：设 \(f(z)=\dfrac{g(z)}{h(z)}\)，\(g,h\) 在 \(z_0\) 处解析，\(z_0\) 是 \(h\) 的一阶零点且 \(g(z_0)\neq 0\)，则
\end{enumerate}

\keyformula{\operatorname{Res}\!\left[\frac{g(z)}{h(z)},z_0\right]=\frac{g(z_0)}{h'(z_0)}}

\subsection{计算步骤与选择}

\begin{center}
\begin{tabular}{@{}lll@{}}
\toprule
奇点类型 & 首选方法 & 备注 \\
\midrule
可去奇点 & 直接写 \(0\) & 不必展开 \\
一阶极点 & \(\lim (z-z_0)f(z)\) 或 \(g(z_0)/h'(z_0)\) & 后者要求 \(h'(z_0)\neq 0\) \\
\(m\) 阶极点（\(m\ge 2\)） & 求导公式 & 求导次数多时改用展开 \\
本性奇点 & 洛朗展开取 \(a_{-1}\) & 唯一可行的方法 \\
\bottomrule
\end{tabular}
\end{center}

\textbf{一般步骤}：先定位全部奇点并判断类型；可去奇点记 \(0\)，极点按阶数选公式，高阶级点与本性奇点回到洛朗展开；闭路内奇点太多时改算闭路外（含 \(\infty\)）的留数，见 \S\ref{sec:inf-residue}。

\corepoint{三类常见错误：(1) \(m\) 阶极点漏掉因子 \(\dfrac{1}{(m-1)!}\)；(2) 用 \(\dfrac{g(z_0)}{h'(z_0)}\) 时忘记验证 \(h'(z_0)\neq 0\)（如 \(h(z)=\sin^2 z\) 在 \(z=0\) 处 \(h'(0)=0\)，是二阶极点，不能套）；(3) 对 \(e^{1/z}\)、\(\sin(1/z)\) 这类本性奇点硬套极点公式——必须展成级数找 \(a_{-1}\)。}

\section{无穷远点处的留数}\label{sec:inf-residue}

\subsection{无穷远点作为孤立奇点}

\begin{itemize}
    \item 无穷远点 \(\infty\) 是 \(f(z)\) 的\textbf{孤立奇点}，当且仅当存在 \(R>0\)，使 \(f(z)\) 在 \(\{|z|>R\}\) 上解析。
    \item 用变换 \(\xi=\dfrac{1}{z}\) 把 \(\infty\) 搬到原点，令 \(g(\xi)=f\!\left(\dfrac{1}{\xi}\right)\)，则
    \[
    \begin{aligned}
    &\infty\ \text{是}\ f\ \text{的可去奇点（}m\text{ 阶极点、本性奇点）}\\
    \iff\ &0\ \text{是}\ g\ \text{的可去奇点（}m\text{ 阶极点、本性奇点）}.
    \end{aligned}
    \]
\end{itemize}

几个典型判定（把 \(\xi=1/z\) 代进去，看 \(g(\xi)\) 在 \(0\) 处的类型）：\(f(z)=az^n+\cdots+a_0\) 的 \(\infty\) 是 \(n\) 阶极点（\(g\) 在 \(0\) 有 \(n\) 阶极点）；\(1/z\) 的 \(\infty\) 是一阶零点（可去奇点）；\(\sin z\) 的 \(\infty\) 是本性奇点（\(g(\xi)=\sin(1/\xi)\)）；而 \(1/\sin(1/z)\) 的 \(\infty\) \emph{不是}孤立奇点，因为 \(g(\xi)=1/\sin\xi\) 的奇点 \(\xi=k\pi\) 以 \(0\) 为聚点。注意"\(\infty\) 是否孤立"与"\(\infty\) 属于哪一类"是两个独立问题。

\subsection{无穷远点留数的定义}

无穷远点的留数不能简单取成 \(g(\xi)=f(1/\xi)\) 在 \(0\) 处的留数（积分方向被反转了一次）。正确做法：把 \(f\) 在环域 \(\{|z|>R\}\) 上展开为 \(f(z)=\sum_{n=-\infty}^{+\infty}a_n z^n\)，再定义

\keyformula{\operatorname{Res}[f(z),\infty]:=-\operatorname{Res}\!\left[\frac{1}{\xi^2}f\!\left(\frac{1}{\xi}\right),0\right]=-a_{-1}}

这里 \(a_{-1}\) 是 \(f\) 在 \(\{|z|>R\}\) 内洛朗展开中 \(z^{-1}\) 的系数。等价地
\[
\operatorname{Res}[f(z),\infty]=-\frac{1}{2\pi i}\oint_{|z|=R}f(z)\,dz\quad(\text{顺时针取正向}),
\]
负号来自"绕无穷远点的正向是顺时针"。

\begin{itemize}
    \item \textbf{定理}：若 \(f(z)\) 在扩充复平面上只有有限多个奇点 \(\{a_1,a_2,\dots,a_n,\infty\}\)，则
\end{itemize}

\keyformula{\sum_{k=1}^{n}\operatorname{Res}[f(z),a_k]+\operatorname{Res}[f(z),\infty]=0}

这条定理把"算 \(n\) 个有限奇点"换成"算 1 个无穷远点"，在大圆弧围道积分中尤其省事：只要求出 \(\infty\) 处的留数，围道内所有奇点留数之和立即得到（差一个负号）。

\corepoint{三个易错点：(1) 定义里必须除以 \(\xi^2\)，只算 \(f(1/\xi)\) 在 \(0\) 的留数是错的方向；(2) \(a_{-1}\) 要取 \emph{环域} \(|z|>R\) 上的展开系数，不能拿 \(|z|<R\) 的展开混用；(3) 全平面留数和为零是对 \emph{扩充} 复平面说的，漏掉 \(\infty\) 就不成立。}

\section{留数在实积分中的应用}

课程要求掌握三类实积分，共同套路是：选一条闭围道，使待求实积分等于围道积分在极限情形下的主部，再用留数定理。

\subsection{形如 \(\int_0^{2\pi}R(\cos\theta,\sin\theta)\,d\theta\) 的积分}

设 \(R(u,v)\) 是二元有理函数，且在单位圆周上无极点。令 \(z=e^{i\theta}\)，则 \(\theta:0\to2\pi\) 对应 \(z\) 沿 \(|z|=1\) 正向绕一圈，且
\[
\cos\theta=\frac{z+z^{-1}}{2},\qquad
\sin\theta=\frac{z-z^{-1}}{2i},
\]
\[
d\theta=\frac{dz}{iz}.
\]

\keyformula{\int_0^{2\pi}R(\cos\theta,\sin\theta)\,d\theta
=2\pi i\sum_{|z_k|<1}\operatorname{Res}\!\left[\frac{1}{iz}R\!\left(\frac{z+z^{-1}}{2},\frac{z-z^{-1}}{2i}\right),z_k\right]}

要点：只取 \textbf{单位圆内} 的极点；遇到 \(\cos2\theta,\sin2\theta\) 用 \(z^2=e^{2i\theta}\) 或先降幂，避免阶数算错。

\subsection{形如 \(\int_{-\infty}^{+\infty}R(x)\,dx\) 的积分}

设 \(R(z)=\dfrac{P(z)}{Q(z)}\) 为有理函数，满足
\begin{enumerate}[label=(\arabic*)]
    \item \(\deg Q\ge \deg P+2\)；
    \item \(Q(z)\) 在实轴上无零点。
\end{enumerate}
取上半平面大围道（实轴段加上半大圆弧 \(C_R\)）。由 \(\deg\) 条件与积分估值不等式，大圆弧上的积分满足
\[
\lim_{R\to\infty}\int_{C_R}R(z)\,dz=0,
\]
于是

\keyformula{\int_{-\infty}^{+\infty}R(x)\,dx=2\pi i\sum_{\operatorname{Im}z_k>0}\operatorname{Res}[R(z),z_k]}

\subsection{形如 \(\int_{-\infty}^{+\infty}R(x)e^{iax}\,dx\)（\(a>0\)）的积分}

设 \(R(z)=\dfrac{P(z)}{Q(z)}\)，\(\deg Q\ge \deg P+1\)，\(R\) 在实轴上的极点都是一阶的。此时大圆弧上的积分由 \textbf{Jordan 引理} 控制（\(e^{iaz}\) 在上半平面指数衰减），得到

\keyformula{\int_{-\infty}^{+\infty}R(x)e^{iax}dx
=2\pi i\!\!\sum_{\operatorname{Im}z_k>0}\!\!\operatorname{Res}[R(z)e^{iaz},z_k]
+\pi i\sum_{x_j\in\mathbb{R}}\operatorname{Res}[R(z)e^{iaz},x_j]}

其中第二项的来历：实轴上的极点要用上半平面小半圆"绕过去"（顺时针），每个小半圆贡献 \(-i\pi\operatorname{Res}\)，移到等式右边即变号成 \(+\pi i\)。

由此取实部、虚部得
\[
\int_{-\infty}^{+\infty}R(x)\cos(ax)\,dx=\operatorname{Re}\{\cdots\},\qquad
\int_{-\infty}^{+\infty}R(x)\sin(ax)\,dx=\operatorname{Im}\{\cdots\}.
\]

\textbf{Dirichlet 积分}（后续傅里叶变换要用）：对 \(a\neq0\),
\[
\int_{-\infty}^{+\infty}\frac{\sin ax}{x}\,dx=\pi\operatorname{sgn}a,
\qquad
\int_0^{+\infty}\frac{\sin ax}{x}\,dx=\frac{\pi}{2}\operatorname{sgn}a .
\]
取 \(R(z)=1/z\)、实轴上一阶极点 \(x=0\)，代入上式即得。

\subsection{计算步骤与易错点}

\begin{enumerate}
    \item 判断属于哪一型，检查条件（\(a>0\)？\(\deg\) 差够不够？有无实轴奇点？单位圆上是否解析？）；
    \item 选出围道与辅助函数（三角有理式取 \(z=e^{i\theta}\)；其余取上半平面半圆，被积函数乘 \(e^{iaz}\)）；
    \item 证明大（小）弧上积分极限为 \(0\) 或算出其贡献；
    \item 求奇点留数，代公式，最后按需取实部/虚部。
\end{enumerate}

\corepoint{四个高频错误：(1) 忘记 \textbf{\(a>0\)} 的条件，对 \(a<0\) 直接套公式会得到相反符号的答案（此时应先令 \(x\mapsto -x\) 或取共轭）；(2) 实轴上的极点仍按"上半平面极点"计入 \(2\pi i\) 项——极性不同，系数应是 \(\pi i\)；(3) \(\deg Q\ge\deg P+2\) 与 \(\deg Q\ge\deg P+1\) 两套条件记混，后者只在乘了 \(e^{iaz}\) 时成立；(4) 三角有理式误取单位圆外的极点。}

\section{辐角原理与儒歇定理}

本节是教材 \S5.4（标 \(*\) 的选学内容），给出"数零点个数"的工具，它本身也是留数定理的直接应用。

\subsection{对数留数与辐角原理}

\textbf{定理}（对数留数定理）：设 \(f(z)\) 在光滑闭曲线 \(C\) 上解析且不为零，在 \(C\) 内除有限个极点 \(b_j\)（阶数 \(\beta_j\)）外处处解析；\(a_k\) 是 \(f\) 在 \(C\) 内的零点（阶数 \(\alpha_k\)）；\(\varphi(z)\) 在 \(C\) 上及 \(C\) 内解析，则
\[
\frac{1}{2\pi i}\oint_C\varphi(z)\frac{f'(z)}{f(z)}\,dz
=\sum_{k=1}^{n}\alpha_k\varphi(a_k)-\sum_{j=1}^{m}\beta_j\varphi(b_j).
\]

\textbf{推论}（取 \(\varphi\equiv1\)）：
\[
\frac{1}{2\pi i}\oint_C\frac{f'(z)}{f(z)}\,dz=Z-P,
\]
其中 \(Z,P\) 分别是 \(f\) 在 \(C\) 内零点、极点的总个数（\(m\) 阶零点算 \(m\) 个）。左端称为 \(f(z)\) 关于 \(C\) 的\textbf{对数留数}。

把左端写成 \(\dfrac{1}{2\pi i}\oint_C d\ln f(z)\)：当 \(z\) 沿 \(C\) 正向绕行一周，\(w=f(z)\) 画出不过原点的闭曲线 \(\Gamma\)，\(\ln|w|\) 复原而无改变量，辐角的改变量是 \(2\pi\) 的整数倍，恰为 \(\Gamma\) 绕原点的圈数。于是

\textbf{辐角原理}：若 \(f(z)\) 在简单闭曲线 \(C\) 上及 \(C\) 内解析，且 \(C\) 上 \(f(z)\neq0\)，则 \(f\) 在 \(C\) 内零点的个数等于 \(z\) 沿 \(C\) 正向绕行一周时 \(f(z)\) 辐角的改变量除以 \(2\pi\)：
\[
Z=\frac{1}{2\pi}\Delta_C\arg f(z).
\]

\subsection{儒歇定理}

\textbf{儒歇（Rouch\'e）定理}：设 \(f(z)\) 与 \(g(z)\) 在简单闭曲线 \(C\) 上及 \(C\) 内解析，且在 \(C\) 上满足
\[
|f(z)-g(z)|<|f(z)|,
\]
则 \(f(z)\) 与 \(g(z)\) 在 \(C\) 内有相同个数的零点。

证明要点：条件即 \(C\) 上 \(\left|\dfrac{g}{f}-1\right|<1\)，故 \(F=g/f\) 把 \(C\) 映入不含原点的圆盘 \(|w-1|<1\)，\(\Gamma\) 绕原点圈数为 \(0\)，由辐角原理 \(g\) 与 \(f\) 的零点数相同。

\textbf{判定方程根数的步骤}：把方程整理成 \(P(z)=0\)；在给定闭曲线（多为 \(|z|=r\)）上取 \(P\) 中"占优"的部分作 \(f\)（最高次项、常数项或系数最大的项），其余作 \(g\)；在边界上验证 \(|P-f|<|f|\)；则 \(P\) 在该区域内的零点数与 \(f\) 相同，换半径即可分出圆环内的根数。

\corepoint{儒歇定理的条件必须 \emph{在闭曲线 \(C\) 上} 验证，而不是在区域内部；\(|f(z)-g(z)|<|f(z)|\) 与"\(|g|<|f|\) 且 \(P=f+g\)"是同一件事的两种写法。取 \(f\) 为最高次项时半径越大越容易占优，取 \(f\) 为常数项时只在小圆上占优——半径选得对不对，直接决定结论。}

\chapter{保形映射}

本章按教材“概念—分式线性映射—其性质—两个重要映射—初等函数构成的映射”五节展开。全章只有一个想法：解析函数在导数不为零处，局部相当于一次“旋转 + 伸缩”的相似变换，故能把区域“形状不变”地搬走；工程上（绕流、平面静电场、温度场）正是靠它把复杂区域化成上半平面或单位圆盘再求解。

\section{保形映射的概念}

\subsection{导数的辐角：旋转角}

设 \(w=f(z)\) 在 \(D\) 内解析，\(z_{0}\in D\)，\(f'(z_{0})\neq0\)（否则 \(\operatorname{Arg}f'(z_{0})\) 无意义）。取过 \(z_{0}\) 的光滑曲线 \(C:z=z(t)\)，\(z(t_{0})=z_{0}\)；割线的极限即切线，故 \(C\) 在 \(z_{0}\) 处的切线与正实轴夹角为 \(\operatorname{Arg}z'(t_{0})\)。映射把它映成过 \(w_{0}=f(z_{0})\) 的光滑曲线
\[
\Gamma:\ w=f[z(t)],\qquad w'(t_{0})=f'[z(t_{0})]\,z'(t_{0}),
\]
\(\Gamma\) 在 \(w_{0}\) 处切向的夹角为 \(\operatorname{Arg}w'(t_{0})\)。相减得

\keyformula{\operatorname{Arg}w'(t_{0})-\operatorname{Arg}z'(t_{0})=\operatorname{Arg}f'(z_{0})}

即像切向与原像切向之差恒等于 \(\operatorname{Arg}f'(z_{0})\)，与曲线形状无关。故称 \(\operatorname{Arg}f'(z_{0})\) 为 \(f(z)\) 在 \(z_{0}\) 处的\textbf{旋转角}。

再取第二条曲线 \(C_{2}\)，像为 \(\Gamma_{2}\)；两式相减消去 \(\operatorname{Arg}f'(z_{0})\)，得“\(C_{1},C_{2}\) 的夹角 \(=\) \(\Gamma_{1},\Gamma_{2}\) 的夹角”——夹角大小与旋转方向都不变，称映射在 \(z_{0}\) 处\textbf{保角}。

\textbf{定理}\quad 若 \(w=f(z)\) 在区域 \(D\) 内解析，则在 \(f'(z)\neq0\) 的点处映射 \(w=f(z)\) 是保角的。

\subsection{导数的模：伸缩率}

由 \(f'(z_{0})=\lim\limits_{z\to z_{0}}\dfrac{w-w_{0}}{z-z_{0}}\) 取模。设 \(\Delta s\) 为 \(C\) 上由 \(z_{0}\) 到 \(z\) 的弧长，\(\Delta s'\) 为像曲线 \(\Gamma\) 上的对应弧长，则

\keyformula{|f'(z_{0})|=\lim_{z\to z_{0}}\frac{|w-w_{0}|}{|z-z_{0}|}
=\lim_{z\to z_{0}}\frac{\Delta s'}{\Delta s}}

比值的极限与曲线形状无关，只与 \(z_{0}\) 有关，且 \(|\Delta w|\approx|f'(z_{0})|\,|\Delta z|\)：\(|f'(z_{0})|>1\) 伸长、\(<1\) 压缩、\(=1\) 近似保长。故 \(|f'(z_{0})|\) 称为\textbf{伸缩率}，此性质称\textbf{伸缩率不变性}。几何上小三角形映成近似相似的曲边三角形，小圆周映成近似圆周。

\subsection{保形映射的定义与判定}

\textbf{定义}\quad 凡在区域 \(D\) 内处处具有保角性和伸缩率不变性的\textbf{一一}映射，称为 \(D\) 内的\textbf{保形映射}。三条件缺一不可：前两条由“解析且 \(f'\neq0\)”保证，单叶性须单独验证——\(w=z^{2}\) 在 \(z\neq0\) 处保角，却在全平面不是一一的。

\textbf{定理}\quad 若 \(w=f(z)\) 在 \(z_{0}\) 处解析且 \(f'(z_{0})\neq0\)，则 \(w=f(z)\) 在 \(z_{0}\) 附近为保形映射。

\textbf{判定步骤}：（1）验证区域内解析；（2）解 \(f'(z)=0\)，把这些点剔除；（3）验证单叶性（能反解出唯一且解析的 \(z=f^{-1}(w)\)）；（4）三条齐备才是该区域上的保形映射。

\corepoint{最常错的是把“解析”当“保形”。\(w=z^{2}\) 在 \(z=0\) 处 \(f'(0)=0\)，该点不保角（顶角被放大 2 倍）。判“某点是否保角”只看 \(f'(z_{0})\neq0\)；判“是否保形映射”还要加单叶性。}

\section{分式线性映射}

\subsection{定义、逆映射与补充对应}

形如

\keyformula{w=\frac{az+b}{cz+d}\qquad(ad-bc\neq0)}

的映射称为\textbf{分式线性映射}（莫比乌斯变换）。\(ad-bc\neq0\) 必要：否则 \(\dfrac{dw}{dz}=\dfrac{ad-bc}{(cz+d)^{2}}\equiv0\)，\(w\) 恒为常数。逆映射 \(z=\dfrac{-dw+b}{cw-a}\) 仍是分式线性映射（行列式仍为 \(ad-bc\)）。

为得到扩充复平面上的一一对应，补两条：\(c=0\) 时 \(w=Az+B\ (A=a/d\neq0)\)，令 \(\infty\leftrightarrow\infty\)；\(c\neq0\) 时 \(z=-d/c\) 是一阶极点，令 \(z=-d/c\leftrightarrow w=\infty\)，又 \(\lim\limits_{z\to\infty}w=a/c\)，令 \(z=\infty\leftrightarrow w=a/c\)。

\subsection{分解为四种基本映射}

\(c\neq0\) 时作除法 \(w=\dfrac{a}{c}+\dfrac{bc-ad}{c^{2}}\cdot\dfrac{1}{z+d/c}\)，故可分解为四步：\(\zeta_{1}=z+\dfrac{d}{c}\)（平移）、\(\zeta_{2}=\dfrac{1}{\zeta_{1}}\)（反演）、\(\zeta_{3}=\dfrac{bc-ad}{c^{2}}\zeta_{2}\)（旋转+伸缩）、\(w=\zeta_{3}+\dfrac{a}{c}\)（平移）。

四种基本映射即：平移 \(w=z+\alpha\)、旋转 \(w=e^{i\theta}z\)、相似 \(w=rz\ (r>0)\)、反演 \(w=1/z\)；前三种导数为 \(1,e^{i\theta},r\)，处处保角。

反演最关键：先取关于单位圆周的对称点 \(\zeta=1/\bar z\)，再取关于实轴的对称点 \(w=\bar\zeta\)。若规定“伸向无穷远的两条曲线在 \(\infty\) 处的夹角等于它们经 \(w=1/z\) 变成的过原点像曲线在原点的夹角”，则 \(w=1/z\) 在扩充复平面处处保角。

\subsection{关于圆周的对称点}

设圆周 \(C:|z-z_{0}|=R\)，若 \(z_{1},z_{2}\) 满足

\keyformula{z_{2}-z_{0}=\frac{R^{2}}{\overline{z_{1}-z_{0}}}}

则称 \(z_{1},z_{2}\) 是\textbf{关于 \(C\) 的一对对称点}。记 \(z_{1}-z_{0}=r_{1}e^{i\theta}\) 可得 \(z_{2}-z_{0}=\dfrac{R^{2}}{r_{1}}e^{i\theta}\)，从而

\[
|z_{1}-z_{0}|\,|z_{2}-z_{0}|=R^{2}.
\]

即：在以 \(z_{0}\) 为起点的\textbf{同一射线}上，到圆心距离之积为 \(R^{2}\)，一内一外。约定\textbf{无穷远点关于圆周的对称点是圆心}；关于实轴的对称即共轭，关于 \(|z|=1\) 的对称点是 \(1/\bar z\)。

两个特征：（1）从圆外点 \(z_{2}\) 作切线（切点 \(T\)），由 \(T\) 作 \(z_{0}z_{2}\) 的垂线，垂足即 \(z_{1}\)；（2）\(z_{1},z_{2}\) 关于 \(C\) 对称的\textbf{充要条件}是过 \(z_{1},z_{2}\) 的任一圆周都与 \(C\) \textbf{正交}。

\subsection{三条基本性质}

\textbf{定理}（一一对应与保角性）\quad 分式线性映射在扩充复平面上是一一对应的保角映射。（由四步分解逐步论证，保角性对复合封闭。）

\textbf{定理}（保圆性）\quad 分式线性映射把圆周映射成圆周，这里把\textbf{直线视为半径无穷大的圆周}。把 \(w=1/z\) 代入圆的一般方程 \(A(x^{2}+y^{2})+Bx+Cy+D=0\)，所得仍为同型方程。由此得判据：给定圆周（直线）上\textbf{无}点映成 \(\infty\) 时像是\textbf{圆周}；有\textbf{一个}点映成 \(\infty\) 时像是\textbf{直线}。

\textbf{定理}（保对称性）\quad 若 \(z_{1},z_{2}\) 关于 \(C\) 对称，则其像 \(w_{1},w_{2}\) 关于像曲线 \(\Gamma\) 也对称。证明：过 \(w_{1},w_{2}\) 的任一圆周必是过 \(z_{1},z_{2}\) 的某圆经映射所得，它与 \(C\) 正交，由保角性也与 \(\Gamma\) 正交，用特征（2）即得。

\corepoint{保圆性说的是“圆周与直线”这一类对象封闭，不是“圆一定变圆”。像里出现直线，一定是原像上 \(z=-d/c\) 被映成了 \(\infty\)；判断时先找 \(cz+d=0\) 的点，看它是否落在所给圆周上。}

\section{唯一决定分式线性映射的条件}

\subsection{三对点定理与交比}

四个系数中用一个非零数同除分子分母可化为三个，即\textbf{只有三个独立常数}，故三个条件即可唯一确定。

\textbf{定理}（存在唯一性）\quad 对扩充 \(z\) 平面上任意三个相异点 \(z_{1},z_{2},z_{3}\) 与扩充 \(w\) 平面上任意三个相异点 \(w_{1},w_{2},w_{3}\)，存在唯一的分式线性映射把 \(z_{k}\) 依次映成 \(w_{k}\ (k=1,2,3)\)。

不必解方程组，直接用\textbf{交比}：

\keyformula{\frac{w-w_{1}}{w-w_{2}}:\frac{w_{3}-w_{1}}{w_{3}-w_{2}}
=\frac{z-z_{1}}{z-z_{2}}:\frac{z_{3}-z_{1}}{z_{3}-z_{2}}}

解出 \(w\) 即得。若某已知点为 \(\infty\)（如 \(w_{3}=\infty\)），只须把含 \(\infty\) 的项整体去掉：
\[
\frac{w-w_{1}}{w-w_{2}}=\frac{z-z_{1}}{z-z_{2}}:\frac{z_{3}-z_{1}}{z_{3}-z_{2}} .
\]

\textbf{定理}（圆到圆）\quad 任一圆都可用分式线性函数映成任一圆：在两圆上各取相异三点，由唯一性定理即得。

\subsection{像域位置的确定：取点法与绕向法}

保圆性只给“边界到边界”，内外需另判。只有两种可能：若 \(C\) 内两点的像一个在内一个在外，像弧与 \(\Gamma\) 的交点就同时是 \(C\) 上一点与线段上一点的像，与一一对应矛盾。

\begin{enumerate}
    \item \textbf{取点法}：在 \(C\) 内任取 \(z_{0}\) 算像 \(w_{0}\)，\(w_{0}\) 在 \(\Gamma\) 内则内部映内部，否则映外部。
    \item \textbf{绕向法}：在 \(C\) 上依序取 \(z_{1},z_{2},z_{3}\)，像为 \(w_{1},w_{2},w_{3}\)。绕向 \(z_{1}\to z_{2}\to z_{3}\) 与 \(w_{1}\to w_{2}\to w_{3}\) \textbf{相同}则内部映内部，相反则映外部。
\end{enumerate}

\(C\) 为圆周、\(\Gamma\) 为直线时，结论改为“\(C\) 内部映成 \(\Gamma\) 的某一侧半平面”，哪一侧仍由上述两法确定。

\subsection{圆弧所围区域的像}

两圆弧所围区域按“交点是否映成 \(\infty\)”分三种像域：无点映成 \(\infty\) 时仍是两圆弧所围区域；圆周上有一点映成 \(\infty\) 时变成一圆弧与一直线所围的区域；两圆\textbf{交点}之一映成 \(\infty\) 时变成角形域。第三条是解题主力——两圆同时成“直线”，区域立刻成角形域，再配幂函数即可张成半平面。

\corepoint{绕向法翻车多在“三点顺序”：顺序取反，绕向就反，内外结论整体颠倒。稳妥做法是用取点法独立验一次——取区域里最好算的点（圆心、原点、\(i\) 等），看像落在哪边。}

\section{两个重要的分式线性映射}

\subsection{将上半平面映成单位圆盘}

设所求为 \(w=\dfrac{az+b}{cz+d}\)，它把实轴映成 \(|w|=1\)。上半平面内取 \(z=\lambda\ (\operatorname{Im}\lambda>0)\) 映成圆心 \(0\)；\(\bar\lambda\) 是 \(\lambda\) 关于实轴的对称点，\(0\) 与 \(\infty\) 关于 \(|w|=1\) 对称，由保对称性 \(z=\bar\lambda\) 必映成 \(\infty\)，故 \(a\lambda+b=0,\ c\bar\lambda+d=0\)，得 \(w=\dfrac{a}{c}\cdot\dfrac{z-\lambda}{z-\bar\lambda}\)；再由 \(z=0\) 时 \(|w|=1\) 定出 \(\left|\dfrac{a}{c}\right|=1\)，即 \(\dfrac{a}{c}=e^{i\theta}\)。故

\keyformula{w=e^{i\theta}\frac{z-\lambda}{z-\bar\lambda}
\qquad(\operatorname{Im}\lambda>0,\ \theta\in\mathbb{R})}

它把上半平面内以 \(\lambda,\bar\lambda\) 为对称点的圆周族 \(\left|\dfrac{z-\lambda}{z-\bar\lambda}\right|=k\ (0<k<1)\) 映成同心圆周族 \(|w|=k\)。

\subsection{将单位圆盘映成单位圆盘}

设 \(z=\alpha\ (|\alpha|<1)\) 映成 \(w=0\)；\(\alpha\) 关于 \(|z|=1\) 的对称点 \(1/\bar\alpha\) 对应 \(w=\infty\)。由 \(a\alpha+b=0,\ c/\bar\alpha+d=0\) 得
\[
w=\frac{a}{c}\cdot\frac{z-\alpha}{z-1/\bar\alpha}
=k\,\frac{z-\alpha}{1-\bar\alpha z},\qquad k=-\frac{a\bar\alpha}{c}.
\]
\(|z|=1\) 时 \(\bar z=1/z\)，算得 \(\left|\dfrac{z-\alpha}{1-\bar\alpha z}\right|=1\)，故 \(|w|=1\) 迫使 \(|k|=1\)，即 \(k=e^{i\varphi}\)。故

\keyformula{w=e^{i\varphi}\frac{z-\alpha}{1-\bar\alpha z}
\qquad(|\alpha|<1,\ \varphi\in\mathbb{R})}

\subsection{使用步骤}

（1）\textbf{定型}：上半平面→圆盘用 \(\dfrac{z-\lambda}{z-\bar\lambda}\)，圆盘→圆盘用 \(\dfrac{z-\alpha}{1-\bar\alpha z}\)；两域同型时先用平移、伸缩化成标准形。（2）\textbf{定参数}：由“哪点映成圆心”唯一确定 \(\lambda\) 或 \(\alpha\)。（3）\textbf{定旋转因子}：剩下一个 \(e^{i\theta}\)（或 \(e^{i\varphi}\)），由第二个条件（某点的像或 \(\arg f'(z_{0})\)）解出。（4）\textbf{验算}：取边界三点看像是否在边界上，取内点看像是否在内部。

\corepoint{\(\theta\) 是整块的 \(e^{i\theta}\)，不是实数倍数。求 \(\arg f'(z_{0})\) 时先用 \(\theta\) 表示 \(f'(z_{0})\) 再让辐角相等；直接给 \(f'(z_{0})\) 赋值会漏掉模的条件（模已被“映成单位圆”固定为 1）。}

\needspace{9\baselineskip}
\section{几个初等函数所构成的映射}

\subsection{幂函数 \(w=z^{n}\)}

\(w=z^{n}\ (n=2,3,\dots)\) 处处可导，除原点外 \(nz^{n-1}\neq0\)，故除原点外处处保角。设 \(z=re^{i\theta}\)，则 \(w=r^{n}e^{in\theta}\)，即 \(\rho=r^{n},\ \varphi=n\theta\)：圆周 \(|z|=r\) 映成 \(|w|=r^{n}\)（单位圆映成单位圆），射线 \(\arg z=\theta_{0}\) 映成 \(\arg w=n\theta_{0}\)；以原点为顶点、顶角 \(\alpha\) 的角形域映成顶角 \(n\alpha\) 的角形域。

\keyformula{w=z^{n}:\quad 0<\arg z<\alpha\ \longmapsto\ 0<\arg w<n\alpha
\qquad\left(0<\alpha<\frac{2\pi}{n}\right)}

特别地 \(0<\arg z<\dfrac{2\pi}{n}\) 映成沿正实轴剪开的 \(w\) 平面 \(0<\arg w<2\pi\)：一边 \(\arg z=0\) 映成正实轴\textbf{上沿}，另一边映成\textbf{下沿}，两域一一对应。反函数 \(z=w^{1/n}\) 的单值解析分支把顶角 \(n\alpha\) 的角形域映成顶角 \(\alpha\) 的角形域。故放大或缩小角形域顶角，就用幂函数或根式。

\subsection{指数函数与对数函数}

\(w=e^{z}\) 以 \(2\pi i\) 为周期，只需在带形域 \(0<\operatorname{Im}z<2\pi\) 中研究。设 \(z=x+iy\)，\(w=e^{x}e^{iy}\)，则 \(|w|=e^{x},\ \arg w=y\)：竖直线 \(x=c\) 映成射线 \(\arg w=c\)；水平线段 \(y=c\) 映成圆周 \(|w|=e^{c}\)（去掉 \(w=0\)）；边界 \(y=0\) 与 \(y=2\pi\) 分别映成正实轴的上沿与下沿。故

\keyformula{w=e^{z}:\quad 0<\operatorname{Im}z<h\ (0<h\le 2\pi)
\ \longmapsto\ 0<\arg w<h}

因 \((e^{z})'=e^{z}\neq0\)，处处保角。反之 \(z=\ln w\) 的单值解析分支把角形域 \(0<\arg w<h\) 保形映成带形域 \(0<\operatorname{Im}z<h\)，分支的值完全由带形域确定——带形域与角形域互化，就靠这一对函数。

\subsection{儒可夫斯基函数（选学）}

\keyformula{w=\frac{1}{2}\left(z+\frac{a^{2}}{z}\right)\qquad(a>0)}

称为\textbf{儒可夫斯基（Joukowski）函数}。改写为 \(w-a=\dfrac{(z-a)^{2}}{2z}\)、\(w+a=\dfrac{(z+a)^{2}}{2z}\)，相除得
\[
\frac{w-a}{w+a}=\left(\frac{z-a}{z+a}\right)^{2},
\]
即由 \(w_{1}=\dfrac{z-a}{z+a}\)、\(w_{2}=w_{1}^{2}\)、\(\dfrac{w-a}{w+a}=w_{2}\) 三步复合：第一步把 \(\pm a\) 映成 \(0\) 与 \(\infty\)，把过 \(\pm a\) 的圆周映成过原点的直线（\(z\) 实数时 \(\dfrac{dw_{1}}{dz}>0\)，故外部映成含正实轴的半平面）；第二步平方映成沿射线剪开的平面；第三步映成沿连接 \(-a\) 与 \(a\) 的圆弧割开的平面。

故它把过 \(\pm a\) 的圆周 \(C\) 的外部一一保形映成“除去连接 \(a\) 与 \(-a\) 的圆弧”的平面。\(C\) 取 \(|z|=a\) 时圆弧退化为线段 \([-a,a]\)（\(\operatorname{Im}w=0\)），这正是机翼剖面（儒可夫斯基翼型）绕流的几何基础：割痕即翼型。

\subsection{构造保形映射的步骤}

把 \(D\) 保形映成标准域（上半平面或单位圆盘），按序试：（1）\textbf{看边界}是直线、圆周或兼有；（2）\textbf{去交点}：用分式线性映射把交点之一送到 \(\infty\)，圆变直线、区域变角形域；（3）\textbf{改顶角}：用 \(w=z^{n}\)（放大 \(n\) 倍）或 \(z=w^{1/n}\)（缩小为 \(1/n\)）把顶角调成 \(\pi\)；（4）\textbf{换带形}：带形域用 \(w=e^{z}\)、角形域用 \(z=\ln w\)；（5）\textbf{收尾}：用上节标准映射送到目标域并写出复合函数。遇带割痕区域，在（2）（3）之间插入儒可夫斯基函数或“平方—平移—开方”展平。

\begin{center}
\begin{tabular}{@{}lll@{}}
\toprule
目标 & 首选工具 & 依据 \\
\midrule
圆（直线）$\to$ 圆（直线） & 分式线性映射 & 保圆性 \\
两圆交点送到 \(\infty\) & 分式线性映射 & 保圆性 + 保角性 \\
角形域顶角放大／缩小 & \(w=z^{n}\)／\(z=w^{1/n}\) & 幂函数 \\
带形域与角形域互化 & \(w=e^{z}\)／\(z=\ln w\) & 指数、对数函数 \\
割痕、圆弧展平 & \(w=\frac{1}{2}(z+a^{2}/z)\) & 儒可夫斯基函数 \\
上半平面 $\leftrightarrow$ 圆盘 & \(e^{i\theta}(z-\lambda)/(z-\bar\lambda)\) & 保对称性 \\
圆盘 $\to$ 自身 & \(e^{i\varphi}(z-\alpha)/(1-\bar\alpha z)\) & 保对称性 \\
\bottomrule
\end{tabular}
\end{center}

\corepoint{四个高频错误：(1) 只验证“解析”就断言保形，漏掉 \(f'(z)\neq0\) 与一一性；(2) 用 \(w=z^{n}\) 忘记角形域要求 \(\alpha<\frac{2\pi}{n}\)，顶角超界后像域绕回、不再单叶；(3) 圆周被映成直线后仍按“圆内”处理内部，忘了直线只有“某一侧”，必须用取点法重新定位；(4) 儒可夫斯基函数的 \(\dfrac{w-a}{w+a}=\left(\dfrac{z-a}{z+a}\right)^{2}\) 漏掉平方，丢掉“双叶到单叶”这关键一步。}

\chapter{傅里叶变换}

积分变换的共同思想是：用一个积分核把"函数"换成"另一个函数"，让微分、卷积等难算的运算在变换域里变成乘法。傅里叶变换把时域信号换到频域，是频谱分析、信号处理与线性系统理论的数学基础。本章按讲义三节展开：可积函数的傅里叶变换、变换的性质、广义傅里叶变换，末尾补上常用变换对与应用方法。

\section{可积函数的傅里叶变换}

\subsection{从傅里叶级数到傅里叶积分}

设 \(f(t)\) 是实数轴上以 \(T\) 为周期的连续（或只有有限个第一类间断点）函数，\(\omega=\dfrac{2\pi}{T}\)。在区间 \(\left[-\dfrac{T}{2},\dfrac{T}{2}\right]\) 上平方可积的复值周期函数空间有正交基 \(\left\{e^{in\omega t}\right\}_{n\in\mathbb{Z}}\)，故
\[
f(t)\sim\sum_{n\in\mathbb{Z}}a_ne^{in\omega t},\qquad
a_n=\frac{1}{T}\int_{-T/2}^{T/2}f(t)e^{-in\omega t}dt .
\]
这是\textbf{离散}频谱：只有 \(\omega\) 的整数倍频率出现。

对一般的（非周期）绝对可积函数 \(f(t)\)，取截断函数 \(f_T(t)=f(t)\chi_{[-T/2,T/2]}(t)\)，把它看作周期 \(T\) 的函数作展开。令 \(T\to+\infty\)：
\begin{itemize}
    \item 相邻频率间隔 \(\omega=\dfrac{2\pi}{T}\to0\)，离散谱线变得无限密，成为\textbf{连续谱}；
    \item 系数 \(T a_n=\displaystyle\int_{-T/2}^{T/2}f(t)e^{-in\omega t}dt\) 趋于一个有界函数在 \(\omega_n=n\omega\) 处的值；
    \item 级数求和按"积分的黎曼和极限"化为积分。
\end{itemize}
于是得到一对积分公式——这就是傅里叶变换与傅里叶逆变换。

\subsection{傅里叶变换与逆变换的定义}

若广义积分 \(\displaystyle\int_{-\infty}^{+\infty}|f(t)|dt\) 收敛，则称实数轴上的复值函数 \(f(t)\) 是\textbf{绝对可积}的。

\begin{itemize}
    \item \textbf{傅里叶变换}：设 \(f(t)\) 在时域上绝对可积，定义
\end{itemize}

\keyformula{F(\omega)=\mathscr{F}[f(t)]=\int_{-\infty}^{+\infty}f(t)e^{-i\omega t}dt}

\begin{itemize}
    \item \textbf{傅里叶逆变换}：设 \(F(\omega)\) 在频域上绝对可积，定义
\end{itemize}

\keyformula{f(t)=\mathscr{F}^{-1}[F(\omega)]=\frac{1}{2\pi}\int_{-\infty}^{+\infty}F(\omega)e^{i\omega t}d\omega}

称 \(F(\omega)\) 为 \(f(t)\) 的\textbf{像函数}（频谱），\(f(t)\) 为 \(F(\omega)\) 的\textbf{原函数}。

\corepoint{傅里叶变换有三种常见约定：(1) 本章约定 \(e^{-i\omega t}\) 配 \(\dfrac{1}{2\pi}\)，系数全在逆变换上；(2) 对称约定 \(\dfrac{1}{\sqrt{2\pi}}\) 两边各一个；(3) 频率用 \(f=\omega/2\pi\) 时不用 \(2\pi\) 因子。三种约定的 \emph{正变换} 结果相差常数，混用会让公式差 \(2\pi\) 的幂。做题、查表必须锁定同一种约定。}

\subsection{傅里叶积分定理}

\textbf{定理}（傅里叶积分定理）：设 \(f(t)\) 是傅里叶可积函数，\(F(\omega)=\mathscr{F}[f(t)]\)，则
\[
\frac{1}{2\pi}\int_{-\infty}^{+\infty}F(\omega)e^{i\omega t}d\omega=
\begin{cases}
f(t), & t\ \text{是}\ f\ \text{的连续点},\\[4pt]
\dfrac{f(t+0)+f(t-0)}{2}, & t\ \text{是}\ f\ \text{的第一类间断点}.
\end{cases}
\]
这与傅里叶级数"收敛于左右极限平均值"的结论完全平行。间断点处逆变换给出的是左右极限的均值，因此逆变换的结果在间断点上与原来的 \(f(t)\) 可以有差别（不影响积分意义下的相等）。

\subsection{频谱}

\(F(\omega)\) 一般是复值函数，记
\[
F(\omega)=|F(\omega)|e^{i\varphi(\omega)} .
\]
\begin{itemize}
    \item \(|F(\omega)|\) 称为\textbf{振幅频谱}（幅频特性），\(\varphi(\omega)=\arg F(\omega)\) 称为\textbf{相位频谱}（相频特性）；
    \item 对非周期函数，\(F(\omega)\) 是频率的连续函数，称为\textbf{频谱密度}：它在频带 \([\omega,\omega+d\omega]\) 上的积分才表示该频段的"含量"，单独一点的 \(F(\omega)\) 没有振幅的量纲意义；
    \item 频谱密度一般是复值的，且 \(|F(-\omega)|=|F(\omega)|\)（对实函数 \(f\)）。
\end{itemize}

\textbf{两条定性结论}（由相似性质直接得出，后续常用）：
\begin{enumerate}
    \item \textbf{时宽与带宽成反比}：脉冲在时域越窄（\(a\) 越大），其频谱在主瓣之外衰减越慢、频带越宽；
    \item \textbf{矩形脉冲的频谱是 sinc 型}：\(\mathscr{F}\left[\chi_{[-a,a]}\right](\omega)=\dfrac{2\sin a\omega}{\omega}\)，在 \(\omega=\dfrac{k\pi}{a}\) 处为零，能量集中在主瓣 \(|\omega|<\dfrac{\pi}{a}\) 内。
\end{enumerate}

\subsection{单位阶跃函数}

\[
u(t)=
\begin{cases}
0, & t<0,\\
1, & t>0,
\end{cases}
\qquad
\mathscr{F}[u(t)]=\pi\delta(\omega)+\frac{1}{i\omega}.
\]
\(u(t)\) 本身不绝对可积，其变换是\textbf{广义}傅里叶变换（见 \S 本章第 3 节），这里先记结论。引入 \(u(t)\) 的意义在于：一切"从 \(t=0\) 开始"的信号都可写成 \(f(t)u(t)\)，这也正是下一章拉普拉斯变换的出发点。

\section{傅里叶变换的性质}

以下总设 \(F(\omega)=\mathscr{F}[f(t)]\)，\(F_1=\mathscr{F}[f_1]\)，\(F_2=\mathscr{F}[f_2]\)，且涉及的函数都绝对可积。

\subsection{线性性质}

对任意复数 \(k_1,k_2\)，
\[
\mathscr{F}[k_1f_1(t)+k_2f_2(t)]=k_1F_1(\omega)+k_2F_2(\omega),
\qquad
\mathscr{F}^{-1}[k_1F_1+k_2F_2]=k_1f_1(t)+k_2f_2(t).
\]
正、逆变换都是线性算子，这是"分解—变换—叠加"整套方法的前提。

\subsection{相似（伸缩）与平移性质}

设 \(a\neq0\)，\(b\in\mathbb{R}\)，则
\begin{equation}
\mathscr{F}[f(at+b)](\omega)=\frac{1}{|a|}e^{\,i\frac{b}{a}\omega}F\!\left(\frac{\omega}{a}\right).
\label{eq:similar}
\end{equation}

推导：令 \(s=at+b\)，\(t=\dfrac{s-b}{a}\)，\(dt=\dfrac{ds}{a}\)，代入定义整理即得；\(a<0\) 时积分限翻转，多出的负号与 \(\dfrac{1}{a}\) 合并成 \(\dfrac{1}{|a|}\)。

两个常用特例：
\begin{itemize}
    \item \textbf{时移}（\(a=1\)）：\(\mathscr{F}[f(t-t_0)]=e^{-i\omega t_0}F(\omega)\)，\(\mathscr{F}[f(t+t_0)]=e^{i\omega t_0}F(\omega)\)；
    \item \textbf{尺度}（\(b=0\)）：\(\mathscr{F}[f(at)]=\dfrac{1}{|a|}F\!\left(\dfrac{\omega}{a}\right)\)。
\end{itemize}

\textbf{频移性质}：
\[
\mathscr{F}\left[e^{i\omega_0t}f(t)\right]=F(\omega-\omega_0).
\]
由频移性质可得\textbf{调制性质}（工程上"搬移频谱"的依据）：
\[
\mathscr{F}[f(t)\cos\omega_0t]=\frac{1}{2}\big[F(\omega-\omega_0)+F(\omega+\omega_0)\big],
\]
\[
\mathscr{F}[f(t)\sin\omega_0t]=\frac{1}{2i}\big[F(\omega-\omega_0)-F(\omega+\omega_0)\big].
\]

\corepoint{时移与频移的符号方向最容易搞反：时域 \emph{右} 移 \(t_0\) 对应频域乘 \(e^{-i\omega t_0}\)（负号）；频域 \emph{右} 移 \(\omega_0\) 对应时域乘 \(e^{+i\omega_0t}\)（无负号）。另外尺度因子是 \(\dfrac{1}{|a|}\) 而不是 \(\dfrac{1}{a}\)，\(a<0\) 时尤为致命。}

\subsection{卷积与卷积定理}

\textbf{定义}：设 \(f_1,f_2\) 在实数轴上绝对可积，称
\[
(f_1*f_2)(t)=\int_{-\infty}^{+\infty}f_1(\tau)f_2(t-\tau)\,d\tau
\]
为 \(f_1\) 与 \(f_2\) 的\textbf{卷积}。

\textbf{封闭性}：两个绝对可积函数的卷积仍绝对可积。由积分估值不等式
\[
\int_{-\infty}^{+\infty}\left|\int_{-\infty}^{+\infty}f_1(\tau)f_2(t-\tau)d\tau\right|dt
\le\int|f_1(\tau)|\left(\int|f_2(t-\tau)|dt\right)d\tau
=\left(\int|f_1|\right)\left(\int|f_2|\right)<\infty .
\]

\textbf{运算律}：卷积满足交换律 \(f*g=g*f\)、结合律 \((f*g)*h=f*(g*h)\)、对加法的分配律 \(f*(g+h)=f*g+f*h\)。

\textbf{卷积定理}：
\[
\mathscr{F}[f_1*f_2](\omega)=F_1(\omega)F_2(\omega),
\]
\[
\mathscr{F}[f_1\cdot f_2](\omega)=\frac{1}{2\pi}\big(F_1*F_2\big)(\omega)
=\frac{1}{2\pi}\int_{-\infty}^{+\infty}F_1(\sigma)F_2(\omega-\sigma)\,d\sigma .
\]

\textbf{意义}：时域的卷积（本来要算一个含参积分）在频域变成\textbf{乘积}；反之，时域的乘积在频域变成卷积（带 \(\dfrac{1}{2\pi}\)）。线性时不变系统的输出就是"输入与冲激响应的卷积"，所以在频域只需相乘：\(Y(\omega)=H(\omega)F(\omega)\)。

\subsection{微分与积分性质}

\textbf{微分性质}：设 \(f(t)\) 是绝对可积的光滑函数，且对 \(k=0,1,\dots,n\) 有 \(\lim_{t\to\pm\infty}f^{(k)}(t)=0\)，则
\[
\mathscr{F}\left[f^{(n)}(t)\right](\omega)=(i\omega)^nF(\omega).
\]
一阶情形由分部积分 \(\displaystyle\int f'(t)e^{-i\omega t}dt=\big[f e^{-i\omega t}\big]_{-\infty}^{+\infty}+i\omega\int f e^{-i\omega t}dt\) 得到，边界项由绝对可积与"导数在无穷远趋于零"保证消失。

\textbf{频域微分}：\(\mathscr{F}\left[(-it)^nf(t)\right](\omega)=F^{(n)}(\omega)\)，即 \(\mathscr{F}[tf(t)]=iF'(\omega)\)。

\textbf{积分性质}：若 \(\displaystyle\int_{-\infty}^{+\infty}f(\tau)d\tau=0\)（等价于 \(F(0)=0\)），则
\[
\mathscr{F}\left[\int_{-\infty}^{t}f(\tau)\,d\tau\right]=\frac{F(\omega)}{i\omega}.
\]
不加 \(F(0)=0\) 的条件时，会出现与 \(u(t)\) 有关的 \(\pi F(0)\delta(\omega)\) 项。

\corepoint{微分性质把"求导"换成"乘 \(i\omega\)"，这是用傅里叶变换解微分方程的核心；但条件"\(f^{(k)}\) 在 \(\pm\infty\) 趋于零"不能省——对 \(u(t)\)、常数、\(\sin\omega_0t\) 这类函数，普通微分性质失效，必须走广义傅里叶变换。}

\subsection{帕塞瓦尔等式（能量积分）}

设 \(f_1,f_2\) 傅里叶可积，则
\[
\int_{-\infty}^{+\infty}\overline{f_1(t)}f_2(t)\,dt
=\frac{1}{2\pi}\int_{-\infty}^{+\infty}\overline{F_1(\omega)}F_2(\omega)\,d\omega .
\]
特别地（取 \(f_1=f_2=f\)）
\[
\int_{-\infty}^{+\infty}|f(t)|^2dt=\frac{1}{2\pi}\int_{-\infty}^{+\infty}|F(\omega)|^2d\omega .
\]

物理意义：左边是信号的总能量（时域），右边是 \(\dfrac{|F(\omega)|^2}{2\pi}\) 在频域上的积分，故称 \(|F(\omega)|^2\) 为\textbf{能量谱密度}。帕塞瓦尔等式即"能量守恒"：变换前后能量不变。它也是计算 \(\int_{-\infty}^{+\infty}\operatorname{sinc}^2\) 类实积分的常用手段——把实积分认成某个已知变换对的能量。

由帕塞瓦尔等式可算得（工程与光学中反复出现的两个积分）：
\[
\int_{-\infty}^{+\infty}\frac{\sin^2x}{x^2}dx=\pi,
\qquad
\int_{-\infty}^{+\infty}\frac{\sin^4x}{x^4}dx=\frac{2\pi}{3}.
\]

\subsection{性质一览}

\begin{center}
\begin{tabular}{@{}lll@{}}
\toprule
性质 & 时域 & 频域 \\
\midrule
线性 & \(k_1f_1+k_2f_2\) & \(k_1F_1+k_2F_2\) \\
相似 & \(f(at)\) & \(\dfrac{1}{|a|}F\!\left(\dfrac{\omega}{a}\right)\) \\
时移 & \(f(t\pm t_0)\) & \(e^{\pm i\omega t_0}F(\omega)\) \\
频移 & \(e^{i\omega_0t}f(t)\) & \(F(\omega-\omega_0)\) \\
调制 & \(f(t)\cos\omega_0t\) & \(\dfrac{1}{2}\left[F(\omega-\omega_0)+F(\omega+\omega_0)\right]\) \\
卷积 & \((f_1*f_2)(t)\) & \(F_1(\omega)F_2(\omega)\) \\
乘积 & \(f_1(t)f_2(t)\) & \(\dfrac{1}{2\pi}(F_1*F_2)(\omega)\) \\
微分 & \(f^{(n)}(t)\) & \((i\omega)^nF(\omega)\) \\
频域微分 & \((-it)^nf(t)\) & \(F^{(n)}(\omega)\) \\
积分 & \(\displaystyle\int_{-\infty}^{t}f(\tau)d\tau\) & \(\dfrac{F(\omega)}{i\omega}\)（需 \(F(0)=0\)） \\
能量 & \(\displaystyle\int|f(t)|^2dt\) & \(\dfrac{1}{2\pi}\displaystyle\int|F(\omega)|^2d\omega\) \\
\bottomrule
\end{tabular}
\end{center}

\section{广义傅里叶变换}

常数、正弦、余弦、阶跃等常见信号都不绝对可积，普通傅里叶变换对它们没有定义；能量集中在某一时刻的脉冲信号也无法用普通函数刻画。解决办法是把函数概念推广为\textbf{广义函数}。

\subsection{单位脉冲函数}

先从密度直观理解：若单位质量集中在原点，则线段 \((a,b]\) 上的质量为
\[
m\big((a,b]\big)=
\begin{cases}
1, & 0\in(a,b],\\
0, & 0\notin(a,b].
\end{cases}
\]
这种质量分布的"密度函数"记作 \(\delta(t)\)，形式上
\[
\delta(t)=
\begin{cases}
0, & t\neq0,\\
\infty, & t=0,
\end{cases}
\qquad
\int_{-\infty}^{+\infty}\delta(t)\,dt=1 .
\]

\textbf{筛选性质}：对任意连续函数 \(\varphi(t)\)，
\[
\int_{-\infty}^{+\infty}\delta(t)\varphi(t)\,dt=\varphi(0),
\qquad
\int_{-\infty}^{+\infty}\delta(t-t_0)\varphi(t)\,dt=\varphi(t_0).
\]
一般地，\(\displaystyle\int_{-\infty}^{+\infty}\delta^{(n)}(t)\varphi(t)dt=(-1)^n\varphi^{(n)}(0)\)。

\textbf{常用运算性质}：
\[
\delta(-t)=\delta(t),\qquad
\delta(at)=\frac{1}{|a|}\delta(t)\ (a\neq0),\qquad
t\,\delta(t)=0,
\]
\[
\delta(t-t_0)=\frac{1}{2\pi}\int_{-\infty}^{+\infty}e^{i\omega(t-t_0)}d\omega
\quad(\text{即}\ \mathscr{F}^{-1}[1]\ \text{的写法}).
\]

\textbf{广义函数观点}（讲义的形式化定义）：实数轴上的广义函数 \(T\) 是紧支集光滑函数空间 \(C_0^\infty(\mathbb{R})\) 到复数域的一个线性映射（泛函）。任一绝对可积函数 \(f(t)\) 都对应泛函
\[
\varphi\mapsto\int_{-\infty}^{+\infty}f(t)\varphi(t)\,dt,
\]
而 \(\delta\) 对应的泛函是 \(\varphi\mapsto\varphi(0)\)。广义函数的导数由分部积分定义：\(\langle T^{(n)},\varphi\rangle=(-1)^n\langle T,\varphi^{(n)}\rangle\)。在此意义下 \(\delta\) 可以合法地求导、作傅里叶变换。

\subsection{常用广义傅里叶变换}

\keyformula{\mathscr{F}[\delta(t)]=1,\qquad \mathscr{F}[1]=2\pi\delta(\omega)}

由 \(\mathscr{F}[1]=2\pi\delta(\omega)\) 与频移性质立刻得到
\[
\mathscr{F}\left[e^{i\omega_0t}\right]=2\pi\delta(\omega-\omega_0),
\]
从而
\[
\mathscr{F}[\cos\omega_0t]=\pi\big[\delta(\omega+\omega_0)+\delta(\omega-\omega_0)\big],
\]
\[
\mathscr{F}[\sin\omega_0t]=i\pi\big[\delta(\omega+\omega_0)-\delta(\omega-\omega_0)\big].
\]

\textbf{符号函数}：\(\operatorname{sgn}t=\dfrac{t}{|t|}\)（\(t\neq0\)，\(\operatorname{sgn}0=0\)），
\[
\mathscr{F}[\operatorname{sgn}t]=\frac{2}{i\omega}.
\]

\textbf{单位阶跃函数}：由 \(u(t)=\dfrac{1+\operatorname{sgn}t}{2}\)，
\[
\mathscr{F}[u(t)]=\pi\delta(\omega)+\frac{1}{i\omega}.
\]

\corepoint{广义变换的三个坑：(1) \(\mathscr{F}[1]=2\pi\delta(\omega)\) 中的 \(2\pi\) 与正变换是否带 \(\dfrac{1}{2\pi}\) 强相关，换约定必须重算；(2) \(\delta\) 不是普通函数，"\(\delta(0)=\infty\)"只是形式写法，真正定义靠筛选性质 \(\int\delta\varphi=\varphi(0)\)；(3) 求 \(\delta(at)\)、\(\delta'\) 之类时不能按普通函数复合求导，要用 \(\delta(at)=\dfrac{1}{|a|}\delta(t)\) 与 \(\int\delta'\varphi=-\varphi'(0)\)。}

\section{常用傅里叶变换对}

下表汇总课程与教材「附录 I 傅氏变换简表」中最常用的变换对（约定同本章：正变换核 \(e^{-i\omega t}\)，逆变换带 \(\dfrac{1}{2\pi}\)）。

\begin{center}
\begin{tabular}{@{}ll@{}}
\toprule
\(f(t)\) & \(F(\omega)=\mathscr{F}[f(t)]\) \\
\midrule
\(\delta(t)\) & \(1\) \\
\(1\) & \(2\pi\delta(\omega)\) \\
\(u(t)\) & \(\pi\delta(\omega)+\dfrac{1}{i\omega}\) \\
\(\operatorname{sgn}t\) & \(\dfrac{2}{i\omega}\) \\
\(e^{-\beta t}u(t)\ (\beta>0)\) & \(\dfrac{1}{\beta+i\omega}\) \\
\(e^{-\beta|t|}\ (\beta>0)\) & \(\dfrac{2\beta}{\beta^2+\omega^2}\) \\
\(\chi_{[-a,a]}(t)\ (\text{矩形脉冲})\) & \(\dfrac{2\sin a\omega}{\omega}\) \\
\(E\chi_{[-\tau/2,\tau/2]}(t)\ (\text{矩形单脉冲})\) & \(\dfrac{2E}{\omega}\sin\dfrac{\omega\tau}{2}\) \\
\(\cos\omega_0t\) & \(\pi\big[\delta(\omega+\omega_0)+\delta(\omega-\omega_0)\big]\) \\
\(\sin\omega_0t\) & \(i\pi\big[\delta(\omega+\omega_0)-\delta(\omega-\omega_0)\big]\) \\
\(e^{-\beta t^2}\ (\beta>0)\) & \(\sqrt{\dfrac{\pi}{\beta}}\,e^{-\frac{\omega^2}{4\beta}}\) \\
\(Ae^{-\alpha t^2}\ (\alpha>0)\ (\text{钟形脉冲})\) & \(A\sqrt{\dfrac{\pi}{\alpha}}\,e^{-\frac{\omega^2}{4\alpha}}\) \\
\(\dfrac{\sin\omega_0t}{\pi t}\ (\text{傅里叶核})\) & \(\begin{cases}1, & |\omega|\le\omega_0,\\ 0, & |\omega|>\omega_0\end{cases}\) \\
\bottomrule
\end{tabular}
\end{center}

\section{傅里叶变换的应用（方法与结论）}

\subsection{求解线性方程}

\textbf{思路}：把微分运算换成乘子。对含 \(t\) 的方程作傅里叶变换（或对空间变量 \(x\) 作变换），利用 \(\mathscr{F}[f^{(n)}]=(i\omega)^nF\)、\(\mathscr{F}[f*g]=F_1F_2\)，把微分方程（组）、甚至偏微分方程化为\textbf{代数方程}；解出像函数后再作逆变换。

\textbf{步骤}：
\begin{enumerate}
    \item 对哪个变量作变换，就要求该变量取遍 \((-\infty,+\infty)\)，并把初始/边界条件一起变换；
    \item 用微分性质、卷积定理把方程化为 \(F(\omega)\) 的代数方程；
    \item 解代数方程得 \(F(\omega)\)，化简（部分分式、配平方）；
    \item 作逆变换；若 \(F\) 是乘积形式，优先用卷积定理回到时域。
\end{enumerate}

傅里叶变换对\textbf{无界域}问题最自然（热传导、波动、静电场在半平面/全平面），对初值在 \(t=0\) 的"单边"问题则不如拉普拉斯变换方便。

\subsection{频谱分析与线性系统}

\begin{itemize}
    \item \textbf{非周期信号的频谱}：由 \(F(\omega)\) 的模与辐角给出幅频、相频特性；\(|F(\omega)|^2\) 是能量谱密度（帕塞瓦尔等式）。周期信号用傅里叶级数（离散谱），非周期信号用傅里叶变换（连续谱密度），两者由 \(T\to\infty\) 联系。
    \item \textbf{线性时不变系统}：设系统冲激响应为 \(h(t)\)，输入 \(f(t)\)，则输出 \(y(t)=h(t)*f(t)\)，在频域 \(Y(\omega)=H(\omega)F(\omega)\)。\(H(\omega)=\mathscr{F}[h(t)]\) 称为\textbf{频率响应}（系统函数），\(|H(\omega)|\) 的模决定各频率成分的放缩，辐角决定相移。
    \item \textbf{滤波}：理想低通滤波器的 \(H(\omega)\) 是矩形窗，其冲激响应是 sinc 型（傅里叶核对），说明"理想滤波"在时域必然拖尾——这正是时宽带宽积的体现。
\end{itemize}

\corepoint{应用层面的两个常见误区：(1) 把 \(F(\omega)\) 直接当"振幅"用——非周期信号的 \(F(\omega)\) 是频谱 \emph{密度}，只有 \(\int_{\omega_1}^{\omega_2}F d\omega\) 才有"含量"意义；(2) 求逆变换时忘记 \(\dfrac{1}{2\pi}\)，或把卷积定理的 \(\dfrac{1}{2\pi}\) 记到 \(\mathscr{F}[f_1*f_2]=F_1F_2\) 上（那个式子没有系数，系数只出现在乘积的变换式中）。}

\chapter{拉普拉斯变换}

常规傅里叶变换要求函数绝对可积，对常数、三角函数、指数增长函数都没有普通意义；而工程中往往只关心 \(t\ge0\) 的一段过程，\(t<0\) 的情形不必知道。把"截断到 \(t\ge0\)"与"再乘一个指数衰减因子"两件事合起来，就得到拉普拉斯（Laplace）变换，简称拉氏变换。它是解线性常微分方程（组）与线性电路暂态过程最顺手的工具。

\section{拉普拉斯变换的定义与存在定理}

\subsection{定义}

设 \(f(t)\) 是定义在 \(t\ge0\) 上（含原点某邻域）的函数，\(s=\beta+i\omega\) 为复参数，定义

\keyformula{F(s)=\mathscr{L}[f(t)]=\int_0^{+\infty}f(t)e^{-st}\,dt}

称 \(F(s)\) 为 \(f(t)\) 的\textbf{像函数}，\(f(t)\) 为 \(F(s)\) 的\textbf{原函数}，记 \(f(t)=\mathscr{L}^{-1}[F(s)]\)。

\textbf{与傅里叶变换的关系}：在收敛域内取 \(s=\beta+i\omega\)，则
\[
F(s)=\mathscr{L}[f(t)]=\mathscr{F}\left[f(t)u(t)e^{-\beta t}\right](\omega).
\]
即"拉氏变换 = 先截断、再乘衰减因子、再作傅里叶变换"。指数因子 \(e^{-\beta t}\) 的作用正是把增长的函数"压"成绝对可积，这就是拉氏变换能处理 \(t^n\)、\(e^{at}\) 的原因。

由定义直接可得几个最基本的像函数：
\[
\mathscr{L}[1]=\mathscr{L}[u(t)]=\frac{1}{s}\ (\operatorname{Re}s>0),\qquad
\mathscr{L}\left[e^{at}\right]=\frac{1}{s-a}\ (\operatorname{Re}s>\operatorname{Re}a),
\]
\[
\mathscr{L}[\cos at]=\frac{s}{s^2+a^2},\qquad
\mathscr{L}[\sin at]=\frac{a}{s^2+a^2}\qquad(\operatorname{Re}s>0).
\]

\subsection{存在定理}

\textbf{定理}：设 \(f(t)\) 在 \(t\ge0\) 上连续（或分段连续），若存在正实数 \(c\) 与 \(M\)，使
\[
|f(t)|\le Me^{ct},\qquad t\ge0,
\]
则 \(F(s)=\mathscr{L}[f(t)]\) 在半平面 \(\operatorname{Re}s>c\) 上收敛（且在该半平面内解析）。

证明：由积分估值不等式
\[
\int_0^{+\infty}\left|f(t)e^{-st}\right|dt
\le M\int_0^{+\infty}e^{-(\operatorname{Re}s-c)t}dt
=\frac{M}{\operatorname{Re}s-c}<\infty\quad(\operatorname{Re}s>c).
\]

满足上述条件的函数称为\textbf{指数增长级函数}。\(t^n\)、\(e^{at}\)、\(\sin at\)、\(\cos at\) 及其乘积都是指数增长级，故都可作拉氏变换；而 \(e^{t^2}\)、\(\dfrac{1}{t}\)（在 \(t=0\) 附近无界过快）不满足条件，不能直接作拉氏变换。

\subsection{收敛域}

\textbf{定理}：设 \(f(t)\) 连续，\(F(s)=\mathscr{L}[f(t)]\)，则
\begin{enumerate}[label=(\arabic*)]
    \item 若积分在 \(s_0\) 处收敛，则它在半平面 \(\operatorname{Re}s>\operatorname{Re}s_0\) 上处处收敛；
    \item 若积分在 \(s_0\) 处发散，则它在半平面 \(\operatorname{Re}s<\operatorname{Re}s_0\) 上处处发散。
\end{enumerate}

证明思路（1）：构造 \(\Phi(t)=\displaystyle\int_0^t f(\tau)e^{-s_0\tau}d\tau\)，由假设 \(\lim_{t\to+\infty}\Phi(t)=F(s_0)\)，故 \(\Phi\) 有界；再由分部积分把 \(e^{-(s-s_0)t}\) 提到积分外，即得 \(\operatorname{Re}s>\operatorname{Re}s_0\) 时收敛。第 (2) 条是 (1) 的逆否命题。

由此，拉氏变换的收敛域一定是形如 \(\operatorname{Re}s>\alpha\) 的\textbf{半平面}（或全平面、或空集），边界 \(\operatorname{Re}s=\alpha\) 上的行为需单独讨论。收敛域与像函数一起才唯一确定原函数（因为不同的 \(f\) 在 \(t<0\) 的取值不影响像函数）。

\corepoint{三个固定认识：(1) 拉氏变换是\textbf{单边}变换，积分从 \(0\) 起，\(t<0\) 的信息被丢掉，所以逆变换得到的原函数默认"\(t<0\) 时取 0"；(2) \(s\) 是复参数，写 \(\operatorname{Re}s>c\) 而不是 \(s>c\)；(3) 本节以后一律默认 \(s\) 取在收敛域内，否则积分不收敛、像函数无意义。}

\section{拉普拉斯变换的性质}

以下设 \(\mathscr{L}[f(t)]=F(s)\)，\(\mathscr{L}[g(t)]=G(s)\)，且 \(s\) 都在相应收敛域内。

\subsection{线性性质}

对任意复数 \(\alpha,\beta\)，
\[
\mathscr{L}[\alpha f(t)+\beta g(t)]=\alpha F(s)+\beta G(s),\qquad
\mathscr{L}^{-1}[\alpha F(s)+\beta G(s)]=\alpha f(t)+\beta g(t).
\]

\subsection{相似、延迟与位移}

\textbf{相似性质}（\(a>0\)）：
\[
\mathscr{L}[f(at)]=\frac{1}{a}F\!\left(\frac{s}{a}\right).
\]

\textbf{延迟性质}（\(\tau>0\)）：
\[
\mathscr{L}\left[f(t-\tau)u(t-\tau)\right]=e^{-s\tau}F(s).
\]

\textbf{位移性质}（\(a\in\mathbb{C}\)）：
\[
\mathscr{L}\left[e^{at}f(t)\right]=F(s-a).
\]

\corepoint{延迟性质里被延迟的部分必须写成 \(f(t-\tau)u(t-\tau)\)：若只写 \(f(t-\tau)\)，则 \(t<\tau\) 时它并不为零，与"函数在 \(t<0\) 处为 0"的单边约定冲突，结果会多出错误项。另外注意"延迟"改的是像函数乘 \(e^{-s\tau}\)，"位移"改的是像函数的自变量 \(s-a\)，两者不可混。}

\subsection{微分性质}

\textbf{原函数求导}：
\[
\mathscr{L}\left[f'(t)\right]=sF(s)-f(0),
\]
\[
\mathscr{L}\left[f^{(n)}(t)\right]=s^nF(s)-s^{n-1}f(0)-s^{n-2}f'(0)-\cdots-f^{(n-1)}(0).
\]

\keyformula{\mathscr{L}\left[f^{(n)}(t)\right]=s^nF(s)-\sum_{k=1}^{n}s^{n-k}f^{(k-1)}(0)}

要求 \(f^{(k)}(t)\ (k\le n)\) 都是指数增长级。这组公式是"用拉氏变换解微分方程"的支点：求导变成乘 \(s\)，初始条件自动进入方程。

\textbf{像函数求导}：
\[
F'(s)=-\mathscr{L}[tf(t)],\qquad
F^{(n)}(s)=(-1)^n\mathscr{L}[t^nf(t)].
\]

\subsection{积分性质}

\textbf{原函数积分}：
\[
\mathscr{L}\left[\int_0^tf(\tau)\,d\tau\right]=\frac{F(s)}{s},
\qquad
\mathscr{L}\left[\underbrace{\int_0^t\!dt\int_0^t\!dt\cdots\int_0^t}_{n\ \text{次}}f(t)dt\right]=\frac{F(s)}{s^n}.
\]

\textbf{像函数积分}：
\[
\int_s^{+\infty}F(\sigma)\,d\sigma=\mathscr{L}\left[\frac{f(t)}{t}\right],
\]
\[
\underbrace{\int_s^\infty\!ds\int_s^\infty\!ds\cdots\int_s^\infty}_{n\ \text{次}}F(s)\,ds
=\mathscr{L}\left[\frac{f(t)}{t^n}\right].
\]

\subsection{卷积与卷积定理}

\textbf{拉氏卷积}（与傅里叶卷积不同：积分上限是 \(t\)，且两个因子都只取 \(t\ge0\) 的部分）：
\[
(f_1*f_2)(t)=\int_0^tf_1(\tau)f_2(t-\tau)\,d\tau
=\int_0^tf_1(t-\tau)f_2(\tau)\,d\tau
\qquad(t\ge0).
\]

\textbf{卷积定理}：
\[
\mathscr{L}[f_1*f_2](s)=F_1(s)F_2(s),
\qquad
\mathscr{L}^{-1}[F_1(s)F_2(s)]=f_1(t)*f_2(t).
\]
即像函数的乘积对应原函数的卷积。这与傅里叶变换的卷积定理完全平行，但拉氏卷积不必担心绝对可积性，工程上更常用。

\subsection{初值定理与终值定理}

\textbf{初值定理}：若 \(\lim_{s\to\infty}sF(s)\) 存在有限，则
\[
f(0^+)=\lim_{s\to+\infty}sF(s).
\]

\textbf{终值定理}：若 \(sF(s)\) 的奇点都位于左半平面（\(\operatorname{Re}s<0\)），且 \(\lim_{s\to0}sF(s)\) 存在有限，则
\[
f(+\infty)=\lim_{s\to0}sF(s).
\]

两个定理的价值在于：不必求逆变换，就能由像函数直接读出原函数的初值与终值（稳定性判断常用）。

\subsection{周期函数的像函数}

设 \(f(t)\) 是以 \(T\) 为周期的函数（\(t\ge0\)），则
\[
\mathscr{L}[f(t)]=\frac{1}{1-e^{-sT}}\int_0^Tf(t)e^{-st}\,dt .
\]
若 \(f(t)\) 是\textbf{交错周期函数}（\(f(t+T)=-f(t)\)），则
\[
\mathscr{L}[f(t)]=\frac{1}{1+e^{-sT}}\int_0^Tf(t)e^{-st}\,dt .
\]

\subsection{性质一览}

\begin{center}
\begin{tabular}{@{}lll@{}}
\toprule
性质 & 时域 & 像函数 \\
\midrule
线性 & \(\alpha f+\beta g\) & \(\alpha F+\beta G\) \\
相似 & \(f(at)\ (a>0)\) & \(\dfrac{1}{a}F\!\left(\dfrac{s}{a}\right)\) \\
延迟 & \(f(t-\tau)u(t-\tau)\) & \(e^{-s\tau}F(s)\) \\
位移 & \(e^{at}f(t)\) & \(F(s-a)\) \\
微分 & \(f^{(n)}(t)\) & \(s^nF(s)-\sum\limits_{k=1}^{n}s^{n-k}f^{(k-1)}(0)\) \\
像函数微分 & \(t^nf(t)\) & \((-1)^nF^{(n)}(s)\) \\
积分 & \(\displaystyle\int_0^tf(\tau)d\tau\) & \(\dfrac{F(s)}{s}\) \\
像函数积分 & \(\dfrac{f(t)}{t}\) & \(\displaystyle\int_s^{+\infty}F(\sigma)d\sigma\) \\
卷积 & \((f_1*f_2)(t)\) & \(F_1(s)F_2(s)\) \\
初值 & \(f(0^+)\) & \(\lim\limits_{s\to\infty}sF(s)\) \\
终值 & \(f(+\infty)\) & \(\lim\limits_{s\to0}sF(s)\) \\
\bottomrule
\end{tabular}
\end{center}

\section{拉普拉斯逆变换}

\subsection{反演积分与留数法}

由复反演公式，\(F(s)\) 的原函数可写成沿竖直直线的复积分
\[
f(t)=\frac{1}{2\pi i}\int_{\beta-i\infty}^{\beta+i\infty}F(s)e^{st}\,ds
\qquad(t>0),
\]
其中 \(\beta\) 取在收敛域内（即 \(F(s)\) 的所有奇点都在直线左侧）。

\textbf{定理}（留数法）：设 \(F(s)\) 在复平面上除有限个孤立奇点 \(s_1,s_2,\dots,s_n\) 外解析，且
\[
\lim_{s\to\infty}F(s)=0,
\]
则
\begin{equation}
f(t)=\mathscr{L}^{-1}[F(s)]=\sum_{k=1}^{n}\operatorname{Res}\left[F(s)e^{st},s_k\right]
\qquad(t>0),
\label{eq:invres}
\end{equation}
而 \(t<0\) 时 \(f(t)=0\)。

\textbf{计算步骤}：
\begin{enumerate}
    \item 验证 \(\lim_{s\to\infty}F(s)=0\)（分母次数高于分子次数即可），并找出 \(F\) 的全部奇点；
    \item 对 \(F(s)e^{st}\) 求各奇点的留数（一阶极点用 \(\lim(s-s_k)F(s)e^{st}\)；\(m\) 阶极点用求导公式）；
    \item 相加即得 \(t>0\) 时的 \(f(t)\)。
\end{enumerate}

\subsection{部分分式法}

把 \(F(s)\) 分解成若干"查表即得"的简单分式之和，是最常用的手段。设 \(F(s)=\dfrac{P(s)}{Q(s)}\)。

\begin{enumerate}
    \item \textbf{先化真分式}：若 \(\deg P\ge\deg Q\)，先作多项式除法，把多项式部分（对应 \(\delta\) 及其导数）与真分式分开；
    \item \textbf{分解分母}：\(Q(s)\) 分解为一次因式与不可约二次因式之积；
    \item \textbf{单极点}：\(\dfrac{A_k}{s-s_k}\)，其中
    \[
    A_k=\lim_{s\to s_k}(s-s_k)F(s)=\frac{P(s_k)}{Q'(s_k)}\quad(\text{若}\ s_k\ \text{是一阶零点});
    \]
    \item \textbf{\(m\) 重极点}：写成 \(\displaystyle\sum_{j=1}^{m}\frac{B_j}{(s-s_0)^j}\)，系数用
    \[
    B_{m-j+1}=\frac{1}{(j-1)!}\lim_{s\to s_0}\frac{d^{\,j-1}}{ds^{\,j-1}}\left[(s-s_0)^mF(s)\right]
    \]
    或直接用待定系数法；
    \item \textbf{共轭复极点}：把 \(\dfrac{Bs+C}{(s+\alpha)^2+\beta^2}\) 配成
    \[
    B\cdot\frac{s+\alpha}{(s+\alpha)^2+\beta^2}+\frac{C-B\alpha}{\beta}\cdot\frac{\beta}{(s+\alpha)^2+\beta^2},
    \]
    逆变换分别是 \(Be^{-\alpha t}\cos\beta t\) 与 \(\dfrac{C-B\alpha}{\beta}e^{-\alpha t}\sin\beta t\)。
\end{enumerate}

\subsection{卷积法}

当 \(F(s)=F_1(s)F_2(s)\) 且两个因子各自的逆变换容易求出时，用卷积定理
\[
f(t)=\mathscr{L}^{-1}[F_1F_2]=f_1(t)*f_2(t)=\int_0^tf_1(\tau)f_2(t-\tau)\,d\tau
\]
求原函数。它把"分解有理分式"的代数工作换成了"算一个卷积积分"，对含 \(e^{-s\tau}\) 的延迟型像函数尤其方便。

\corepoint{逆变换的三个易错点：(1) 留数法求的是 \(F(s)e^{st}\) 的留数，漏掉 \(e^{st}\) 就退回成求 \(F\) 的留数，结果必错；(2) 留数法要求 \(\lim_{s\to\infty}F(s)=0\)，对 \(F(s)=1\)（对应 \(\delta(t)\)）或多项式像函数不适用，必须另作处理；(3) 部分分式法必须先化成真分式，假分式的多项式部分对应的是 \(\delta(t)\) 及其导数，不能当作普通分式查表。}

\section{常用拉普拉斯变换对}

下表按教材「附录 II 拉氏变换简表」整理，约定 \(t\ge0\)（\(t<0\) 时各原函数取 0），\(m>-1\)，\(\Gamma\) 为伽马函数。

\begin{center}
\begin{tabular}{@{}ll@{}}
\toprule
\(f(t)\) & \(F(s)=\mathscr{L}[f(t)]\) \\
\midrule
\(1\ (u(t))\) & \(\dfrac{1}{s}\) \\
\(t\) & \(\dfrac{1}{s^2}\) \\
\(t^n\ (n=1,2,\dots)\) & \(\dfrac{n!}{s^{n+1}}\) \\
\(t^m\ (m>-1)\) & \(\dfrac{\Gamma(m+1)}{s^{m+1}}\) \\
\(\sqrt{t}\) & \(\dfrac{\sqrt{\pi}}{2s^{3/2}}\) \\
\(\dfrac{1}{\sqrt{t}}\) & \(\sqrt{\dfrac{\pi}{s}}\) \\
\(e^{at}\) & \(\dfrac{1}{s-a}\) \\
\(t^me^{at}\) & \(\dfrac{\Gamma(m+1)}{(s-a)^{m+1}}\) \\
\(\sin at\) & \(\dfrac{a}{s^2+a^2}\) \\
\(\cos at\) & \(\dfrac{s}{s^2+a^2}\) \\
\(\sinh at\) & \(\dfrac{a}{s^2-a^2}\) \\
\(\cosh at\) & \(\dfrac{s}{s^2-a^2}\) \\
\(e^{at}\sin bt\) & \(\dfrac{b}{(s-a)^2+b^2}\) \\
\(e^{at}\cos bt\) & \(\dfrac{s-a}{(s-a)^2+b^2}\) \\
\(t\sin at\) & \(\dfrac{2as}{(s^2+a^2)^2}\) \\
\(t\cos at\) & \(\dfrac{s^2-a^2}{(s^2+a^2)^2}\) \\
\(\dfrac{1}{a^2}\left(1-\cos at\right)\) & \(\dfrac{1}{s(s^2+a^2)}\) \\
\(\dfrac{1}{a^3}\left(at-\sin at\right)\) & \(\dfrac{1}{s^2(s^2+a^2)}\) \\
\(\dfrac{1}{a}\sin at-\dfrac{1}{b}\sin bt\) & \(\dfrac{b^2-a^2}{(s^2+a^2)(s^2+b^2)}\) \\
\(\cos at-\cos bt\) & \(\dfrac{(b^2-a^2)s}{(s^2+a^2)(s^2+b^2)}\) \\
\(\delta(t)\) & \(1\) \\
\(\delta^{(n)}(t)\) & \(s^n\) \\
\(f(t-\tau)u(t-\tau)\) & \(e^{-s\tau}F(s)\) \\
\(f(t)\ (T\ \text{为周期})\) & \(\dfrac{1}{1-e^{-sT}}\displaystyle\int_0^Tf(t)e^{-st}dt\) \\
\bottomrule
\end{tabular}
\end{center}

\section{拉普拉斯变换在解方程与电路中的应用}

\subsection{解线性常微分方程（组）}

利用微分性质，拉氏变换把"带初始条件的线性常微分方程"化成"像函数的代数方程"，解出像函数再逆变换即可。整个流程不需要先求通解、再代初始条件。

\textbf{步骤}：
\begin{enumerate}
    \item 设 \(\mathscr{L}[y(t)]=Y(s)\)，对方程两边作拉氏变换；
    \item 用微分性质把 \(y^{(k)}\) 换成 \(s^kY(s)-s^{k-1}y(0)-\cdots-y^{(k-1)}(0)\)，初始条件在这一步自动代入；
    \item 解出 \(Y(s)\)（有理分式）；
    \item 部分分式分解（或留数法）作逆变换，得 \(y(t)\)；
    \item 回代检验初值与终值（用初值、终值定理快速自查）。
\end{enumerate}

对\textbf{方程组}：设 \(\mathscr{L}[x]=X(s)\)、\(\mathscr{L}[y]=Y(s)\)，同样把每个方程变成 \(X,Y\) 的线性代数方程，解线性方程组后再逐个逆变换。这一步在电路（网孔法、节点法）里就是"把微分方程组变成复数代数方程组"。

\subsection{电路的 \(s\) 域模型}

对电路元件的时域关系作拉氏变换，得到\textbf{运算电路}（\(s\) 域模型）：

\begin{center}
\begin{tabular}{@{}lll@{}}
\toprule
元件 & 时域关系 & \(s\) 域关系 \\
\midrule
电阻 \(R\) & \(u=Ri\) & \(U(s)=RI(s)\) \\
电感 \(L\) & \(u=L\dfrac{di}{dt}\) & \(U(s)=sLI(s)-Li(0)\) \\
电容 \(C\) & \(i=C\dfrac{du}{dt}\) & \(U(s)=\dfrac{1}{sC}I(s)+\dfrac{u(0)}{s}\) \\
\bottomrule
\end{tabular}
\end{center}

即：电感等效为阻抗 \(sL\) 串联一个电压源 \(-Li(0)\)（或并联电流源 \(\dfrac{i(0)}{s}\)）；电容等效为阻抗 \(\dfrac{1}{sC}\) 串联电压源 \(\dfrac{u(0)}{s}\)（或并联电流源 \(Cu(0)\)）。初始条件由此变成"外加电源"，之后完全按直流电路的方法（串并联、节点法、网孔法）列代数方程。

\textbf{传递函数与响应}：
\begin{itemize}
    \item 令所有初始值为零，定义\textbf{传递函数}（网络函数）
    \[
    H(s)=\frac{Y(s)}{X(s)},
    \]
    它只由电路结构与参数决定，与输入无关；
    \item \textbf{冲激响应} \(h(t)=\mathscr{L}^{-1}[H(s)]\)，则任意输入下的零状态响应为 \(y(t)=h(t)*x(t)\)，即 \(Y(s)=H(s)X(s)\)；
    \item 全响应可分解为\textbf{零输入响应}（初始条件引起，令 \(X=0\)）与\textbf{零状态响应}（输入引起，令初始值为 0）之和；
    \item \(H(s)\) 的极点位置决定暂态衰减快慢与稳定性：极点全在左半平面（\(\operatorname{Re}s<0\)）时系统稳定，这与终值定理的条件一致。
\end{itemize}

\corepoint{解方程与电路应用中的高频错误：(1) 作变换时漏掉初始条件项（\(-f(0)\)、\(-s f(0)-f'(0)\) 等），这是最常见的失分点；(2) 逆变换前忘记把 \(Y(s)\) 化成真分式；(3) 混淆"零输入"与"零状态"——前者只由初始储能决定，后者只由外加激励决定，叠加时初始条件不能重复计入；(4) 用终值定理判稳态值时，忘记先确认 \(sF(s)\) 的极点全在左半平面（含虚轴极点时会得出错误结论）。}


\end{document}